% !TeX program = lualatex
\documentclass[11pt,a4paper]{ltjsarticle}
\usepackage[margin=25mm,headheight=17pt]{geometry}
\usepackage[haranoaji]{luatexja-preset}
\usepackage{amsmath,amssymb,amsthm,mathtools,mathrsfs,booktabs,array,enumitem,longtable}
\usepackage[unicode,hidelinks]{hyperref}
\usepackage{fancyhdr}
\pagestyle{fancy}\fancyhf{}\fancyfoot[C]{\thepage}
\fancyhead[L]{\small\nouppercase{\leftmark}}
\renewcommand{\sectionmark}[1]{\markboth{#1}{}}
\setlength{\emergencystretch}{3em}\allowdisplaybreaks[2]
\setlength{\parskip}{.15em}
\newtheoremstyle{jp}{.6em}{.6em}{\normalfont}{}{\bfseries}{.}{.5em}{}
\theoremstyle{jp}
\newtheorem{theorem}{定理}[section]
\newtheorem{proposition}[theorem]{命題}
\newtheorem{lemma}[theorem]{補題}
\newtheorem{corollary}[theorem]{系}
\newtheorem{definition}[theorem]{定義}
\newtheorem{remark}[theorem]{注意}
\newtheorem{example}[theorem]{例}
\newtheorem{question}[theorem]{問題}
\renewcommand{\proofname}{証明}
\newcommand{\kk}{\mathbb C}\newcommand{\PP}{\mathbb P}
\newcommand{\OO}{\mathcal O}\newcommand{\A}{\mathcal A}
\newcommand{\B}{B_{\mathrm{VK}}}\newcommand{\Blo}{B_{p,q}}
\newcommand{\HH}{\operatorname{HH}}\newcommand{\Hom}{\operatorname{Hom}}
\newcommand{\End}{\operatorname{End}}\newcommand{\Sym}{\operatorname{Sym}}
\newcommand{\Perf}{\operatorname{Perf}}\newcommand{\Der}{\operatorname{Der}}
\newcommand{\Inn}{\operatorname{Inn}}\newcommand{\gr}{\operatorname{gr}}
\newcommand{\rk}{\operatorname{rank}}\newcommand{\Spec}{\operatorname{Spec}}
\newcommand{\id}{\operatorname{id}}\newcommand{\Ext}{\operatorname{Ext}}
\newcommand{\Tor}{\operatorname{Tor}}\newcommand{\Ann}{\operatorname{Ann}}
\newcommand{\coker}{\operatorname{coker}}\newcommand{\Supp}{\operatorname{Supp}}
\newcommand{\pd}{\operatorname{pd}}\newcommand{\depth}{\operatorname{depth}}
\newcommand{\Frac}{\operatorname{Frac}}\newcommand{\Aut}{\operatorname{Aut}}
\newcommand{\Out}{\operatorname{Out}}\newcommand{\GL}{\operatorname{GL}}
\newcommand{\PGL}{\operatorname{PGL}}\newcommand{\Hilb}{\operatorname{Hilb}}
\newcommand{\dashsim}{\mathrel{\overset{\mathrm{bir}}{\sim}}}
\setcounter{tocdepth}{1}

\hypersetup{pdftitle={重み付き量子射影圏の変形環と対称性}}
\title{重み付き量子射影圏の変形環と対称性\\
\large 全重みの内在的な純粋部分、半直交補完と障害の比較}
\author{}\date{2026年10月30日\quad 日本語公表用原稿X・v40}
\begin{document}\maketitle
\begin{abstract}
重み $(1,1,1,m)$ の量子射影圏の全変形環、関数斜体、および半直交成分の対称性を調べる。
非線形CGG共鳴族の固定非零点のhullは、正共鳴で
$\mathbb C[[\gamma]]/(\gamma^{j+1})$、偶数端点で滑らかな $m+2$ 次元環となる。
端点の全hullは通常Calabi--YauのAS正則Ore族で実現され、その関数斜体は
Weyl型、非冪根量子トーラス型、有限位数型に完全分類できる。
PIファイバーはZariski稠密な非構成可能集合を成し、非PI端点へ非有界なPI次数で集積する。
別の引戻し型Ore族では
$\mathbb C[[t_1,t_2,a,b]]/(a^{m+2},ab,b^2)$ という非完全交叉hullを得る。
CGG族内の射影圏と導来圏の同値は $(m,j)$ の一致で特徴付けられる。
重み三の二つの量子VK成分は同じ三十五次元の滑らかなhullを持つが、
$HH^1$ のGerstenhaber作用で導来非同値と判定できる。
外的導分を保持する部分は一方で滑らかな四次元空間となり、
他方では正確に $\mathbb C[[u_1,\ldots,u_5,\gamma]]/(\gamma^2)$ となる。
全ての変形は $w$ フィルトレーションを保つ形に直せ、
純粋三変数部分の三方向と原始射の接合の三十二方向に全形式的に還元される。
純粋部分の対称性保持曲線は滑らかであるが、接合後にはその原点の平方零近傍だけが残る。
二重構造全体を実際の $\mathbf G_a$ 作用を持つ全次数の族で実現する。
後者の標準六項補完接線は内在的な二次対称性障害を持ち、
その初速度の任意の弧に沿う相対 $HH^2$ の捩れは正確に
$\mathbb C[[t]]/(t^2)$ である。
さらに全ての共鳴 $m=2j+r\ge3$ で、純粋三変数部分は任意のArtin基底上で内在的に持ち上がる。
純粋部分とその同定を固定した相対変形は、次元 $(m^2+9m+28)/2$ の滑らかな形式空間となる。
純粋部分の外的導分が全ての共鳴へ持ち上がるため、通常hullも
$R_{N_m}[[u_1,\ldots,u_{(m^2+9m+28)/2}]]$ と分解する。
従って同じ重みの全共鳴は、偶数端点を含めて同型な通常変形環を持つ。
代数と導分を同時に記録する障害複体から、外側の結合性障害と
内側の対称性障害を結ぶ自然な推移写像を構成し、標準補完の障害線上で非零と示す。
\end{abstract}
\tableofcontents\clearpage
\section{序論}

同じ関数斜体を持つ非可換射影空間を、射影導来圏の段階でどこまで区別できるか。
さらに、その半直交成分だけを取り出した後にも、量子化の違いは残るか。
本稿は重み $(1,1,1,m)$ の明示的な非線形CGG量子化族について、
固定非零点の変形環、Serre関手、およびGerstenhaber作用からこの問題を調べる。
四項成分では、代数自身の滑らかな変形基底の中に、
対称性を保持する条件によって非被約な二重構造が現れることを示す。
さらに引戻し型のOre量子化を加え、非完全交叉の変形環と、全変形の明示的なAS正則実現を得る。

対象は整数 $j,r\ge0$、$m=2j+r\ge3$ と
$\alpha\in\mathbb C$、$s\in\mathbb C^*$ による
$\mathcal S_{j,r}(\alpha,s)$ である。積は次節で完全に定義する。
各代数はAS正則四次元であり、標準 $(m+3)$ 項のfull strong exceptional collectionを持つ。
全ての族の次数零関数斜体は、$\mathbb C(C)$ 上の第一Weyl斜体に同型となる。
その一方、射影圏のアーベル圏同値と有界導来圏同値は共に
\[
 \boxed{(m,j)=(m',j')}
\]
によって特徴付けられる。固定した $m$ に対し、相異なる圏は
$\lfloor m/2\rfloor+1$ 個である。分類対象は、ここで明示したCGG族である。

\subsection{全共鳴の変形環と端点の相違}

標準全列の有限自己準同型代数を $A_{j,r}$ とすると、固定非零量子化のminiversal環は
\[
 R_{A_{j,r}}\simeq
 \begin{cases}
 \mathbb C[[\gamma]]/(\gamma^{j+1})&r>0,\\
 \mathbb C[[u_1,\ldots,u_{m+2}]]&r=0
 \end{cases}
\]
となる。第一の場合は任意に高い次数で初めて現れる障害を与え、
$j\ge1$ では一次変形があるにもかかわらず、全ての一パラメータ形式変形が自明である。
第二の場合は $m=2j$ の偶数重みの端点であり、$m+2$ 次元の滑らかな変形基底が現れる。
端点の証明では、混合Poisson写像の核をHamiltonianベクトル場で尽くし、
調和多項式の分解によって余核がただ一方向であることを示す。
その方向を非零混合点の横断パラメータが消すため、全局所環が正則になる。

Poissonスキームから固定量子代数への移送には、可換基点の形式同型を
$h$ 進に小さい量子化弧へ適用し、$\mathbb C((h))$ 上で再中心化する方法を用いる。
全局所環を使うため、接空間の一致だけに依存しない。
得られるのは有限Beilinson代数、従って射影導来圏の変形環である。
無限次元の次数付き環の全ての非次数付き変形を分類しているわけではない。


\subsection{全変形の座標環実現と非完全交叉hull}

偶数端点の $m+2$ 次元の形式変形は、三変数AS正則代数 $B_{\gamma,s}$ と
混合項 $H$ から作る $B_{\gamma,s}[w;\sigma,x^{-1}[\,-,H]]$ によって全て実現できる。
正規元 $x$ の正規化自己同型が $B$ のNakayama自己同型と一致するため、
中心元 $J=xw+H$ の固定性から拡大のNakayama自己同型が恒等写像となる。

これは各複素パラメータでAS正則となる有限パラメータ族である。
その次数零関数斜体を全パラメータで分類する。
混合項 $H$ の $m+1$ 方向は関数斜体から消えるが、非冪根の場合には残る $\gamma$ を完全に区別できる。
冪根の場合には中心変数を混ぜる整数変換により、同じ位数のパラメータが同じ斜体を与える。
この族のPI集合は構成可能でなく、任意に高いPI次数のファイバーが非PI端点へ集積する。

一方、共点三直線に基づく微分Ore量子化では、固定非零点のhullが
\[
 \mathbb C[[t_1,t_2,a,b]]/(a^{m+2},ab,b^2)
\]
となる。二つの滑らかな方向はOre係数 $A,B$ の変化で実現される。
横断Artin環の長さは $m+3$、socle次元は二であり、全hullはCohen--Macaulay型二の非完全交叉環である。
従って本稿は、一変数の高次障害、滑らかな全実現、および二変数の障害の相互作用を、
実際の非可換重み付き射影圏で比較できる。

\subsection{Serre関手と内在的な反標準環}

通常のCalabi--Yau代表に対して、反標準Hochschild--Serre環の全体を
$Z(\mathcal S)^{(m+3)}$ と同定する。
$r\ge2,r=1,r=0$ に応じて、この環は定数環、$D_{m+4}$ 型、$A_{m+2}$ 型となる。
正次数部分がある場合のstacky Projは、それぞれ
$(2,2,m+2)$ 型のorbifold曲線と $\mathbb P_{\mathrm{stack}}(2,m+1)$ である。
中心から自然変換を作るだけでなく、正twist双加群の全ての準同型が中心から来ることを証明する。
この全体性によって、中心の表示を射影圏に内在する不変量として用いられる。

中心の部分解消、central reduction、primitive idealについては、
独立に読める公表稿XI\cite{VKPrimitive36}にまとめる。
本稿では導来分類に必要な中心と反標準環の計算を扱う。

\subsection{同じ滑らかな変形基底を持つ二つの量子VK圏}

重み三では $(j,r)=(0,3),(1,1)$ の二つの非零族がある。
標準四項 $(0,1,3,4)$ を切り出した五十六次元代数を $B_0,B_1$ と書く。
両者は
\[
 \dim HH^\bullet(B_i)=(1,1,35,39),\qquad
 R_{B_i}\simeq\mathbb C[[z_1,\ldots,z_{35}]]
\]
を満たす。従って、Hochschild群の次元も、miniversal基底の抽象的な形式環も一致する。
それでも、二つの導来圏は同値ではない。

区別に使うのは、一次元空間 $HH^1(B_i)$ の非零元 $\xi$ による
\[
 [\xi,-]:HH^2(B_i)\longrightarrow HH^2(B_i)
\]
の冪零Jordan型である。二つの場合のブロックはそれぞれ
\[
 (13,10,7,5),\qquad (11,9,6,5,3,1)
\]
であり、第一冪の階数だけでも $31$ と $29$ に分かれる。
生成元を非零スカラー倍してもこのJordan型は変わらず、
Gerstenhaber括弧の導来不変性から、これは標準例外列の選択によらない区別となる。
可換基点も含めると、重み三の本CGG族の四項成分は正確に三つの導来同値類に分かれる。

標準六項補完の変形についても、二つの四項成分の滑らかな基底の中で、
補完像はそれぞれ被約点と長さ二の非被約点となる。
この補完像の相違と、Gerstenhaber作用による内在的な非同値証明は別の主張である。
任意の別の例外列による補完可能性まで分類したものではない。


さらに、相対Hochschild複体を縮約すると、対称性の消失を一列の微分で記述できる。
一般変形のHochschild次元は $(1,0,34,39)$ となる。
外的導分を保持する部分の接空間は四次元と六次元であり、
二次近傍はそれぞれ無関係な四変数と、六変数に一つの平方関係を課したものとなる。
代数自身の変形は無障害でも、基点の外的導分を保持する条件には二次障害が現れる。
原始射に支持を持つ積補正を用いると、二次障害の零方向は全次数で同時に積分できる。
さらに平方零方向を含む全次数の族を構成することで、対称性保持部分の全環は
\[
 R_{Z_{B_0}}\simeq\mathbb C[[u_1,\ldots,u_4]],\qquad
 R_{Z_{B_1}}\simeq
       \mathbb C[[u_1,\ldots,u_5,\gamma]]/(\gamma^2)
\]
と完全に決まる。
第二の環は滑らかな五次元部分の二重構造であり、
基底全体で二十四乗が零となる導分、従って $\mathbf G_a$ 作用を持つ。
その構成では、有理係数の有限行列から得た解について、
全ての結合則と導分方程式を多項式恒等式として直接確認する。
有限jetから全完備局所環を推測しているのではない。

この二重構造には、純粋な三変数部分からの説明もある。
全ての変形と基底変換を $w$ フィルトレーションを保つ形へ還元し、
純粋部分への制限の相対変形空間が三十二次元で滑らかであることを示す。
純粋部分自体は、固定した四対象の合成の像から復元できる。
その変形基底は滑らかな三次元、対称性保持部分は滑らかな曲線
$\operatorname{Spf}\mathbb C[[\gamma]]$ である。
原始射を接合すると、この曲線は $B_0$ では原点に、
$B_1$ では原点の平方零近傍に制限される。
従って上の第二の同型は、純粋曲線の環 $\mathbb C[[\gamma]]$ 上で取ることができる。

通常の三十五変数hullの中でも形式座標を選び、対称性保持イデアルを
\[
 I_{B_1}=(z_1,\ldots,z_{29},\gamma^2)
\]
と書ける。標準六項への補完像は、この二重構造の横断長さ二の切片となる。
その接線を初速度とするどの積分弧でも外的導分は二次で失われ、
相対Hochschild群には長さ二の捩れが現れる。
さらに一般の弧の捩れの長さも、$z_i$ の付値と $\gamma$ の付値の二倍の最小値で決まる。

\subsection{全重みの通常hullの分解と障害の自然な比較}

四項成分の純粋三変数部分が残る理由は、重み三に固有ではない。
全ての $m=2j+r\ge3$、$j,r\ge0$ で、任意のArtin変形における純粋部分は
合成の像から一意かつ自然に復元できる。
純粋部分の変形とその同定を固定した相対変形は、
\[
 k_m=\frac{m^2+9m+28}{2}
\]
次元の滑らかな形式空間となり、相対自己同型を持たない。
証明は任意の間隔に有効な反対称化と、CGG積の有限フィルトレーションによる。
さらに、三本の隣接射を反転する普遍局所化と、半古典Poisson構造の安定化群から
$HH^1(N_m)=\mathbb C[\log(1+N)]$ を全ての重みで証明する。
この生成導分は各四項成分へ持ち上がるため、通常miniversal環にも
\[
 R_{B_{j,r}}\simeq R_{N_m}[[u_1,\ldots,u_{k_m}]]
\]
という分解が成立する。特に、同じ重みの全共鳴は偶数端点を含めて同型な通常hullを持つ。
この同型は導分の作用や対称性保持部分を同一視しない。
一般の $m$ に対する純粋部分の環 $R_{N_m}$ 自体は、まだ決定していない。

また、六項側の結合性障害と四項側の対称性障害の対応を、cochainの段階で説明する。
積 $\mu$ と導分 $\delta$ の同時変形を記録する複体を
\[
 \mathcal P^n=C^n\oplus C^{n-1},\qquad
 \partial(f,g)=(df,[\delta,f]-dg)
\]
と置く。これはAssDer対の既存の変形理論\cite{DasMandal2026}の複体である。
全充満部分圏への制限はこの複体と両立し、
外側の結合性障害が内側で境界になったとき、
その増強された障害は内側の $\operatorname{coker}[\delta,-]$ の一意の類へ移る。
標準六項補完では、これにより一次元の結合性障害線から
内在的対称性障害への自然な非零線形写像を得る。
この比較は通常の $HH^3$ 全体を対象とする写像ではなく、
導分の補正を記録した障害と、その標準補完の障害線上の帰結である。

\subsection{先行理論と証明の構成}

CGG積と可解Lie作用による量子化にはLecoutre--Sierra\cite{LecoutreSierra2019}、
AS正則代数の標準列とSerre双対性にはMinamoto--Mori\cite{MinamotoMori2011}、
次数付きtwistにはZhang\cite{Zhang1996}の理論を用いる。
Hochschild複体とそのGerstenhaber構造の導来不変性はKeller\cite{KellerDerived2003}、
dg Lie代数の変形関手の不変性はManetti\cite{Manetti2005}による。
Hochschild--Serre環の一般構成はLin--Zhang\cite{LinZhangSerre2025}に帰する。

本稿の新しい部分は、この全重み族の固定非零点の変形環、
反標準対角の全体性と具体的表示、それらによる族内同値分類、
ならびに量子VK四項成分の変形とGerstenhaber作用、
全重みの純粋部分への還元と通常hullの分解、対称性保持部分の全非被約構造の決定、
および導分付き障害を使った標準補完の自然な比較である。
可換基点のhullと $r>0$ のPoisson局所環は、公表稿IX\cite{VKPoisson34}を入力とする。
偶数端点の全局所環、全変形のOre実現、引戻し型の固定hullも本稿で証明する。
Ore拡大のNakayama公式はLiu--Wang--Wu\cite{LWWEndpoint37}、中心固定性はReyes--Rogalski\cite{ReyesRogalski2022}による。
四項成分のJordan型は、有理数上の有限Hochschild複体から厳密に決定し、
商座標を使う方法と全コサイクル空間を使う方法で独立に検証する。

公表稿I--IXのLO型関係代数、Prym幾何、相対potentialの問題は、
本稿の明示CGG族の分類とはそれぞれ別の問題として保持する。

\section{量子化族の定義と使用する変形定理}
\label{sec:quantum-setup-v35}

基礎体は $\mathbb C$ とする。$\operatorname{qgr}S$ は有限生成次数付き右加群の圏を
有限次元加群で割った商圏、$\mathcal O_S(i)$ は $S(i)$ の像を表す。
右加群間の射を左からの乗法で表す場合、$a:i\to j$ と $b:j\to k$ の合成は $b\star a$ である。
有限検算では経路順に $a\star b$ と書く反対の規約も用いるが、
その場合は同時に反対代数を取り、Hochschild群と変形環を変えない。

\subsection{有限CGG積}

$j,r\ge0$、$m=2j+r\ge3$、$d=m+r+1$ とし、
\[
 R_m=\mathbb C[x,y,z,w],\qquad
 \deg(x,y,z,w)=(1,1,1,m),\qquad q=xz-y^2
\]
と置く。次の重み零の導分を考える。
\begin{align*}
 N&=x\partial_y+2y\partial_z-q^jz^r\partial_w,\\
 X_\alpha&=\alpha x\partial_x+(\alpha+1)y\partial_y+(\alpha+2)z\partial_z
                      +(m\alpha+d)w\partial_w.
\end{align*}
$N$ は局所冪零であり、$[N,X_\alpha]=N$ である。
$R_m$ と同じ次数付きベクトル空間に
\begin{equation}
 a\star_{\alpha,s}b
  =\sum_{n\ge0}s^nN^n(a)\binom{X_\alpha}{n}(b),\qquad
 \binom{X}{n}=\frac{X(X-1)\cdots(X-n+1)}{n!}
 \label{eq:CGG-product-v35}
\end{equation}
を入れた代数を $\mathcal S_{j,r}(\alpha,s)$ と定める。
各 $a$ に対して和は有限なので、任意の複素数 $s$ で定義される。
下付き文字を省略した $\mathcal S$ と、上付き $(q)$ を付けた
$\mathcal S_{j,r}^{(q)}$ は、本稿では同じ $q^jz^r$ 型を意味する。

\begin{proposition}[量子化族の基本性質]
\label{prop:CGG-input-v35}
上の積は結合的であり、全 $\alpha,s$ について
$\mathcal S$ はnoetherianなAS正則四次元整域である。
\[
 H_{\mathcal S}(t)=\frac1{(1-t)^3(1-t^m)},\qquad \ell=m+3
\]
であり、$\ell$ はGorenstein parameterである。
$(\mathcal O,\ldots,\mathcal O(m+2))$ はfull strong exceptional collectionである。
$T_s=1+sN$ は自己同型で、全てのスカラー冪 $T_s^c$ も有限和として定まり、
Nakayama自己同型は
\[
 \mu=T_s^{(m+3)\alpha+m+r+3}
\]
である。特に $s\ne0$ の通常のCalabi--Yau代表は
\[
 \alpha_0=-\frac{m+r+3}{m+3}
\]
で与えられる。
\end{proposition}

これらは公表稿IX\cite[非線形アフィン量子化の節]{VKPoisson34}で証明した性質である。
結合性は $[N,X]=N$ に対するCGG恒等式、AS正則性は反復Ore表示、
Nakayama式の最後の平行移動項の消滅は次数付き自己同型の剛性による。
通常のCY性をPoissonテンソルの発散だけから推論したものではない。
標準列の存在はAS正則代数のtilting定理\cite{MinamotoMori2011}からも従う。

\subsection{同じ共鳴内のパラメータと有限代数}

$s\ne0$ を固定し、$T=1+sN$ と置く。
重み付きEuler場 $E$ を使うと $X_{\alpha+\beta}=X_\alpha+\beta E$ であり、
Vandermonde恒等式から、斉次 $a,b$ に対して
\[
 a\star_{\alpha+\beta,s}b
       =T^{\beta\deg b}(a)\star_{\alpha,s}b
\]
となる。これは左Zhang twistであり、$\operatorname{qgr}$ 圏を変えない。
さらに、有限例外列の自己準同型代数は、各射空間上の
\begin{equation}
 G_{ik}(a)=T^{-\beta i}(a),\qquad
 a\in\operatorname{Hom}(\mathcal O(i),\mathcal O(k))
 \label{eq:finite-twist-iso-v35}
\end{equation}
によって直接同型になる。実際、$a:i\to j$、$b:j\to k$ に対して
\[
 G_{jk}(b)\star_{\alpha+\beta,s}G_{ij}(a)
 =T^{-\beta i}(b)\star_{\alpha,s}T^{-\beta i}(a)
 =G_{ik}(b\star_{\alpha,s}a).
\]
従って全冪等元と任意の指定cornerを同時に保つ。
経路順の積規約では対応する式を $F_{ij}=T^{\beta j}$ と書ける。
また、補助次数 $(0,1,2,d)$ のスカラー作用は全ての非零 $s$ を互いに移す。
後の定理で $s=1$ またはCY値を選ぶ際には、この二つの同型を使う。

標準列の有限代数を
\[
 A_{j,r}(\alpha,s)=\operatorname{End}\left(
                 \bigoplus_{a=0}^{m+2}\mathcal O_{\mathcal S}(a)\right)
\]
と書く。次元は全パラメータで
\[
 \dim A_{j,r}=\binom{m+6}{4}+15
\]
である。特殊ファイバー $s=0$ を $\Lambda_m$ と書けば、
これは $X_m=\mathbb P_{\mathrm{stack}}(1,1,1,m)$ の同じ標準列の代数である。

\subsection{共鳴点のPoisson局所環と可換基点のhull}

$\mathscr P_m\subset H^0(X_m,\wedge^2T_{X_m})$ を
$[\pi,\pi]=0$ の全二次式で定義したPoissonスキームとする。
このイデアルは根基化しない。
$N_m=\binom{m+2}{2}$ と置く。以下が本稿の形式移送に必要な既知の入力である。

\begin{proposition}[可換基点と非零共鳴点の全形式環]
\label{prop:Poisson-local-input-v35}
\begin{enumerate}
\item $\Lambda_m$ のminiversal変形環は、HKRを接写像とする非標準的な同型により
\[
 R_{\Lambda_m}\simeq\widehat{\mathcal O}_{\mathscr P_m,0}
\]
となる。$HH^1(\Lambda_m)$ の次元は $N_m+9$ である。
\item $r>0$ では、式\eqref{eq:CGG-product-v35}の一次反対称部分が定める
射影Poisson点 $\pi$ は、接線共鳴成分の非零の稠密軌道に属する。
その全完備局所環と座標変換軌道の次元は
\[
 \widehat{\mathcal O}_{\mathscr P_m,\pi}
   \simeq\mathbb C[[z_1,\ldots,z_{N_m+8},\gamma]]/(\gamma^{j+1}),
 \qquad \dim\operatorname{Orb}(\pi)=N_m+8
\]
である。
\end{enumerate}
\end{proposition}

第一の主張は公表稿IX\cite[一般重みの二次的miniversal基底]{VKPoisson34}による。
Euler発散によるPoisson持上げを量子化すると、HKR接写像を持つ全Jacobi基底上の族を得る。
一次方向間のGerstenhaber括弧が定める独立な二次障害式がJacobi二次式と一致するため、
二次生成性とNakayamaの補題から
miniversal理想に隠れた高次生成元がないことが従う。
第二の主張は同稿の接線共鳴系列の全形式近傍の定理とアフィン持上げの命題による。
ここで $(c,D,[F])$ は、$\mathbf x=(x,y,z)$、
$p=(\nabla c-\mathbf x\times D)/3$ として
\[
 \Pi=p_1\partial_y\wedge\partial_z+p_2\partial_z\wedge\partial_x
                      +p_3\partial_x\wedge\partial_y,
 \qquad [\Pi+F\wedge\partial_w]
\]
を表すEuler商の標準座標であり、$F\sim F+g\mathbf x$（$g\in\mathbb C[x,y,z]_m$）とする。
その基準点は
\[
 c=xq,\quad D=-\frac{4m+3-6j}{2m}(x\partial_y+2y\partial_z),\quad
 F=q^j(0,yz^r/2,z^{r+1})
\]
であり、CGGの $N\wedge X_\alpha$ の射影像と、非零スカラーおよび重み付き座標変換を除いて一致する。
一般点の長さだけでなく、表示した\emph{全}完備局所環を使用する点が重要である。
$r=0$ の局所環は後の定理\ref{thm:endpoint-hull-v36}で別に証明する。

\subsection{多項式不変量、中心と共通の関数斜体}

\[
 c_r=(-1)^r\frac{(2r)!!}{(2r+1)!!},\qquad
 J=x^{r+1}w+q^j\sum_{k=0}^r\binom rk\frac{q^{r-k}y^{2k+1}}{2k+1},\qquad
 \Delta=\frac{J^2+c_r^2q^d}{x^{r+1}}
\]
と置く。$\Delta$ は多項式であり、既知の不変環の表示は
\[
 \ker N=\mathbb C[x,q,J,\Delta]
 \simeq\frac{\mathbb C[X,Q,U,V]}{(X^{r+1}V-U^2-c_r^2Q^d)}
\]
である\cite[非線形アフィン族の多項式中心]{VKPoisson34}。
これらの不変元同士の星積は通常積なので、以後の中心式における積は可換である。
通常CY値かつ $s\ne0$ では
\[
 Z(\mathcal S)=
 \begin{cases}
 \mathbb C&r\ge2,\\
 \mathbb C[q,J]&r=0,\\
 \mathbb C[\mathsf S,\mathsf P,\mathsf R]/
  (\mathsf R^2-\mathsf S\mathsf P^2-\frac{16}{81}\mathsf S^{m+3})&r=1,
 \end{cases}
\]
ただし $r=1$ の生成元は
\[
 \mathsf S=q^2\Delta,\qquad
 \mathsf P=2xJ\Delta^{(m+3)/2},\qquad
 \mathsf R=q(x^2\Delta+J^2)\Delta^{(m+3)/2}
\]
である。中心全体の生成性は $x$ に関する飽和と零重量単項式の完全な分類によって証明される。
局所化後の中心だけから全中心を推測することはしない。

最後に、$t=x^{-1}\star y$、$u=q/x^2$、$v=J/x^d$ とすると
\[
 [t,u]=2su,\qquad [t,v]=dsv,\qquad [u,v]=0.
\]
$d=2h+1$ と書き、$C=v^2/u^d$、$\eta=v/u^h$、
$\partial=\eta^{-1}t/s$ とすれば $[\partial,\eta]=1$ である。
逆に $u=\eta^2/C$、$v=\eta u^h$、$t=s\eta\partial$ なので、これらは斜体全体を生成する。
従って全ての非零族の次数零関数斜体は $\mathbb C(C)$ 上の第一Weyl斜体となる。
この一致は、本稿で求める射影圏の同値を意味しない。

\section{非零量子化の変形環と任意次数の障害}
\label{sec:fixed-quantum-hulls-v35}

本節では、量子化係数を非零に固定した後の変形問題を扱う。
対象は重み付き斉次座標環そのものの無次数変形ではなく、
標準例外列の有限次元自己準同型代数、従って対応する射影導来圏の変形である。
原点における変形環の同定を非零の複素数で直接評価することはしない。
代わりに、形式弧に沿って係数をLaurent級数体へ移し、
その上で完備局所環を再中心化する。

\subsection{結果}

$m=2j+r\ge3$、$j\ge0$、$r>0$ とし、前節のCGG代数を
$\mathcal S_{j,r}(\alpha,s)$ と書く。
標準例外列の有限代数を
\[
 A_{j,r}(\alpha,s)=
 \operatorname{End}\left(\bigoplus_{a=0}^{m+2}
              \mathcal O_{\mathcal S_{j,r}(\alpha,s)}(a)\right)
\]
と置く。以下の変形は、標数零の可換局所Artin基底上の平坦な結合代数変形である。
原始冪等元を持ち上げて固定しても同じ変形問題を得る。
これは冪等元の持上げと共役一意性による。
miniversal環はSchlessingerのhullの意味で用いる
\cite{Schlessinger1968}。

\begin{theorem}[固定非零量子化のminiversal環]
\label{thm:fixed-quantum-hull-v35}
$s\ne0$ とする。$A=A_{j,r}(\alpha,s)$ について
\[
 \dim HH^1(A)=1,\qquad
 \dim HH^2(A)=
 \begin{cases}0&j=0,\\1&j\ge1\end{cases}
\]
であり、そのminiversal変形環は非標準的に
\[
 \boxed{R_A\simeq\mathbf C[[\gamma]]/(\gamma^{j+1})}
\]
と同型である。$j=0$ の右辺は $\mathbf C$ と解釈する。
結論は標準Calabi--Yau値だけでなく、全ての $\alpha\in\mathbf C$ に成立する。
\end{theorem}

\begin{corollary}[Hochschild群の次元]
\label{cor:fixed-quantum-HH-v35}
$\epsilon_j=0$（$j=0$）、$\epsilon_j=1$（$j\ge1$）と置く。
上の有限代数について $HH^n(A)=0$（$n>3$）であり、
\[
 (\dim HH^0,\dim HH^1,\dim HH^2,\dim HH^3)=
 \begin{cases}
 (1,1,\epsilon_j,6+\epsilon_j)&m=3,\\
 (1,1,\epsilon_j,3+\epsilon_j)&m\ge4.
 \end{cases}
\]
\end{corollary}

\begin{proof}
可換基点ではスタック版HKRと非負線束の高次コホモロジー消滅により、
$HH^{>3}(A_0)=0$ である。非自明慣性成分も三次以外には寄与しない。
有限自由の相対Hochschild複体を特殊化する階数不等式からgeneric fiberでも消滅し、
非零係数での同型により全ての非零ファイバーへ移る。
この有限複体のEuler特性は一定であり、可換基点の次元
\[
 1,\quad N_m+9,\quad m^2+6m+26,\quad
 \begin{cases}41&m=3,\\10+\binom{m+5}{2}&m\ge4\end{cases}
\]
から、それぞれ $-6,-3$ となる。主定理と $HH^0(A)=\mathbf C$ を代入すればよい。
\end{proof}

この定理の有限代数は、標準列全体から作る。
四対象cornerの全変形環、または $\mathcal S$ の無次数全変形環について
同じ式を主張してはいない。
また、$HH^1(A)\ne0$ なので、miniversal環を自己同型を含む変形問題の
fineな表現環と呼ぶこともしない。

\begin{corollary}[任意次数で初めて現れる障害]
\label{cor:arbitrary-obstruction-v35}
$j\ge1$ とする。非零の一次変形類は
$\mathbf C[t]/(t^{j+1})$ まで延長できるが、
$\mathbf C[t]/(t^{j+2})$ へは延長できない。
従って $j=1$ では二次障害、$j\ge2$ では二次障害が消えた後の高次障害が現れる。
一方、$\mathbf C[[t]]$ 上の全一パラメータ形式変形は自明である。
\end{corollary}

\begin{proof}
非零一次類を持つ写像は $\gamma\mapsto at+O(t^2)$、$a\ne0$ と書ける。
その $(j+1)$ 乗の最初の項は $a^{j+1}t^{j+1}$ なので、延長可能な次数は正確に上記の通りである。
完備被約環 $\mathbf C[[t]]$ への写像では $\gamma$ の像は零である。
後述の形式切片補題\ref{lem:algebra-slice-v35}を用いると、
任意の形式変形には軌道変数だけが残る。
これは恒等写像を持ち上げる形式的な基底変換に対応するので、変形は自明である。
\end{proof}

\subsection{形式弧に沿った再中心化}

以下の補題では $k$ を標数零の体、$O=k[[h]]$、$K=k((h))$ とする。

\begin{lemma}[形式同型の弧に沿う移送]
\label{lem:arc-translation-v35}
$X,Y$ を有限型 $k$ スキーム、$x\in X(k)$、$y\in Y(k)$ とし、
\[
 \widehat{\mathcal O}_{X,x}\simeq\widehat{\mathcal O}_{Y,y}
\]
を固定する。$a\colon\operatorname{Spec}O\to X$ が閉点を $x$ に送る弧であるとき、
この同型は閉点を $y$ に送る弧 $b$ を定め、
\[
 \widehat{\mathcal O}_{X_K,a_K}\simeq
 \widehat{\mathcal O}_{Y_K,b_K}
\]
を誘導する。有限個の形式的に滑らかな変数を付加した形式スキームにも同じ主張が成立する。
\end{lemma}

\begin{proof}
局所的なアフィン埋込みを取り、形式同型とその逆を形式冪級数で表示する。
閉点を原点に移すと $a(h)$ の各座標は $hO$ に入る。
任意の形式冪級数 $F$ に対し
\[
 F(a(h)+z)=\sum_{\nu}
       \frac{\partial^\nu F(a(h))}{\nu!}z^\nu
\]
の各係数は $h$ 進収束する。
実際、固定した $z$ 次数 $|\nu|$ において、元の全次数 $n$ の項の
$h$ 進付値は少なくとも $n-|\nu|$ であり、$n\to\infty$ で無限大となる。

この代入は任意の局所Artin $K$ 代数にも意味を持つ。
その極大イデアルを $\mathfrak n$、$\mathfrak n^L=0$ とすると、
変位には高々 $L-1$ 個の冪零因子しか現れない。
それらの係数の $h$ 進付値は下から有界なので、同じ収束の議論が使える。
形式同型と逆写像の合成恒等式、および定義イデアルの等式は、この代入後も保たれる。
従って両辺のArtin $K$ 点の形式近傍が同型となり、完備局所環の同型を得る。
形式的に滑らかな変数も同じ議論で扱える。
\end{proof}

\begin{remark}
上の議論は $h$ 進に小さい点への評価である。
形式級数の収束半径を仮定せず、非零の複素数 $h$ を形式同型に代入する議論でもない。
\end{remark}

\begin{lemma}[斉次錐と太い滑らかな枝]
\label{lem:cone-arc-v35}
$P\subset V$ を斉次イデアルで定義されたアフィン錐とする。
$\pi\in P(k)$ において
\[
 \widehat{\mathcal O}_{P,\pi}
       \simeq k[[z_1,\ldots,z_d,\gamma]]/(\gamma^\ell)
\]
と仮定する。$p(h)=h\pi+O(h^2)\in P(O)$ に対して
\[
 \widehat{\mathcal O}_{P_K,p_K}
       \simeq K[[z_1,\ldots,z_d,\gamma]]/(\gamma^\ell)
\]
である。
\end{lemma}

\begin{proof}
$h$ は $K$ 上で可逆なので、錐のスカラー作用で $p(h)$ を
$p(h)/h=\pi+O(h)$ に移す。
この弧に $\pi$ での局所同型を適用する。
$O$ は被約なので、$\gamma^\ell=0$ を満たす $\gamma$ の像は零である。
滑らかな変数の像だけが $hO$ に入る弧となり、前補題によってそれらを平行移動すればよい。
\end{proof}

\subsection{代数構造スキームと最初のHochschild差分}

$E=\bigoplus_{a=0}^{m+2}\mathbf C e_a$ とし、
標準列代数の underlying $E$ 双加群を $V$ とする。
この双加群上の単位的な結合積のスキームを $\mathfrak A$ と書く。
基底変換群は
\[
 G=\prod_{a<b}\operatorname{GL}(e_bVe_a)
\]
である。相対正規化Hochschild複体を $C_E^\bullet$ とすると、
$\operatorname{Lie}G=C_E^1$、積の接空間は $Z_E^2$、
軌道の接空間は $B_E^2$ となる。

\begin{lemma}[軌道方向を分離する形式切片]
\label{lem:algebra-slice-v35}
$a\in\mathfrak A(k)$ に対応する有限代数のminiversal環を $R_a$ とし、
$b_a=\dim B_E^2(a)$ と置く。このとき
\[
 \widehat{\mathcal O}_{\mathfrak A,a}
       \simeq R_a[[u_1,\ldots,u_{b_a}]].
\]
\end{lemma}

\begin{proof}
$C_E^1$ で $Z_E^1$ の補空間を選び、その指数を形式的な基底変換として用いる。
$C_E^2$ では $B_E^2$ の補空間を選ぶ。
各次数で積の $B_E^2$ 成分を一意に消すことで、形式逆関数定理により
全結合積の形式近傍を、$b_a$ 個の軌道変数と横断切片に分離する。
切片上の族は全変形に対して形式的に滑らかで、接空間は
$Z_E^2/B_E^2=HH^2(a)$ なのでminiversalである。
\end{proof}

\begin{lemma}[最初の差分による上界]
\label{lem:first-differential-v35}
$A_h$ を有限自由 $O$ 代数とし、積を
$\mu_h=\mu_0+h\mu_1+O(h^2)$ と書く。
相対正規化Hochschild複体が各次数で有限自由であるとする。
このとき
\[
 \dim_K HH^n(A_h\otimes_O K)
 \le
 \dim_k H^n\bigl(HH^\bullet(A_0),[\mu_1,-]\bigr).
\]
\end{lemma}

\begin{proof}
特殊ファイバーの複体をコホモロジーと可縮な二項複体に線形分解する。
可縮部分の可逆行列式は $O$ 上でも単元なので、その部分を形式的Gaussian消去で除く。
残る有限自由複体の台は $HH^\bullet(A_0)[[h]]$ で、その微分は
\[
 h\delta_1+h^2\delta_2+\cdots,\qquad
 \delta_1=[\mu_1,-]
\]
となる。$K$ 上で微分を $h$ で割る。
その各微分の階数は特殊化 $h=0$ での階数以上だから、
各次数で核次元から像次元を引くと主張が従う。
必要な次数の有限切断にこの議論を適用すれば十分である。
\end{proof}

\subsection{CGG弧への適用}

まずCalabi--Yau値
\[
 \alpha=-\frac{m+r+3}{m+3}
\]
を固定し、$s=h$ とした有限代数を $A_h$ とする。
CGG積は各斉次成分上有限和なので、これは $\mathbf C[h]$ 上の
同じ underlying $E$ 双加群を持つ平坦な代数族である。
特殊ファイバーは可換な重み付き射影スタックの標準列代数 $A_0=\Lambda_m$ である。
一次類は、非零定数倍を除いて、接線共鳴成分の非零代表点 $\pi$ である。
非零定数倍は錐の作用で局所環を変えない。

既に証明した可換基点のminiversal環の同定は、HKRを接写像として
\[
 R_{A_0}\simeq\widehat{\mathcal O}_{\mathscr P_m,0}
\]
を与える。形式切片補題を合わせると
\[
 \widehat{\mathcal O}_{\mathfrak A,A_0}
 \simeq
 \widehat{\mathcal O}_{\mathscr P_m,0}[[u_1,\ldots,u_{b_0}]]
\]
となる。弧 $A_h$ のPoisson座標は $h\pi+O(h^2)$ である。
接線成分定理と錐の弧補題により、そのgeneric fiberにおいて
\begin{equation}
\label{eq:generic-framed-local-v35}
 \widehat{\mathcal O}_{\mathfrak A_K,A_h\otimes K}
 \simeq
 K[[u_1,\ldots,u_{b_0},z_1,\ldots,z_{N_m+8},\gamma]]/
       (\gamma^{j+1}).
\end{equation}

次に、この滑らかな変数の全てが基底変換に対応することを示す。
$HH^1(A_0)=H^0(X_m,T_{X_m})$ の次元は $N_m+9$ であり、
HKRの下で $[\mu_1,-]$ は $[\pi,-]$ に対応する。
接線成分の被約部分には次元 $N_m+8$ の座標変換軌道が稠密にあるので、
この写像の核は一次元である。$HH^0(A_0)=\mathbf C$ からの差分は零である。
従って最初の差分補題から
\[
 \dim_K HH^1(A_h\otimes K)\le1
\]
を得る。

逆向きの不等式には、既知のCGG自己同型 $T_h=1+hN$ を用いる。
有限代数の各成分上で
\[
 D_h=\frac1h\log(1+hN)
       =N-\frac h2N^2+\frac{h^2}{3}N^3-\cdots
\]
は有限和であり、$A_h$ の $O$ 線形導分である。
実際、$T_h^t$ が全てのスカラー $t$ に対して自己同型なので、
$t=0$ で微分して $\log T_h$ を得る。
$D_h$ の特殊化 $N$ は、$X_m$ 上の非零の大域ベクトル場を与え、
$HH^1(A_0)$ において非零である。

$D_h$ が $K$ 上で内的だと仮定し、$D_h=\operatorname{ad}(b(h))$ と書く。
$b(h)$ の負次数の最初の係数は $A_0$ の中心に入る。
$Z(A_0)=\mathbf C$ なので、そのスカラー項を差し引ける。
負次数を有限回除けば $b(h)\in A_h$ としてよい。
特殊化すると $N=\operatorname{ad}(b(0))$ となり矛盾する。
よって
\[
 \boxed{\dim_K HH^1(A_h\otimes K)=1}.
\]

標準列代数の中心は、どのファイバーでもスカラーのみである。
実際、中心元は冪等元と可換なので対角成分に入り、隣接する非零射によって
全対角係数が等しくなる。
従って $\dim\operatorname{im}(d^0)=m+2$ は一定である。
軌道次元は
\[
 b_a=\dim C_E^1-\dim HH^1(A_a)-(m+2)
\]
だから、generic fiberでは
\[
 b_\eta=b_0+N_m+8.
\]
式\eqref{eq:generic-framed-local-v35}から軌道方向を形式切片で除くと、
generic fiberのminiversal環は
\[
 K[[\gamma]]/(\gamma^{j+1})
\]
となる。ここで環の同型型を確認しておく。
generic fiberのhullを $R_\eta$ とすると、形式切片補題によって
\[
 R_\eta[[v_1,\ldots,v_{b_\eta}]]
 \simeq K[[z_1,\ldots,z_{b_\eta},\gamma]]/(\gamma^{j+1})
\]
である。従って $\dim R_\eta=0$、$j=0$ では埋込次元零、
$j\ge1$ では埋込次元一となる。
後者では $R_\eta\simeq K[[t]]/(t^n)$ と書ける。
Hilbert--Samuel重複度は形式冪級数変数の付加で変わらないので、
両辺の重複度を比較して $n=j+1$ を得る。

最後に、非零の量子化係数の同型を使う。
補助次数 $(0,1,2,m+r+1)$ のスカラー作用は $K$ 上で
$A_h\otimes K\simeq A_1\otimes K$ を与える。
有限代数のHochschildコホモロジーと形式切片は体拡大と可換である。
後者については、有限次元の積の空間で選んだ線形切片と結合式は
全て $\mathbf C$ 上に定義されているので、その交叉スキームを体拡大して完備化すればよい。
従って $A_1$ のminiversal環は埋込次元零（$j=0$）または一（$j\ge1$）である。
後者を $\mathbf C[[\gamma]]/(f)$ と書くと、体拡大後の消滅次数から
$\operatorname{ord}_\gamma f=j+1$ を得る。
ゆえに $f$ は $\gamma^{j+1}$ の単元倍であり、定理の環が従う。
全非零 $s$ も同じである。

全 $\alpha$ への拡張は、有限 $\mathbf Z$ 代数の明示的な同型でも確認できる。
$T=1+sN$ とし、同じ $(j,r,s)$ に対してパラメータを $\alpha$ から
$\alpha+\beta$ に変える。斉次元 $a,b$ について、二項係数のVandermonde恒等式から
\[
 a\star_{\alpha+\beta}b
       =T^{\beta\deg b}(a)\star_\alpha b
\]
を得る。$N$ は局所冪零なので、任意の複素数を指数に持つ $T$ の冪も
各斉次成分上の有限和である。
経路順に $a\in\operatorname{Hom}(i,j)$、
$b\in\operatorname{Hom}(j,k)$ の積を書く規約で、
\[
 F_{ij}(a)=T^{\beta j}(a)
\]
と定める。$T$ が $\star_\alpha$ の自己同型であることを使うと
\[
 \begin{aligned}
 F_{ij}(a)\star_{\alpha+\beta}F_{jk}(b)
 &=T^{\beta(k-j)}T^{\beta j}(a)\star_\alpha T^{\beta k}(b)\\
 &=T^{\beta k}(a\star_\alpha b)
 =F_{ik}(a\star_\alpha b).
 \end{aligned}
\]
従って $F$ は全ての頂点冪等元を固定する有限代数の同型である。
通常の右加群の自己準同型環で合成を $b\star a$ と書く場合には、
対応する明示同型は始点を用いた
\[
 G_{ij}(a)=T^{-\beta i}(a)
\]
である。実際、$\deg a=j-i$ なので
\[
 \begin{aligned}
 G_{jk}(b)\star_{\alpha+\beta}G_{ij}(a)
 &=T^{\beta(j-i)-\beta j}(b)\star_\alpha T^{-\beta i}(a)\\
 &=T^{-\beta i}(b\star_\alpha a)
 =G_{ik}(b\star_\alpha a).
 \end{aligned}
\]
従ってこちらの規約でも全ての頂点冪等元を固定する同型を得る。
いずれの規約でも、任意の頂点部分集合のcornerを同時に保つ。
これにより全 $\alpha$ について結論を得る。

\subsection{射影導来圏の区別}

Hochschild複体の $B_\infty$ 高次構造は導来Morita同値で保存される
\cite[Theorem~3.2 and its corollary]{KellerDerived2003}。
標数零では、そのshiftしたdg Lie構造のMaurer--Cartan変形関手は
quasi-isomorphismで不変である\cite[Corollary~3.2]{Manetti2005}。
従ってminiversal環の同型型も保存される。
ここでは有限代数を次数零のdg代数と見てよい。
同じ基底上の最小 $A_\infty$ 構造では、次数の理由から $m_2$ 以外の演算は現れない。
また、標準tilting対象は自己拡大を持たず、Artin変形に沿って持ち上がるので、
その有限代数の変形は対応する導来圏の変形を記述する。

\begin{corollary}[同じ重みの異なる共鳴は導来非同値]
\label{cor:resonance-derived-inequivalence-v35}
$m$ を固定する。$0\le j\le\lfloor(m-1)/2\rfloor$ に対して、
$s\ne0$ の射影導来圏
\[
 D^b\bigl(\operatorname{qgr}\mathcal S_{j,m-2j}(\alpha,s)\bigr)
\]
は、異なる $j$ に対して互いに $\mathbf C$ 線形導来Morita同値でない。
\end{corollary}

\begin{proof}
標準列のtilting同値により $\operatorname{Perf}(A_{j,m-2j})$ と同一視する。
同値ならminiversal環の長さも等しいが、それは $j+1$ なので $j$ が一致する。
\end{proof}

この結論は、多項式中心の相違を射影圏の不変量と見なしたものではない。
全ての族が同じ次数零Weyl関数斜体を持つこととも両立する。
奇数重み $m=2\ell-1$ の最後の共鳴では、固定した射影圏のminiversal環の長さも
$\ell$ となり、標準CY代表の中心は $D_{2\ell+3}$ 型である。

\subsection{論証の射程}

この移送に必要なのは、可換基点のhullの原点形式同型、そのHKR接写像、
CGG族の多項式性、非零係数の同型、および接線点の全局所環である。
CGG積が特定のKontsevich公式と全次数で同じであるとは仮定しない。
全局所環の式を知っているため、弧が持つ高次Poisson補正は
被約な滑らかな枝の中にとどまり、そのgeneric fiberの厚みを変えない。

ここまでの定理は $r>0$ の接線成分に関する。
$r=0$ の有理型は次の小節で、四対象cornerは後の節で別に扱う。
無限次元斉次環の全無次数変形と、自己同型を含む変形スタックの完全な表示は本稿の範囲外である。
補助的な低重みの導分行列の階数検査は、全整数に対するこの形式移送の証明を置き換えない。

\subsection{低重みにおける有理数上の導分検証}

標準列のquiverでは、次数一の三生成元の矢印が $m+2$ 箇所にあり、
次数 $m$ の生成元の矢印が三箇所にある。
従って冪等元を固定する導分の未知数は
$9(m+2)+3(N_m+1)$ 個である。
三つの二次関係と三つの次数 $m+1$ の混合関係にLeibniz則を適用すると、
その未知数に関する有理数行列を得る。
順序付きCGG単項式を用いて全関係を正確に評価し、
有理数Gaussian消去で導分の次元を求めた結果は次の通りである。
ここでは $s=1$、$\alpha=-(m+r+3)/(m+3)$ としている。
\[
\begin{array}{c|c|c|c|c|c}
(m,j,r)&\text{未知数}&\text{方程式}&\operatorname{rank}_{\mathbf Q}
 &\dim\operatorname{Der}_E&\dim HH^1\\\hline
(3,0,3)&78&180&72&6&1\\
(3,1,1)&78&180&72&6&1\\
(4,0,4)&102&234&95&7&1\\
(4,1,2)&102&234&95&7&1\\
(5,2,1)&129&294&121&8&1
\end{array}
\]
可換基点 $(m,s)=(3,0)$ では、同じ $180\times78$ 行列の階数は $54$ となり、
$\dim\operatorname{Der}_E=24$、$\dim HH^1=19$ を回復する。
各非零例では、頂点ごとの内的導分が $m+2$ 次元の部分空間を張り、
$s^{-1}\log(1+sN)$ がこれらから独立な導分であることも、行列上で検証した。
さらに $\alpha=0,-1$、$s=-1/2$、および $m=3$ のもう一つの座標表示を含む
計十ケースを、標準ライブラリのみを用いる付属コードと証明書に収録した。
検証モードは証明書を生成せず、再計算の全JSONを既存証明書と比較する。

\subsection{偶数重みの端点は滑らかな変形空間を持つ}
\label{subsec:endpoint-hulls-v36}

接線成分定理の仮定 $r>0$ を外すと、結論は変わる。
$m=2j\ge4$ の端点 $r=0$ では、一変数の非被約環の代わりに
$m+2$ 個の自由な変形変数が現れる。
この相違は、固定した非零の量子化を取った後にも残る。
以下では $N_d=\binom{d+2}{2}$ と置く。

\begin{theorem}[端点の全局所環と固定量子化のhull]
\label{thm:endpoint-hull-v36}
$m=2j\ge4$ とし、
\[
 q=xz-y^2,\quad c_*=xq,\quad
 a_0=-\frac{m+3}{2m},\quad
 D_0=a_0(x\partial_y+2y\partial_z),\quad
 F_0=q^j(0,y/2,z)
\]
と置く。Poissonスキームの点 $Q_0=(c_*,D_0,[F_0])$ について、
\[
 \widehat{\mathcal O}_{\mathscr P_m,Q_0}
 \simeq\mathbf C[[t_1,\ldots,t_{N_{m+1}+8}]].
\]
この点の重み付き座標変換軌道の次元は $N_m+8$、
無限小安定化群の次元は一である。
さらに $s\ne0$、$\alpha\in\mathbf C$ に対し、標準全列の有限代数
$A=A_{j,0}(\alpha,s)$ のminiversal環とHochschild群は
\[
 \boxed{R_A\simeq\mathbf C[[u_1,\ldots,u_{m+2}]]},\qquad
 (\dim HH^0,\dim HH^1,\dim HH^2,\dim HH^3)
       =(1,1,m+2,m+5),
\]
および $HH^{>3}(A)=0$ を満たす。
\end{theorem}

\begin{proof}
最初に全Poissonスキームの局所環を求める。
基底スキーム $D(c)=0$ の二変数切片を
\[
 q_\gamma=xz-y^2+\gamma z^2,\quad
 c_\gamma=xq_\gamma,\quad
 D_{a,\gamma}=a((x+2\gamma z)\partial_y+2y\partial_z)
\]
と取る。七つの軌道方向とこの切片による形式的表示は、$a_0\ne0$ なので
接線成分の場合と同様に有効である。
新しい基底座標
\[
 \kappa=\frac{m+3+2ma}{3}
\]
を用いると、全混合写像は正確に
\begin{equation}
\label{eq:endpoint-mixed-v36}
 M_{\kappa,\gamma}(F)
 =xL_\gamma(F)+\kappa K_\gamma(F),
 \quad
 L_\gamma(F)=\frac{m+3}{2}F(q_\gamma)
                 -q_\gamma\operatorname{div}F,
 \quad
 K_\gamma(F)=q_\gamma F_x-\frac x2F(q_\gamma)
\end{equation}
となる。両写像はEuler商
$W_m=S_{m+1}^{\oplus3}/S_m(x,y,z)$ 上で定義される。

\medskip\noindent
\textbf{混合核と一つの余核方向。}
任意の非退化二次形式 $q$ に対し、上の $L$ は
\[
 W_m\longrightarrow S_{m+2}
\]
の写像である。その核は
\begin{equation}
\label{eq:endpoint-hamiltonians-v36}
 \ker L=\{[\nabla q\times\nabla H]:H\in S_{m+1}\},
 \qquad \dim\ker L=N_{m+1},
\end{equation}
であり、余核は $q^{j+1}$ の類で生成される一次元空間である。
この主張を次数に依存しない計算で確かめる。
線形座標変換により $q=x^2+y^2+z^2$ としてよい。
$k=j+1$、$n=m+2=2k$ と置く。
Fischer双線形形式 $\langle f,g\rangle=f(\partial)g|_0$ に対し
\[
 \langle q^k,L(F)\rangle
 =\sum_i\left\langle
 (m+3)\partial_i q^k-x_i\Delta q^k,F_i\right\rangle=0
\]
である。実際、$\Delta q^k=2k(2k+1)q^{k-1}$ かつ $m+3=2k+1$ である。
また $\langle q^k,q^k\rangle\ne0$ なので $q^k$ は像に入らない。

逆に、斉次多項式の調和分解
\[
 S_n=\bigoplus_{e=0}^{k}q^e\mathcal H_{n-2e},
 \qquad \mathcal H_\ell=\ker(\Delta:S_\ell\to S_{\ell-2})
\]
を用いる。
この分解は
$\Delta(q^e h_\ell)=2e(2\ell+2e+1)q^{e-1}h_\ell$
から次数に関する帰納法で得られる。
$H=q^e h_\ell$、$\ell=n-2e$ に対し
\[
 L(\nabla H)=\ell(\ell+1)H.
\]
従って $\ell>0$ の全成分は像に入り、余核は正確に一次元である。

$\nabla q\times\nabla H$ は発散零で $q$ に接するので $L$ の核に入る。
この写像はEuler商上で単射である。
実際、このベクトル場が $g\,(x,y,z)$（$g\in S_m$）に等しければ、発散を取って
$(m+3)g=0$、従ってベクトル場は零となる。
さらに $dq\wedge dH=0$ なら、$\nabla q$ の成分が正則列をなすため
$dH=f\,dq$ と書ける。Euler公式から $(m+1)H=2qf$ を得るので $q\mid H$ であり、
同じ議論を $H/q$ に繰り返せる。
$\deg H=m+1$ は奇数だから最後に一次式へ達し、その式も $q$ で割れるため零である。
余核の次元から
\[
 \dim\ker L
 =3N_{m+1}-N_m-(N_{m+2}-1)=N_{m+1}
\]
となり、これで式\eqref{eq:endpoint-hamiltonians-v36}も証明された。
$q_\gamma$ は $x'=x+\gamma z$ という可逆な座標変換で $x'z-y^2$ に直るため、
核と余核の以上の記述は $\gamma$ を含む局所基底上でも自由に成立する。

\medskip\noindent
\textbf{非零混合点での全方程式。}
$\kappa=0$ では $M=xL_\gamma$ であり、その核の基底には
実際の列 $\nabla q_\gamma\times\nabla H$ を取れる。
像の可逆小行列を選び、$\kappa,\gamma$ を含む局所環上でその変数を消去する。
消去後の全残差は $\kappa=0$ 上で恒等的に零になるため、いずれも $\kappa$ で割れる。

$H_0=q^j y$ とすると
\[
 F_{H_0}=\nabla q\times\nabla H_0=q^j(-x,0,z)
       =2F_0-q^j(x,y,z),
 \qquad K_0(F_{H_0})=-xq^{j+1}.
\]
最後の元は $xL_0$ の像に入らない。
従って $[F_0]$ における $\kappa$ 微分は混合余核で非零であり、
消去後の残差の一つは $\kappa$ と単元の積である。
全残差が $\kappa$ で割れることと合わせると、残差イデアルは正確に $(\kappa)$ となる。
$\gamma$ と $N_{m+1}$ 個の核変数には関係が残らない。
七軌道変数を戻して、局所環は正則で次元 $N_{m+1}+8$ と分かる。
これはZariski接空間だけの計算ではなく、全次数の残差イデアルの決定である。

\medskip\noindent
\textbf{軌道と安定化群。}
基底 $(c_*,D_0)$ の軌道は七次元である。
混合項を $[F_H]/2$ としたときの基本一形式は
\[
 \alpha_H=\frac12\{2q\,d(xw+H)-(m+1)(xw+H)\,dq\}.
\]
実際、混合部分については
$(\nabla q\times\nabla H)\times\mathbf x
 =2q\nabla H-(m+1)H\nabla q$ から従い、
$w$ を含む部分は $c_*,D_0$ の表示から従う。
混合核の表示 $H\mapsto[\nabla q\times\nabla H]$ の下で、
$w$ の次数 $m$ の移動は、非零の定数倍を除いて $H\mapsto H+xg$（$g\in S_m$）を与える。
これは $N_m$ 次元の部分空間である。
$H_0=q^j y$ は $x$ で割れないので、$w$ の倍率変更から来る
$[F_0]$ 方向はこの部分空間と独立である。
従って軌道の次元は少なくとも $7+N_m+1=N_m+8$ である。
一方、座標変換群のLie代数の次元は $N_m+9$ であり、
\[
 \mathcal N_j=x\partial_y+2y\partial_z-q^j\partial_w
\]
はこのPoissonテンソルを固定する非零の大域ベクトル場である。
従って軌道の次元は正確に $N_m+8$、安定化群のLie代数は一次元となる。

\medskip\noindent
\textbf{固定した量子化への移送。}
まず標準CY値 $\alpha=-1$ で $s=h$ としたCGG族を用いる。
可換基点のHKR接写像を持つhullの同定、および形式弧に沿った再中心化補題を適用する。
Poisson座標は $h\pi+O(h^2)$ であり、錐のスカラー作用で
$\pi+O(h)$ に移すと、今証明した滑らかな形式近傍の内部にある。
従ってgeneric fiberにおける代数構造スキームの完備局所環は滑らかで、
その次元は $b_0+N_{m+1}+8$ となる。
ここで $b_0=\dim B_E^2(A_0)$ である。

Poisson安定化群の計算と最初のHochschild差分によって $HH^1$ の次元は高々一であり、
明示的な外的導分 $h^{-1}\log(1+h\mathcal N_j)$ によって逆の不等式を得る。
従ってgeneric fiberで $HH^1=1$、軌道次元は $b_0+N_m+8$ である。
形式切片補題によりhullは正則で、その次元は
\[
 (b_0+N_{m+1}+8)-(b_0+N_m+8)=N_{m+1}-N_m=m+2
\]
となる。
非零 $s$ のスカラー同型と $\alpha$ の明示的なZhang同型で、全パラメータに結論を移せる。
$HH^{>3}$ の消滅とEuler特性 $-3$ により、残りのHochschild次元も従う。
\end{proof}

\begin{remark}[端点の変形の意味]
端点は有理型 $R(2,m+1)$ の滑らかな点であり、全categorical hullは $m+2$ 次元である。
次節では、この抽象的なhullの全ての方向を、有限個の係数を持つ通常CYのAS正則座標環で実現する。
固定された無限次元次数付き環の全ての非次数付き変形を分類する主張とは区別する。
\end{remark}

\section{偶数端点の全変形を実現するCalabi--Yau座標環}
\label{sec:endpoint-quantization-v37}

偶数端点の滑らかなhullは、形式的な存在だけでなく、
有限個の係数を持つAS正則座標環によって実現できる。
構成の要点は、三変数代数の正規元とNakayama自己同型を一致させることである。
その後の混合項は、任意次数の多項式を許すOre導分として組み込める。

\subsection{二次代数と任意の混合項}

$s,\gamma$ をパラメータとし、
\[
 D=1-\gamma s^2,\qquad
 a=\frac{s^2}{D},\qquad b=\frac{s}{D},\qquad
 c=\frac{1+\gamma s^2}{D}
\]
と置く。以下の族は $D$ が可逆な開集合上に定義する。
次数一の生成元 $x,y,Z$ を持つ代数 $B_{\gamma,s}$ を
\begin{equation}
\label{eq:endpoint-B-relations-v37}
\begin{aligned}
 yx&=cxy-bx^2-2\gamma b\,xZ,\\
 Zx&=cxZ-2bxy+ax^2,\\
 Zy&=yZ+axy-bxZ
\end{aligned}
\end{equation}
で定義する。大文字 $Z$ は、CGG表示で用いた $z$ と区別するためである。
自己同型と二次の元を
\begin{equation}
\label{eq:endpoint-sigma-Q-v37}
\begin{aligned}
 \sigma(x)&=x,\\
 \sigma(y)&=cy-bx-2\gamma bZ,\\
 \sigma(Z)&=cZ-2by+ax,\\
 Q&=\frac{xZ-sxy}{D}-y^2+\gamma Z^2
\end{aligned}
\end{equation}
と置く。

\begin{proposition}[三変数部分の構造]
\label{prop:endpoint-base-v37}
環 $\mathbb C[s,\gamma,D^{-1}]$ 上で、$B_{\gamma,s}$ は
$x^i y^k Z^\ell$ を基底に持つ次数付き自由加群である。
$x$ は正則な正規元であり、
\[
 ux=x\sigma(u),\qquad
 B_{\gamma,s}/(x)\simeq\mathbb C[s,\gamma,D^{-1}][y,Z]
\]
が成立する。ここで商の式は係数環の任意の基底変換後にも成立する。
$Q$ は中心元であり、$\sigma(Q)=Q$ である。

各複素数点 $D\ne0$ のファイバーはnoetherianなAS正則三次元整域で、
Hilbert級数 $(1-t)^{-3}$、Gorenstein parameter $3$、
Nakayama自己同型 $\mu_B=\sigma$ を持つ。
\end{proposition}

\begin{proof}
式\eqref{eq:endpoint-B-relations-v37}を左辺から右辺へ簡約する。
次数付き辞書式順序 $Z>y>x$ で簡約は停止し、唯一の重なりは $Zyx$ である。
\[
 c^2-4\gamma b^2=1,\qquad
 c-1=2\gamma a,\qquad a(c+1)=2b^2
\]
を代入すると、この重なりの二通りの簡約結果は一致する。
従ってDiamond lemmaにより主張した基底を得る。
式\eqref{eq:endpoint-sigma-Q-v37}が関係式を保存することは直接代入で分かり、
逆写像は $\sigma$ の式で $s$ を $-s$ に置き換えたものである。
正規性、$Q$ の中心性、および $\sigma(Q)=Q$ も三つの関係式から従う。
これらは $\mathbb Q(\gamma,s)$ における恒等式であり、
分母は $D$ の冪だけなので、表示した係数環上で成立する。

AS正則性とNakayama自己同型を、二種類の反復Ore表示で確認する。
$\gamma\ne0$ のとき、$\rho^2=\gamma$ とし、
\[
 t=\frac{1-s\rho}{1+s\rho},\qquad
 u=y+\rho Z+\frac{x}{2\rho},\qquad
 v=y-\rho Z-\frac{x}{2\rho},\qquad
 k=\frac{t^{-1}-t}{4\gamma}
\]
と置く。$D\ne0$ により $t$ は定義され、非零である。
可逆な線形変換によって、関係式と中心元は
\begin{equation}
\label{eq:endpoint-diagonal-Ore-v37}
 ux=t xu,\qquad vx=t^{-1}xv,\qquad
 uv-vu=kx^2,\qquad Q=-uv-\frac{t}{4\gamma}x^2
\end{equation}
となる。
従って $B=\mathbb C[x][u;\tau][v;\tau_v,\partial]$ であり、
\[
 \tau(x)=tx,\quad
 \tau_v(x)=t^{-1}x,\quad \tau_v(u)=u,\quad
 \partial(x)=0,\quad\partial(u)=-kx^2.
\]
量子平面のNakayama自己同型は $x\mapsto t^{-1}x$、$u\mapsto tu$ である。
Ore拡大の定理\cite[Theorem~3.3]{LWWEndpoint37}から
$\mu_B(x)=x$、$\mu_B(u)=tu$ を得る。
Nakayama自己同型は中心を各点で固定する
\cite[Lemma~4.5]{ReyesRogalski2022}。
式\eqref{eq:endpoint-diagonal-Ore-v37}の $Q$ を固定することと整域性から、
$\mu_B(v)=t^{-1}v$ も従う。これは式\eqref{eq:endpoint-sigma-Q-v37}の $\sigma$ である。

$\gamma=0$ のときは
\[
 A=\mathbb C[x][y;\delta],\quad\delta(x)=-s x^2,\qquad
 B=A[Z;\tau,\partial]
\]
と書ける。ここで
\[
 \tau(x)=x,\quad\tau(y)=y-sx,\qquad
 \partial(x)=-2sxy+s^2x^2,\quad\partial(y)=s^2xy.
\]
Jordan平面のNakayama自己同型は
$\mu_A(x)=x$、$\mu_A(y)=y-2sx$ である
\cite[\S4.2]{LWWEndpoint37}。
従って $\mu_B(x)=x$、$\mu_B(y)=y-sx$ となり、
中心元 $Q=xZ-sxy-y^2$ の固定性から
$\mu_B(Z)=Z-2sy+s^2x$ を得る。これも $\sigma$ である。
両Ore表示はnoetherian性、整域性、AS正則性を与える。
次数一のOre拡大を三回取った表示からGorenstein parameterは $3$ である。
\end{proof}

一般に、次の補題が混合項の量子化を与える。
ここでOre拡大の規約は $wu=\sigma(u)w+\delta(u)$ とする。

\begin{lemma}[中心元を持つOre拡大]
\label{lem:central-Ore-CY-v37}
$B$ を連結次数付きnoetherian AS正則整域とする。
正則な斉次正規元 $x$ の正規化自己同型を
$ux=x\sigma(u)$ で定め、
$\mu_B=\sigma$ および $B/(x)$ の可換性を仮定する。
斉次元 $H\in B$ について $\deg H>\deg x$ とし、
\[
 \delta_H(u)=x^{-1}(uH-Hu)
\]
と置く。右辺は局所化を必要とせず、$x$ による一意の左割算を意味する。
この写像は $B$ に値を取る $\sigma$ 導分である。
\[
 S_H=B[w;\sigma,\delta_H],\qquad
 \deg w=\deg H-\deg x
\]
は通常のCalabi--Yau代数であり、$J=xw+H$ は中心元となる。
\end{lemma}

\begin{proof}
$B/(x)$ が可換なので $uH-Hu\in xB$ である。
正規性と $x$ の正則性を使うと
\[
 \delta_H(uv)=\sigma(u)\delta_H(v)+\delta_H(u)v
\]
を得る。また
\[
 \delta_H(x)=H-\sigma(H),\qquad\delta_H(H)=0.
\]
$u\in B$ に対して $x\sigma(u)=ux$ を用いれば
\[
 (xw+H)u=uxw+uH=u(xw+H)
\]
となる。さらに
\[
 w(xw+H)=xw^2+(H-\sigma(H))w+\sigma(H)w=(xw+H)w.
\]
従って $J$ は中心元である。

Ore拡大の定理\cite[Theorem~3.3]{LWWEndpoint37}より、
$S_H$ はtwisted Calabi--Yauで、
$\mu_{S_H}|_B=\sigma^{-1}\mu_B=\operatorname{id}$ となる。
Nakayama自己同型の中心固定性から
\[
 x\mu_{S_H}(w)+H=\mu_{S_H}(J)=J=xw+H.
\]
$x$ はOre拡大内でも正則なので $\mu_{S_H}(w)=w$ である。
従ってNakayama自己同型は恒等写像であり、通常のCalabi--Yau性を得る。
\end{proof}

\begin{theorem}[任意の混合項の代数的量子化]
\label{thm:endpoint-Ore-family-v37}
$m\ge1$、$H\in(B_{\gamma,s})_{m+1}$ とし、
\begin{equation}
\label{eq:endpoint-S-family-v37}
 S_{\gamma,s,H}
 =B_{\gamma,s}[w;\sigma,\delta_H],\qquad \deg w=m
\end{equation}
と定める。これは $H$ のPBW係数と $s,\gamma,D^{-1}$ に代数的に依存する族であり、
$x^i y^k Z^\ell w^e$ を基底に持つ。
各複素数点でのファイバーはnoetherianなAS正則四次元整域で、
\[
 H_S(t)=\frac1{(1-t)^3(1-t^m)},\qquad \ell=m+3,\qquad \mu_S=\operatorname{id}.
\]
$Q$ と $J=xw+H$ は中心元である。
従って標準列 $(\mathcal O_S,\ldots,\mathcal O_S(m+2))$ は
full strong exceptional collectionである。
\end{theorem}

\begin{proof}
命題\ref{prop:endpoint-base-v37}と補題\ref{lem:central-Ore-CY-v37}を適用する。
$Q$ は $B$ の中心にあり、$\sigma(Q)=Q$、$\delta_H(Q)=0$ なので、
全代数 $S$ の中心にも入る。
Ore拡大は左 $B$ 加群として $1,w,w^2,\ldots$ を基底に持つため、
PBW基底とHilbert級数を得る。
Gorenstein parameterは次数 $m$ のOre変数を加えることで $3+m$ となる。
標準列についてはAS正則代数のtilting定理\cite{MinamotoMori2011}を適用する。

係数環上では、PBW左割算による $\delta_H$ の係数は $H$ の係数に線形で、
分母には $D$ の冪だけが現れる。従ってこれは実際の有限型の代数族である。
任意の基底変換後も各斉次成分は同じPBW基底で自由となる。
\end{proof}

\begin{remark}
$x$ の正規性を使う場所は三変数部分 $B$ である。
全代数では $wx=xw+H-\sigma(H)$ なので、一般の $H$ に対して
$x$ は正規元ではない。$\sigma$ が全代数の自己同型に延長するとも主張しない。
また、冪零元を持つ係数環上では「AS正則」を絶対的な大域次元の意味で用いない。
上の主張は、代数族の次数ごとの自由性と、各体上のファイバーのAS正則性である。
\end{remark}

\subsection{端点CGG代数との明示的な一致}

ここから $m=2j\ge4$、$n=m+1=2j+1$、$\epsilon=(-1)^j$ とする。
$\gamma=0$ で
\[
 H_*=Q^j y
\]
を選ぶと、$S_{0,s,H_*}$ は端点のCY代表
$\mathcal S_{j,0}(-1,s)$ と同型になる。
生成元の対応は
\begin{equation}
\label{eq:endpoint-CGG-identification-v37}
 x\longmapsto x,\qquad y\longmapsto y,\qquad
 Z\longmapsto z+sy,\qquad w\longmapsto w.
\end{equation}
この対応で $Q\mapsto q=xz-y^2$、$J\mapsto xw+q^j y$ である。
実際、CGG積から得られる混合関係は正確に
\begin{equation}
\label{eq:endpoint-CGG-mixed-v37}
\begin{aligned}
 wx&=xw+s xQ^j,\\
 wy&=(y-sx)w,\\
 wZ&=(Z-2sy+s^2x)w-sQ^j(Z-sy)
\end{aligned}
\end{equation}
であり、$\delta_{Q^j y}$ の定義と一致する。
三つの二次関係も式\eqref{eq:endpoint-B-relations-v37}の $\gamma=0$ への特殊化である。

miniversal族には、二変数の混合項を用いた方が便利である。
$Q\equiv-y^2\pmod{xB}$ なので
\[
 G=x^{-1}(Q^j y-\epsilon y^n)\in(B_{0,s})_m
\]
はPBW左割算によって定義された有限の多項式である。
明示的には
\[
 G=\sum_{\ell=0}^{j-1}(-1)^\ell
 Q^{j-1-\ell}(y-sx)^{2\ell}(Z-sy)y.
\]
これは $Q$ と $y$ の可換性、
$Q+y^2=x(Z-sy)$、$y^{2\ell}x=x(y-sx)^{2\ell}$ から従う。
従って
\begin{equation}
\label{eq:endpoint-binary-CGG-v37}
 S_{0,s,\epsilon y^n}\xrightarrow{\ \sim\ }\mathcal S_{j,0}(-1,s),
 \qquad (x,y,Z,w)\longmapsto(x,y,z+sy,w+G)
\end{equation}
を得る。式の右辺の $G$ では $Z=z+sy$ と読む。
この同型は $xw+\epsilon y^n$ を $xw+q^j y$ に送る。
一般に $H'=H+xG$ としたとき、
$S_H\to S_{H'}$、$w\mapsto w+G$ が同型となることからも従う。

\subsection{明示族のminiversality}

次の二変数斉次式を取る。積は常に順序 $y$、$Z$ のPBW積である。
\begin{equation}
\label{eq:endpoint-miniversal-H-v37}
 H_{\mathbf h}
 =\epsilon y^n+\sum_{k=1}^{n}h_k y^{n-k}Z^k.
\end{equation}
$s=1$ に固定し、
\[
 U=\operatorname{Spec}\mathbb C[\gamma,h_1,\ldots,h_n,(1-\gamma)^{-1}],
 \qquad o=(0,\ldots,0)
\]
上の代数族
$\mathcal S_U=S_{\gamma,1,H_{\mathbf h}}$ を考える。
標準列の有限代数族を $\mathcal A_U$ と書く。

\begin{theorem}[全端点変形のAS正則実現]
\label{thm:endpoint-miniversal-quantization-v37}
$m=2j\ge4$ とする。
$\mathcal A_U$ の $o$ における完備化は、
端点の有限代数 $A_{j,0}(-1,1)$ のminiversal変形である。
すなわち、自然なKodaira--Spencer写像は同型であり、
\[
 \widehat{\mathcal O}_{U,o}
 =\mathbb C[[\gamma,h_1,\ldots,h_{m+1}]]
\]
はそのhullを与える。

従って、任意の局所Artin基底上の端点の標準列代数の変形は、
この実際のAS正則座標環族からの引戻しとして得られる。
同じ結論は対応する射影導来圏の形式変形に適用される。
全ての非零 $s$ および任意の $\alpha$ に対しても、
特殊ファイバーの明示的な有限代数同型を用いれば同じminiversal族を得る。
\end{theorem}

\begin{proof}
まず $s$ を変数に戻す。$s=0$ のファイバーは全ての
$\gamma,\mathbf h$ に対して可換な重み付き多項式環である。
一次交換子を $\{f,g\}=\left.s^{-1}[f,g]\right|_{s=0}$ とすると、
三変数部分では
\begin{equation}
\label{eq:endpoint-semiclassical-v37}
 \{x,y\}=x(x+2\gamma Z),\qquad
 \{y,Z\}=xZ,\qquad
 \{Z,x\}=-2xy.
\end{equation}
四変数部分のPoisson括弧は
\[
 q_\gamma=xZ-y^2+\gamma Z^2,\qquad F=xw+H_{\mathbf h}
\]
をCasimirとするJacobian括弧である。
その規格化を
$\iota_\Pi(dx\wedge dy\wedge dZ\wedge dw)=dq_\gamma\wedge dF$
とすれば、式\eqref{eq:endpoint-semiclassical-v37}と一致する。
重み付きEuler場との縮約で得られる射影一形式は
\begin{equation}
\label{eq:endpoint-rational-form-v37}
 2q_\gamma\,dF-nF\,dq_\gamma.
\end{equation}
従って、これは定理\ref{thm:endpoint-hull-v36}のHamilton型混合核を
量子化した族である。
その証明で用いた一形式 $\alpha_H$ とは全体の非零定数倍だけ異なる。
この差はPoisson錐のスカラー作用で除かれ、全局所環と軌道商を変えない。

端点のPoisson近傍において、$\gamma$ は基底切片の自由変数である。
混合方向のうち、$w\mapsto w+g$ と $w$ の倍率変更が与える軌道方向は
\[
 xS_m+\mathbb C H_*,\qquad H_*=q^j y,
\]
であり、定理\ref{thm:endpoint-hull-v36}の軌道次元計算から、
これら以外の混合軌道方向はない。
$H_*\equiv\epsilon y^n\pmod{x}$ なので
\[
 S_n/(xS_m+\mathbb C H_*)
 \simeq \bigoplus_{k=1}^{n}\mathbb C\, y^{n-k}Z^k.
\]
式\eqref{eq:endpoint-rational-form-v37}での $H_*$ と
$\epsilon y^n$ の違いは、式\eqref{eq:endpoint-binary-CGG-v37}の
$s=0$ への特殊化である座標変換によって除かれる。
従って $\gamma,h_1,\ldots,h_n$ のPoisson一次類は
座標変換軌道による商の基底を成す。

この一次独立性を固定非零量子化へ移す。
$O=\mathbb C[[s]]$、$K=\mathbb C((s))$ とし、
$\gamma=\mathbf h=0$ の有限代数弧を $\mathcal A_s$ とする。
各パラメータ微分による積の二コサイクルは $s$ で割り切れる。
それを $s$ で割った二コサイクルのHKR先頭項は、直前のPoisson一次類である。
可換特殊ファイバーのHochschild複体の可縮部分を $O$ 上で消去すると、
残る微分は $s d_1+s^2d_2+\cdots$、$d_1=[\pi,-]$ となる。
微分を $s$ で割り、さらに $d_1$ の可縮部分を消去すれば、
残る複体は
\[
 E_2^\bullet[[s]],\qquad
 E_2^\bullet=H^\bullet(HH^\bullet(\Lambda_m),[\pi,-]),
\]
を台に持ち、微分は $s$ で割り切れる。
次数二では $\dim E_2^2=m+2$ であり、上の $m+2$ 個のコサイクルの先頭項が
その基底を成す。

一方、定理\ref{thm:endpoint-hull-v36}により
$\dim_K HH^2(\mathcal A_s\otimes K)=m+2$ である。
従って、消去後の複体で次数二へ入る微分と次数二から出る微分は、
$K$ 上でともに階数零である。
$O$ は整域なので、両行列は $O$ 上でも零である。
この段階では境界による商に新たな $s$ 進torsionはなく、
像の飽和性を仮定する必要はない。
先頭項が基底を成すコサイクルは $O$ 上でも可逆な係数行列を持つため、
$K$ 上のHochschild二次類の基底となる。
これによりKodaira--Spencer写像は $K$ 上で同型である。

最後に、この結論を $s=1$ へ移すのは形式級数の評価ではなく、
次の代数的な同型である。
\begin{equation}
\label{eq:endpoint-family-scaling-v37}
\begin{gathered}
 \gamma'=s^2\gamma,\qquad h'_k=s^k h_k,\\
 S_{\gamma,s,H_{\mathbf h}}
 \xrightarrow{\ \sim\ }S_{\gamma',1,H_{\mathbf h'}},\qquad
 (x,y,Z,w)\longmapsto(x,sy,s^2Z,s^n w)
 \quad(s\ne0).
\end{gathered}
\end{equation}
実際、$Q$ は $s^2Q'$ に、$H_{\mathbf h}$ は $s^nH_{\mathbf h'}$ に移り、
$\sigma$ と $\delta_H$ の式も対応する。
パラメータの一次変換は $K$ 上で可逆なので、$s=1$ の
Kodaira--Spencer写像が体拡大後に同型となる。
有限次元線形写像の忠実な体拡大により、これは $\mathbb C$ 上でも同型である。

定理\ref{thm:endpoint-hull-v36}のhullは滑らかで次元 $m+2$ である。
そのminiversal族への分類写像は接空間上の同型を与えるので、
形式逆関数定理により、$\mathcal A_U$ の完備化自体がminiversalとなる。
$s\ne0$ は式\eqref{eq:endpoint-family-scaling-v37}で処理できる。
任意の $\alpha$ は式\eqref{eq:finite-twist-iso-v35}の有限代数同型で処理できる。
\end{proof}

\begin{remark}[得られた実現の範囲]
本定理は、滑らかな $m+2$ 次元の全categorical hullを、
通常CYのAS正則座標環の有限パラメータ族で実現する。
ただし、固定した $\alpha\ne-1$ の無限次元次数付き座標環そのものの
全変形関手を記述したものではない。
その点との同一視には特殊ファイバーの有限代数のZhang同型を用いており、
CGGの自己同型が一般の混合項に延長すると仮定しない。
$HH^1=1$ なので、パラメータを自己同型込みのfineなモジュライと呼ぶこともできない。
一方、族自体は $1-\gamma\ne0$ 上に代数的に定義されているため、
各パラメータを複素数に取った実際のAS正則代数を持つ。
その原点完備化が全形式変形を記述する、という二つの主張を区別している。
\end{remark}

\section{端点変形族の関数斜体とPI次数の分類}
\label{sec:endpoint-birational-v38}

端点の全変形を実現する座標環族には、二つの異なる型のパラメータがある。
混合項 $H$ は射影座標環の変形に寄与するが、関数斜体からは消える。
これに対し、二次部分のパラメータ $\gamma$ は、ほとんどの場合に関数斜体を区別する。
一の冪根に対応する点では別の現象が起こり、関数斜体の同型類は
その冪根の位数だけで決まる。
以下ではこれらを、表示された中心関数を固定しない
$\mathbb C$ 代数としての同型問題として扱う。

\subsection{混合項の消去と次数零の商斜体}

$m=2j\ge4$、$s\in\mathbb C^\times$、$1-\gamma s^2\ne0$ とする。
定理\ref{thm:endpoint-Ore-family-v37}の座標環を $S_{\gamma,s,H}$ と書き、
その次数零関数斜体を $\mathscr D_{\gamma,s,H}$ と書く。
ここで関数斜体とは、非零斉次元を反転した次数付き商環の零次部分である。

\begin{proposition}[中心の有理パラメータと関数斜体]
\label{prop:endpoint-division-v38}
中心有理元
\[
 \nu=\frac{J}{Q^j},\qquad
 \chi=\frac{J^2}{Q^{m+1}}
\]
の次数はそれぞれ $1,0$ である。$K=\mathbb C(\chi)$ とすると、
\[
 \mathscr D_{\gamma,s,H}\simeq
 \begin{cases}
  \operatorname{Frac}\bigl(K[\eta,\partial]/(\partial\eta-\eta\partial-1)\bigr),
     &\gamma=0,\\[2pt]
  K(X)(U;\tau),\quad \tau(X)=tX,
     &\gamma\ne0,
 \end{cases}
 \qquad
 t=\dfrac{1-s\sqrt\gamma}{1+s\sqrt\gamma}.
\]
後者の規約は $UX=tXU$ である。いずれの表示も $H$ に依存しない。
全商斜体は $\mathscr D_{\gamma,s,H}(\nu)$ であり、$\nu$ は中心の超越元である。
\end{proposition}

\begin{proof}
$B=B_{\gamma,s}$ の正規元 $x$ を反転した環に、独立な中心変数 $J$ を加える。
\[
 S_{\gamma,s,H}\longrightarrow B[x^{-1}][J],\qquad
 w\longmapsto x^{-1}(J-H)
\]
は準同型である。非零元 $\sum_{a=0}^e b_aw^a$ の像の最高 $J$ 次係数は
$b_ex^{-e}\ne0$ なので、これは単射である。
従って全商斜体は $\operatorname{Frac}(B)(J)$ に等しい。
ここでは $x$ が全代数 $S$ の正規元であるとは仮定していない。

$\gamma\ne0$ では $\rho^2=\gamma$、
$u=y+\rho Z+x/(2\rho)$、$v=y-\rho Z-x/(2\rho)$ と置く。
式\eqref{eq:endpoint-diagonal-Ore-v37}より
\[
 ux=txu,\qquad uv=-Q-\frac{t}{4\gamma}x^2.
\]
従って $v$ を消去でき、
$\operatorname{Frac}(B)=\mathbb C(Q)(x)(u;\tau)$ である。
$X=x/\nu$、$U=u/\nu$ と置き、
\[
 Q=\frac{\nu^2}{\chi},\qquad J=\frac{\nu^{m+1}}{\chi^j}
\]
を使うと、全商斜体は $K(X)(U;\tau)(\nu)$ に等しい。
PBW基底と $J$ の超越性から、この表示には追加の関係はない。
次数零の生成元 $\chi,X,U$ と次数一の中心元 $\nu$ からなるため、
零次斜体は主張の通りである。

$\gamma=0$ では $Q=xZ-sxy-y^2$ によって $Z$ を消去する。
$X=x/\nu$、$Y=y/\nu$ は $[Y,X]=-sX^2$ を満たす。
よって $\eta=X^{-1}$、$\partial=s^{-1}Y$ と置けば
$[\partial,\eta]=1$ となる。逆に $X,Y$ を回復できるので、第一Weyl斜体を得る。
\end{proof}

\subsection{スカラー交換子による分類}

$\mathbb C$ を中心に含む斜体 $F$ に対して
\[
 G_{\mathbb C}(F)=[F^\times,F^\times]\cap\mathbb C^\times
\]
と置く。角括弧は、乗法群の交換子から生成される部分群を意味する。
これはAlev--Dumasのスカラー交換子不変量である
\cite{AlevDumas1994Bir38,RichardSolotar2006Bir38}。
本例に必要な値を、中心の有理関数を固定するという仮定なしに確認する。

\begin{lemma}[二種類の関数斜体の不変量]
\label{lem:endpoint-AD-v38}
$K$ を $\mathbb C$ の体拡大、$t\in\mathbb C^\times$ とする。
$D_t(K)=K(X)(U;UX=tXU)$ と第一Weyl斜体 $D_{
\mathrm W}(K)$ に対して
\[
 G_{\mathbb C}(D_t(K))=\langle t\rangle,\qquad
 G_{\mathbb C}(D_{\mathrm W}(K))=\{1\}
\]
が成立する。第一式は $t$ が一の冪根の場合も含む。
\end{lemma}

\begin{proof}
量子斜体を反復Laurent級数斜体
$K((X))((U;\tau))$ に埋め込む。
各非零元の先頭項は一意に $aX^pU^q$、$a\in K^\times$ と書け、
先頭項を取る写像は乗法を保つ。
二つの先頭単項式の交換子は
\[
 (aX^pU^q)(bX^rU^h)(aX^pU^q)^{-1}(bX^rU^h)^{-1}
 =t^{qr-ph}.
\]
従って交換子の有限積が複素数スカラーであれば、それは $t$ の整数冪である。
逆に $UXU^{-1}X^{-1}=t$ なので、第一式を得る。

Weyl斜体には、微分作用素環
$K(\eta)[\partial;d/d\eta]$ の位数と先頭係数を使う。
非零作用素 $P$ の先頭記号を $\operatorname{lc}(P)\xi^{\operatorname{ord}(P)}$
とすると、先頭記号は乗法的である。
Ore分数へ延長した写像は
\[
 D_{\mathrm W}(K)^\times\longrightarrow K(\eta)(\xi)^\times
\]
という群準同型で、値域は可換群である。
従って交換子部分群はその核に入る一方、複素数スカラー $a$ の像は $a$ 自身である。
これより第二式が従う。
\end{proof}

一の冪根の場合には、中心の超越元が一つあることを使った同型が必要である。

\begin{lemma}[中心変数を混ぜる整数変換]
\label{lem:endpoint-torsion-change-v38}
$t$ を位数 $\ell$ の一の冪根とし、$\gcd(k,\ell)=1$ とする。
任意の体 $k_0$ が $t$ を含むとき、
\[
 D_t(k_0(\chi))\simeq D_{t^k}(k_0(\chi))
\]
が $k_0$ 斜体として成立する。
この同型は、表示された中心関数 $\chi$ を固定するとは限らない。
\end{lemma}

\begin{proof}
整数 $a,b$ を $kb-\ell a=1$ となるように選び、
\[
 X'=X^k\chi^a,\qquad U'=U,\qquad
 \chi'=X^\ell\chi^b
\]
と置く。$\chi'$ は中心にあり、$U'X'=t^kX'U'$ である。
指数変換の行列と逆行列は
\[
 \begin{pmatrix}k&0&a\\0&1&0\\\ell&0&b\end{pmatrix},\qquad
 \begin{pmatrix}b&0&-a\\0&1&0\\-\ell&0&k\end{pmatrix}
\]
である。従って
\[
 X=(X')^b(\chi')^{-a},\qquad U=U',\qquad
 \chi=(X')^{-\ell}(\chi')^k.
\]
$X$ と $\chi$ は可換なので、これらの逆変換には追加の量子係数が生じない。
整数指数が可逆に変換されることと量子Laurent単項式の一次独立性から、
$X',U',\chi'$ は表示した量子トーラス斜体を生成し、追加の関係を持たない。
\end{proof}

\begin{theorem}[端点族の双有理同型類の完全分類]
\label{thm:endpoint-birational-classification-v38}
$m=2j\ge4$ と $s\ne0$ を固定する。
任意の許される斉次混合項 $H,H'$ と、
$1-\gamma s^2\ne0$、$1-\gamma's^2\ne0$ に対し、
$\mathscr D_{\gamma,s,H}$ と $\mathscr D_{\gamma',s,H'}$ が
$\mathbb C$ 斜体として同型であるための必要十分条件は、次のいずれかである。
\begin{enumerate}
 \item $\gamma=\gamma'=0$。
 \item $\gamma,\gamma'\ne0$ に対応する $t,t'$ がともに一の冪根でなく、
       $\gamma=\gamma'$。
 \item $\gamma,\gamma'\ne0$ に対応する $t,t'$ がともに一の冪根であり、
       $\operatorname{ord}(t)=\operatorname{ord}(t')$。
\end{enumerate}
平方根の選択は、この判定に影響しない。
\end{theorem}

\begin{proof}
$\gamma=0$ では全て同じWeyl斜体になる。
$\gamma\ne0$ では $s\sqrt\gamma\ne0$ なので $t\ne1$ であり、
補題\ref{lem:endpoint-AD-v38}によりWeyl斜体との同型は存在しない。

両方が量子型の場合、同型から $\langle t\rangle=\langle t'\rangle$ が従う。
一方が有限位数、他方が無限位数ならこれは不可能である。
両方が無限位数なら、同じ無限巡回群の生成元は逆元を除いて一意なので
$t'=t$ または $t^{-1}$ である。逆に $X\mapsto X^{-1}$ によって
$t$ と $t^{-1}$ の斜体は同型になる。
\[
 s^2\gamma=\left(\frac{1-t}{1+t}\right)^2
\]
であり、同じ右辺を持つ $t'$ は $t,t^{-1}$ だけなので、
これは $\gamma=\gamma'$ と同値である。

両方が有限位数の場合、スカラー交換子群の等号は位数の等号と同値である。
同じ位数の二つの原始冪根は $t'=t^k$、$\gcd(k,\ell)=1$ と書けるので、
補題\ref{lem:endpoint-torsion-change-v38}が十分性を与える。
$\sqrt\gamma$ の符号変更は $t$ を $t^{-1}$ に変えるだけである。
\end{proof}

例えば $s=1$、位数 $5$ の場合には
$\gamma=-5+2\sqrt5$ と $\gamma'=-5-2\sqrt5$ が異なるにもかかわらず、
関数斜体は同型である。$t'=t^2$ と選べば
\[
 X'=X^2\chi,\qquad U'=U,\qquad \chi'=X^5\chi^3
\]
が明示同型を与える。
この有限位数での同一視には、中心関数を動かすことが本質的に使われている。
一方、非冪根の領域では $\gamma$ の異なるファイバーは全て双有理非同型である。

\subsection{PIファイバーと非PI端点への集積}

\begin{proposition}[PI条件と次数]
\label{prop:endpoint-PI-v38}
$S_{\gamma,s,H}$ が多項式恒等式を満たすことと、
$\gamma\ne0$ かつ $t$ が一の冪根であることは同値である。
位数を $\ell$ とすると
\[
 \operatorname{PIdeg}S_{\gamma,s,H}
 =\operatorname{PIdeg}\mathscr D_{\gamma,s,H}=\ell,
 \qquad
 Z(\mathscr D_{\gamma,s,H})=\mathbb C(\chi,X^\ell,U^\ell).
\]
この族で現れる有限位数は $\ell\ge3$ である。
\end{proposition}

\begin{proof}
$t$ の位数が $\ell$ なら、$X^\ell,U^\ell$ は中心にあり、
\[
 \{X^aU^b:0\le a,b<\ell\}
\]
は $\mathbb C(\chi,X^\ell,U^\ell)$ 上の基底である。
一次独立性は量子Laurent単項式の指数を $\ell$ で割った剰余から従う。
この基底の張る空間は有限次元の整域で、積で閉じているため斜体となる。
よって分数を全て含む。
$X,U$ の両方との可換性を基底係数ごとに比較すると、中心に残るのは
剰余 $(a,b)=(0,0)$ だけである。従って中心上の次元は $\ell^2$ である。

$t$ が無限位数なら、$K(X)((U;\tau))$ における級数展開で
$X$ と可換な元は $U$ の零次項だけを持つ。
さらに $U$ と可換である条件は $f(tX)=f(X)$ であり、
$f\in K(X)$ に対してこれは $f\in K$ を意味する。
よって中心は $K$ で、中心上の次元は無限である。
Weyl斜体の中心も特性零では $K$ であり、中心上の次元は無限である。

全商斜体は、これらに中心の超越元 $\nu$ を加えたものである。
従って有限位数の場合の中心上次元はやはり $\ell^2$ であり、
他の場合は無限のままである。
素PI環に対するPosnerの定理
\cite[Theorem~4.1]{Bresar2010Bir38}を適用すると、
座標環のPI条件とPI次数について主張を得る。
逆向きには、有限次元中心斜体の部分代数がPIであることを使う。
最後に $t=1$ は $s\sqrt\gamma=0$ を要求し、
$t=-1$ は $1=-1$ を要求するため、この量子型には現れない。
\end{proof}

\begin{corollary}[PI次数が非有界に増える代数的変形族]
\label{cor:endpoint-PI-accumulation-v38}
$s=1$ とし、miniversal族の有限型の底空間を
\[
 \mathcal U=\operatorname{Spec}
 \mathbb C[\gamma,h_1,\ldots,h_{m+1},(1-\gamma)^{-1}]
\]
とする。そのPIファイバーの集合は
\[
 \gamma=-\tan^2\left(\frac{\pi k}{\ell}\right),
 \qquad \ell\ge3,\quad 1\le k<\ell,\quad \gcd(k,\ell)=1
\]
を満たす超平面の可算和である。各点のPI次数は $\ell$ である。
この集合はZariski稠密であるが構成可能ではない。
非PIファイバーの集合もZariski稠密である。
さらに
\[
 \gamma_\ell=-\tan^2(\pi/\ell)\longrightarrow0,
 \qquad \operatorname{PIdeg}S_{\gamma_\ell,1,H}=\ell
\]
なので、非PIなCGG端点の任意の複素解析的近傍に、
PI次数が任意に大きいファイバーが存在する。
\end{corollary}

\begin{proof}
$t=e^{2\pi i k/\ell}$ を逆Cayley変換に代入すれば表示した式を得る。
$t\ne-1$ なので分母は零にならない。
異なる位数は定理\ref{thm:endpoint-birational-classification-v38}によって区別される。
特に $\gamma_\ell$ は無限個の異なる値を取り、零に収束する。
$\gamma$ 直線上の無限集合はZariski稠密であり、
$h_1,\ldots,h_{m+1}$ は任意なので、PI集合は $\mathcal U$ で稠密である。
補集合も $\gamma$ 直線上に無限個の値を持つため、同様に稠密である。
可算無限で余集合も無限な集合は複素直線の構成可能集合にはならない。
$h_1=\cdots=h_{m+1}=0$ という切断に制限すれば、
$\mathcal U$ のPI集合も構成可能でないと分かる。
最後に $H=(-1)^j y^{m+1}$ を固定すれば、この収束列の極限は
命題\ref{prop:endpoint-division-v38}のWeyl型CGG端点である。
\end{proof}

最小の位数について、$z=s^2\gamma$ の可能な値は次の多項式で表される。
\[
\begin{array}{c|c|c}
 \ell&z\text{ の方程式}&\operatorname{PIdeg}\\ \hline
 3&z+3=0&3\\
 4&z+1=0&4\\
 5&z^2+10z+5=0&5\\
 6&3z+1=0&6
\end{array}
\]
これらは、原始 $\ell$ 乗根の円分多項式と
$z(1+t)^2-(1-t)^2$ の終結式から、有理数上で求められる。
付属コードはこの計算、整数変換とその逆、交換子の先頭指数を厳密に検証する。
全ての位数に対する分類は、円分多項式の有限計算からの外挿ではなく、
補題\ref{lem:endpoint-torsion-change-v38}のB\'ezout恒等式による。

\begin{remark}[圏の同型問題との区別]
この節で分類したのは次数零関数斜体であり、座標環の同型類や
有界導来圏の同値類ではない。
特に $H$ の方向が関数斜体から消えることは、端点の $m+2$ 次元hullにおける
その方向が自明であることを意味しない。
また、関数斜体の非同型だけから、任意の導来同値の不存在は結論しない。
今回の結果は、同じ明示的なAS正則CY変形族の内部で、
Weyl型、非PI量子トーラス型、有限PI次数型を完全に区別するものである。
\end{remark}

\section{引戻し量子化の非完全交叉変形環}
\label{sec:pullback-fixed-hulls-v37}

本節では、二つの独立な障害変数を持つ固定量子圏を構成する。
接線共鳴族の一変数切断環と異なり、一般の引戻し量子化には、
二次元の被約変形空間と、Cohen--Macaulay型二の非被約厚みが同時に残る。
入力は、公表稿IX\cite{VKPoisson34}の共点三直線における全局所環と、
その引戻し族の微分Ore量子化である。

\subsection{対象と主定理}

$m\ge3$、$N_m=\binom{m+2}{2}$ とする。$x,y,z$ の重みを一、$w$ の重みを $m$ とし、
\[
 c=xy(x+y),\qquad L=Ax+By,\qquad f=u x^{m+1}
\]
を取る。$\mathbf C[x,y,w]$ 上の重み一次の導分を
\begin{align*}
 \delta(x)&=\frac{c_y}{3}+\frac{m}{3(m+3)}xL,\\
 \delta(y)&=-\frac{c_x}{3}+\frac{m}{3(m+3)}yL,\\
 \delta(w)&=f-\frac{m}{m+3}Lw
\end{align*}
と定める。$\operatorname{div}\delta=0$ であり、
\[
 S_m(A,B,u;s)=\mathbf C[x,y,w][z;s\delta]
\]
はHilbert級数 $((1-t)^3(1-t^m))^{-1}$ のAS正則4-Calabi--Yau代数である。
この構成では $x,y,w$ は互いに可換で、$[z,v]=s\delta(v)$ である。
標準全列の有限代数を
\[
 \mathcal A_m(A,B,u;s)
 =\operatorname{End}\left(\bigoplus_{i=0}^{m+2}\mathcal O_{S_m}(i)\right)
\]
と書く。

\begin{theorem}[二変数の非完全交叉厚みを持つ固定量子圏]
\label{thm:pullback-fixed-hull-v37}
各 $m\ge3$ に対し、$(2,5)$ を含む非空Zariski開集合
$U_m\subset\mathbf A^2_{A,B}$ が存在する。
$(A,B)\in U_m$、$u,s\ne0$ とすると、$\mathcal A=\mathcal A_m(A,B,u;s)$ の
miniversal環は非標準的に
\[
 \boxed{R_{\mathcal A}\simeq
 \frac{\mathbf C[[t_1,t_2,a,b]]}{(a^{m+2},ab,b^2)}}
\]
となる。また
\[
 \dim HH^\bullet(\mathcal A)=
 \begin{cases}
 (1,2,4,9)&m=3,\\
 (1,2,4,6)&m\ge4,
 \end{cases}
 \qquad HH^{>3}(\mathcal A)=0.
\]
特に、固定された $A=2,B=5,u=s=1$ は、全ての $m\ge3$ に対して結論を満たす。
\end{theorem}

一つの明示的な十分条件として、次の開集合を $U_m$ に取れる。
\[
 \kappa_1=1+\frac m3(1-B),\quad
 \kappa_2=1+\frac m3(1+A),\quad
 \kappa_3=1+\frac m3(1-A+B),
\]
\begin{align*}
 \Delta_m(A,B)={}&(B-2A)(A-2B)(A+B)
                  \kappa_1\kappa_2\kappa_3\,[m(B+2)+6]\\
 &\hspace{12mm}\cdot\prod_{k=0}^{m}[m(B+2)-3k],\\
 U_m={}&\{\Delta_m\ne0\}\setminus
       \bigcup_{\substack{\nu\in\mathbf Z_{>0}^3\\|\nu|=m+3}}
       \{\kappa=\nu\}.
\end{align*}
これらは共点基底切片、全ての次数 $m$ の移動の単射性、
残す四行のSchur切片、および三角形側の非共鳴を保証する。
例えば最高重量で消す二行の小行列式は
\[
 -\frac{[m(1-B)+3][m(B+2)+6]}9
\]
である。他の重量では標的への全射性が $\kappa_i\ne0$ から従うため、
各点の近傍で可逆小行列を選べる。
$(2,5)$ では $\kappa_1<0$、$\kappa_2,\kappa_3>0$、
$m(B+2)-3k=7m-3k>0$ であり、全ての条件が成立する。
この $U_m$ は十分な範囲を明示するもので、全例外パラメータの分類ではない。

\subsection{非零混合点の全Poisson局所環}

\begin{lemma}[残余混合係数を反転した局所環]
\label{lem:pullback-nonzero-local-v37}
上の開集合を十分小さく取る。$u_0\ne0$ とし、
\[
 \pi=(c,D,[F]),\qquad
 D=(Ax+By)\partial_z,\qquad F=u_0x^{m+1}\partial_z
\]
を重み付きPoissonスキーム $\mathscr P_m$ の点とする。このとき
\[
 \widehat{\mathcal O}_{\mathscr P_m,\pi}
 \simeq
 \frac{\mathbf C[[z_1,\ldots,z_{N_m+9},a,b]]}
 {(a^{m+2},ab,b^2)}.
\]
\end{lemma}

\begin{proof}
公表稿IXの共点切片では、六つの $\mathrm{SL}_3$ 軌道変数を除いた基底は
\[
 \mathbf C[[A-A_0,B-B_0,a,b]]/(ab,b^2)
\]
であり、混合写像は $M_0+aM_1+bM_2$ となる。
同稿の全重みの核計算により、次数 $m$ の $w$ 移動は
混合核の階数 $N_m$ の自由直和因子をなし、残る一方向は
$x^{m+1}\partial_z$ で生成される。
Schur消去後の残余方程式は、単元倍を除いて正確に
\[
 a^{m+2}u=0
\]
である。$b$ の一次項は重量の差によって消え、他の $b$ 項は $ab=b^2=0$ により消える。

この消去を、$F=0$ の完備局所環へ非零の複素数 $u$ を代入する操作と解釈してはならない。
実際、基底 $(c,D)$ の座標変換は $F$ に依存せず、
混合写像のSchur消去と源の基底変換は全て $F$ の係数について線形である。
従って、基底方向の完備化と、$F$ の係数に関する多項式環との上で
同じ消去が成立する。任意の有限の $F$ の値で改めて完備化してよい。
$u=u_0+v$ と置けば $u_0+v$ は単元なので、残余イデアルは $(a^{m+2})$ となる。
六つの軌道変数、$A,B$、$N_m$ 個の移動変数、$v$ を合わせると、
滑らかな変数は $N_m+9$ 個である。
\end{proof}

\begin{lemma}[Poisson安定化群の二方向]
\label{lem:pullback-poisson-stabilizer-v37}
補題\ref{lem:pullback-nonzero-local-v37}の点における
重み付き座標変換群の軌道次元は $N_m+7$ であり、
安定化Lie代数は
\[
 \mathbf C\,x\partial_z\oplus\mathbf C\,y\partial_z
\]
である。
\end{lemma}

\begin{proof}
座標変換群のLie代数は、Euler場を除いて
$\mathfrak{sl}_3$、$w\partial_w$、$g_m\partial_w$ からなり、次元は $N_m+9$ である。
三つの相異なる共点直線の積 $c=xy(x+y)$ の $\mathrm{SL}_3$ 軌道は六次元である。
その安定化Lie代数は $x\partial_z,y\partial_z$ の二方向であり、
両者は $D=L\partial_z$ も保存する。

公表稿IXの核計算では、全ての次数 $m$ の $w$ 移動の像が $N_m$ 次元であり、
$F_{\mathrm{ex}}=x^{m+1}\partial_z$ の類はその像に入らない。
特に二変数移動では、$V=(H_c+LE_{xy})/3$ として
\[
 V(x^{m-k}y^k)=
 \left(\frac{m(A+1)}3-k\right)x^{m+1-k}y^k+
 \left(\frac{m(B+2)}3-k\right)x^{m-k}y^{k+1}.
\]
$(A,B)=(2,5)$ では後者の係数は $7m/3-k>0$ なので単射であり、
余核は $x^{m+1}$ で生成される。非二変数移動の単射性は重量ごとの核の表示から従う。
$u_0\ne0$ なので $w$ の倍率がさらに独立な一方向を与える。
従って全軌道の次元は少なくとも $6+N_m+1=N_m+7$ である。

一方、発散零代表 $\widehat\pi=\partial_z\wedge\delta$ は
$x\partial_z,y\partial_z$ によって保存される。
例えば $[x\partial_z,\delta]=-\delta(x)\partial_z$ であり、
これを $\partial_z$ と外積すれば零となる。
従って安定化群の次元は少なくとも二で、軌道次元の上界も $N_m+7$ となる。
上下界が一致する。
\end{proof}

\subsection{Artin厚みを保つ形式移送と軌道切片}

次の二補題は、一変数のArtin環に限らず使える。

\begin{lemma}[定数Artin厚みを持つ錐の弧]
\label{lem:artin-cone-arc-v37}
$k$ を標数零の体、$T$ を剰余体 $k$ の有限局所 $k$ 代数とする。
斉次アフィン錐 $P$ の点 $\pi$ において
\[
 \widehat{\mathcal O}_{P,\pi}\simeq T[[z_1,\ldots,z_d]]
\]
とする。$p(h)=h\pi+O(h^2)$ が $P(k[[h]])$ の弧なら、
$K=k((h))$ 上で
\[
 \widehat{\mathcal O}_{P_K,p_K}
 \simeq(T\otimes_kK)[[z_1,\ldots,z_d]].
\]
\end{lemma}

\begin{proof}
錐のスカラー作用で $p(h)$ を $p(h)/h=\pi+O(h)$ に移す。
形式弧に沿う移送補題\ref{lem:arc-translation-v35}を適用する。
$k[[h]]$ は被約なので、$T$ の極大イデアルは全て零へ写る。
従って弧は滑らかな $z$ 方向だけを動かし、その値を形式的に平行移動すればよい。
\end{proof}

\begin{lemma}[滑らかな軌道方向だけを除く切片]
\label{lem:artin-orbit-slice-v37}
有限局所 $k$ 代数 $T$ の剰余体を $k$ とし、
$R=T[[z_1,\ldots,z_d]]$ とし、$\operatorname{char}k=0$ を仮定する。
滑らかな代数群の作用に対し、閉点の軌道の次元を $q$ とする。
このとき、軌道に横断するminiversal形式切片の環は
\[
 T[[z_{q+1},\ldots,z_d]]
\]
と取れる。
\end{lemma}

\begin{proof}
群のLie代数から、軌道接写像が単射となる $q$ 方向 $v_1,\ldots,v_q$ を選ぶ。
標数零なので形式指数 $\exp(\sum t_iv_i)$ を使え、その逆は $\exp(-\sum t_iv_i)$ である。
軌道は被約スキームの像であるから、$T$ の冪零元は軌道上で零となる。
従って軌道の接空間は $z$ 方向に含まれる。
$z$ の線形座標変換により、最初の $q$ 個の $z$ の軌道接写像を可逆にできる。
切片を $z_1=\cdots=z_q=0$ と定める。

任意の形式的な点について、群の上の $q$ 座標を使って、作用後の最初の $q$ 個の
$z$ 座標を零にする方程式を解く。そのJacobianは閉点で可逆なので、
形式逆関数定理により一意な解を得る。
従って全形式近傍は、この切片と $q$ 個の軌道変数との積になる。
切片の接空間は軌道接空間を割った空間であり、miniversal性を持つ。
\end{proof}

体拡大から元の複素係数へ同型型を戻す際には、次の初等的な補題を使える。
これは有限次元だけの環に限らない。

\begin{lemma}[複素係数の形式同型の下降]
\label{lem:complex-formal-descent-v37}
$R,R'$ を剰余体 $\mathbf C$、有限埋込次元の完備Noether局所
$\mathbf C$ 代数とする。ある拡大体 $K/\mathbf C$ に対して、$K$ 代数として
\[
 R\widehat\otimes_{\mathbf C}K
 \simeq R'\widehat\otimes_{\mathbf C}K
\]
なら、$R\simeq R'$ である。同型の選択は非標準的である。
\end{lemma}

\begin{proof}
$R=\mathbf C[[X_1,\ldots,X_e]]/(f_1,\ldots,f_r)$、
$R'=\mathbf C[[Y_1,\ldots,Y_{e'}]]/(g_1,\ldots,g_{r'})$ と表示する。
全ての $f_i,g_j$ の係数は可算個なので、これらを含む可算部分体
$k_0\subset\mathbf C$ を取れる。
与えられた同型と逆同型について、各変数の像を形式冪級数に持ち上げる。
定義関係が定義イデアルへ入ることと、二つの合成が恒等写像になることを表す
有限個のイデアル所属の係数冪級数も選ぶ。
これら全ての係数を $k_0$ へ添加して得る体 $K_0\subset K$ は可算である。

$\mathbf C$ の $k_0$ 上の超越次数は非可算なので、$K_0/k_0$ の
超越基底を $\mathbf C$ の代数的独立な集合へ送り、代数的部分の埋込みを延長することで、
$k_0$ 上の体の埋込み $K_0\hookrightarrow\mathbf C$ が存在する。
これを、同型・逆同型および所属を表す全係数に適用する。
各係数恒等式は有限個の体演算からなるので保たれ、$f_i,g_j$ の係数は固定される。
変数の像を $\mathbf C$ 線形に延長すれば、元の二つの完備局所環の間の
互いに逆な同型を得る。
\end{proof}

\subsection{量子側の $HH^1$ と主定理の証明}

Ore族の積はPBW基底上で
\[
 a\star_s b=\sum_{r\ge0}\frac{s^r}{r!}
                  (\partial_z^r a)(\delta^r b)
\]
と書ける。各対について有限和なので、$s=h$ は有限代数の多項式的な量子化弧を与える。
$h\ne0$ 上では、$x,y,w$ を固定して $z\mapsto hz$ とすることで
$S_m(A,B,u;h)\simeq S_m(A,B,u;1)$ を得る。
また $w\mapsto uw$ により、非零 $u$ も一へ規格化できる。

\begin{lemma}[量子側の二つの外的導分]
\label{lem:pullback-quantum-HH1-v37}
主定理の仮定の下で $\dim HH^1(\mathcal A)=2$ である。
生成元は、座標環の自己同型
\[
 z\longmapsto z+tx,\qquad z\longmapsto z+ty
\]
を $t=0$ で微分した導分から得られる。
\end{lemma}

\begin{proof}
$x,y,w$ は可換なので、上の二つの置換はOre関係を保ち、逆も同じ形で与えられる。
次数も保つため、標準有限代数の全頂点を固定する自己同型を誘導する。
対応する導分を $D_x,D_y$ と書けば、生成元上で
\[
 D_x(z)=x,\quad D_y(z)=y,\qquad
 D_x(x)=D_x(y)=D_x(w)=D_y(x)=D_y(y)=D_y(w)=0.
\]
PBW多項式上の通常の $x\partial_z,y\partial_z$ と、そのまま同一視する必要はない。

冪等元を固定する内的導分は、対角のスカラー元から来る。
従って各次数一の隣接射空間にはスカラー倍で作用する。
上の二つの導分の非零線形結合は、$x,y$ を零へ送り $z$ を非零の一次式へ送るため、
内的ではない。この議論は $\mathbf C((h))$ 上でも成立する。
よって $HH^1$ の次元は少なくとも二である。

逆向きには、最初のHochschild差分補題\ref{lem:first-differential-v35}を使う。
可換基点では $HH^1=H^0(T_{X_m})$、一次差分は $[\pi,-]$ である。
補題\ref{lem:pullback-poisson-stabilizer-v37}からその核は二次元であり、
$HH^0=\mathbf C$ からの差分は零なので、generic fiberで $HH^1\le2$ となる。
非零係数間の同型によって全ての $s\ne0$ へ移せば結論を得る。
\end{proof}

\begin{proof}[定理\ref{thm:pullback-fixed-hull-v37}の証明]
$K=\mathbf C((h))$ とする。可換標準列代数のhullとPoisson錐の原点完備環の同型は、
HKRを接写像として持つ。従って、代数構造スキーム $\mathfrak A$ の可換点における
全完備局所環は、Poisson錐の原点環に $b_0$ 個の軌道変数を付け加えたものである。
ここで $b_0$ は可換点でのHochschild境界 $B_E^2$ の次元である。

Ore弧のPoisson座標は $h\pi+O(h^2)$ である。
補題\ref{lem:pullback-nonzero-local-v37}と
補題\ref{lem:artin-cone-arc-v37}、形式弧の移送から、generic fiberの全局所環は
\[
 \widehat{\mathcal O}_{\mathfrak A_K,\mathcal A_h\otimes K}
 \simeq
 \frac{K[[z_1,\ldots,z_{b_0+N_m+9},a,b]]}{(a^{m+2},ab,b^2)}
\]
となる。

中心はどの標準有限代数でもスカラーのみで、$\dim\operatorname{im}d^0=m+2$ は一定である。
可換点の $HH^1$ は $N_m+9$ 次元、量子側では補題により二次元なので、
基底変換軌道の次元は
\[
 b_\eta=b_0+N_m+7
\]
となる。補題\ref{lem:artin-orbit-slice-v37}を適用すれば、二つの滑らかな変数が残り、
generic fiberのhullは
\[
 \frac{K[[t_1,t_2,a,b]]}{(a^{m+2},ab,b^2)}
\]
である。

非零係数の同型は $\mathcal A_h\otimes K\simeq\mathcal A_1\otimes K$ を与える。
有限代数の形式切片は体拡大と可換なので、補題\ref{lem:complex-formal-descent-v37}により、
この同型型は $\mathbf C$ 上へ降下する。全ての非零 $s,u$ も同じである。

得られたhullの埋込次元は四なので $HH^2$ の次元は四である。
$HH^{>3}=0$ は、可換点での消滅、有限自由Hochschild複体の特殊化、
および非零係数間の同型による。
同じ有限複体のEuler特性は $m=3$ で $-6$、$m\ge4$ で $-3$ であり、
$HH^0=1,HH^1=2,HH^2=4$ を代入すれば表示した $HH^3$ の次元を得る。
\end{proof}

\begin{corollary}[二つの滑らかな方向の明示実現]
\label{cor:pullback-reduced-parameters-v37}
主定理の点 $(A_0,B_0,u,s)$ を固定する。
$A=A_0+p,B=B_0+q$ とするOre族は、hullの被約化
\[
 (R_{\mathcal A})_{\mathrm{red}}\simeq\mathbf C[[t_1,t_2]]
\]
を形式的にパラメータ付ける。すなわち、hullからこの族への写像を一つ選べば、
誘導される写像
$(R_{\mathcal A})_{\mathrm{red}}\longrightarrow\mathbf C[[p,q]]$
は同型である。
\end{corollary}

\begin{proof}
まずPoisson側の二つの接ベクトル $\partial_A\pi,\partial_B\pi$ は、
座標変換軌道の接空間を法として独立である。
実際、この線形結合を軌道接ベクトルと仮定すると、その $c$ 成分は零なので、
$\mathfrak{sl}_3$ 成分は $x\partial_z,y\partial_z$ の線形結合である。
両者は $D$ を固定し、$w$ 移動と倍率も $D$ を変えない。
一方、上の二つのパラメータは $D=(Ax+By)\partial_z$ を独立に変えるため、
線形結合の係数はともに零となる。

量子側への独立性の移送も確認する。
$O=\mathbf C[[h]]$ 上で可換点の可縮部分を除いたHochschild複体を取り、
次数一から二への微分を
\[
 d_h=h(D_1+hD_2+\cdots),\qquad D_1=[\pi,-]
\]
と書く。$HH^1$ の次元は、$\ker D_1$ とgeneric fiberの両方で二である。
次数零の微分は中心の次元が一定一であるため零なので、
$\operatorname{rank}_K(d_h/h)=\operatorname{rank}_{\mathbf C}D_1=N_m+7$ となる。
従って $d_h/h$ の像は $O$ 上の自由な直和因子である。
実際、$D_1$ の可逆小行列を消去すれば、残余行列は $K$ 上で階数零、
従って $O$ 上でも零となる。

Ore族の二つのKodaira--Spencerコサイクルは $h$ で割り切れ、
上の縮約複体における商の先頭項は $\partial_A\pi,\partial_B\pi$ である。
これらが $K$ 上のコホモロジーで従属なら、係数を一つの $h$ 冪で規格化し、
少なくとも一つを単元とした $O$ 係数の関係を取れる。
$d_h/h$ の像が直和因子なので、この関係を $h=0$ へ特殊化すると、
$\partial_A\pi,\partial_B\pi$ の非零線形結合が $\operatorname{im}D_1$ に入る。
これは前段と矛盾する。
非零 $h$ と一との同型は $A,B$ に依存しないので、固定された非零 $s$ でも
二つのKodaira--Spencer類は独立である。

形式族はhullの性質により $R_{\mathcal A}\to\mathbf C[[p,q]]$ を与える。
標的は被約なので、この写像はhullの被約化を通る。
主定理により被約化は二変数の形式冪級数環であり、
今証明した接写像の独立性から二変数間のJacobianが可逆である。
形式逆関数定理によって結論を得る。
\end{proof}

\subsection{変形環の構造と検証}

$n=m+2$ と置く。横断Artin環
\[
 T_m=\mathbf C[a,b]/(a^n,ab,b^2)
\]
の基底は $1,a,\ldots,a^{n-1},b$、socleは
$\langle a^{n-1},b\rangle$ である。従って長さは $m+3$、socle次元は二である。
主定理の環は $\mathbf C[[t_1,t_2]]$ 上の階数 $m+3$ の自由加群であり、
二次元Cohen--Macaulay環、型二、Hilbert--Samuel重複度 $m+3$ を持つ。
特にGorensteinではなく、完全交叉でもない。
関係イデアルのHilbert--Burch行列は
\[
 \begin{pmatrix}
 b&a^{n-1}\\
 -a&0\\
 0&-b
 \end{pmatrix}
\]
と取れる。高さ二に対して最小生成元は三個である。

二つの障害方向は異なる振舞いを示す。
$b$ 方向は $b^2=0$ によって二次で止まり、$a$ 方向の最初の純粋な障害は
$a^{m+2}=0$ によって次数 $m+2$ に現れる。
両者は二次関係 $ab=0$ でも結び付く。
一方、被約な形式基底への変形は $a=b=0$ を通り、二つの滑らかな方向を残す。

付属の有理数検算は、全標準有限代数の導分を矢印像から計算し、
Ore関係のLeibniz式の全係数を照合する。
$(A,B,u,s)=(2,5,1,1)$ における結果は次の通りである。
\[
 \begin{array}{c|r|r|r|r|r}
 m&\text{未知数}&\text{方程式}&\operatorname{rank}_{\mathbf Q}
 &\dim\operatorname{Der}_E&\dim HH^1\\\hline
 3&78&180&71&7&2\\
 4&102&234&94&8&2\\
 5&129&294&120&9&2\\
 6&159&360&149&10&2
 \end{array}
\]
別係数 $(3,7,2,-1/2)$ と可換基点も照合し、
二つの明示導分が内的導分から独立なことを行列上で確かめた。
$A,B$ の関係微分を二列追加した行列では、全ての非零係数例で
階数が正確に二増え、二つのKodaira--Spencer類の独立性も確認できる。
Poisson安定化群も別の多項式Lie微分行列で検査し、
$m=3,4,5,6$ の軌道階数 $17,22,28,35$ と二次元の核を得る。
これらの有限計算は、本文の全重みの構造的証明を補助する。

本節は、固定した標準有限代数、従って対応する射影導来圏の全結合変形を扱う。
無限次元の座標環の全非次数付き変形や、
パラメータ $(A,B)$ 間の射影圏の完全な同値分類を主張するものではない。

\section{反標準Hochschild--Serre環と射影導来圏の区別}
\label{sec:anticanonical-serre-v35}

この節では、斉次座標環の中心を、射影圏に内在する不変量へ移す。
中心そのものはZhang twistで変わり得るため、中心の相異だけから
射影圏の非同値を推論することはできない。
本族では通常のCalabi--Yau代表を選び、Serre関手に対応する
Hochschild--Serre環の一つの対角を計算することで、この問題を解決する。
Hochschild--Serre環の一般的な定義とMorita不変性は
Lin--Zhang\cite[Definition~3.3 and Theorem~3.4]{LinZhangSerre2025}による。
以下の内容はその一般構成を本族に適用し、対角環を具体的に決定するものである。

\subsection{反標準対角の定義}

$j,r\ge0$、$m=2j+r\ge3$、$s\in\mathbb C^*$ とし、
\[
 \ell=m+3,\qquad d=m+r+1,\qquad
 \alpha_0=-\frac{m+r+3}{m+3}
\]
と置く。この節で $A=\mathcal S_{j,r}^{(q)}(\alpha_0,s)$ は
通常の4-Calabi--Yau代表を表し、
$\mathcal C=\operatorname{Perf}_{\mathrm{dg}}(\operatorname{qgr}A)$ と書く。
その有界導来圏の標準tilting対象とBeilinson代数を
\[
 T=\bigoplus_{i=0}^{\ell-1}\mathcal O_A(i),\qquad
 B=\operatorname{End}(T)
\]
とする。$B$ は有限次元かつ有限global dimensionを持ち、
$\mathcal C\simeq\operatorname{Perf}_{\mathrm{dg}}(B)$ である。
AS正則代数のSerre双対性と標準tilting列について
\cite[Proposition~4.4, Theorems~4.12 and~4.14]{MinamotoMori2011}を用いる。
通常CY条件により、Serre関手は
\[
 \mathbb S_{\mathcal C}\simeq(-\ell)[3]
\]
となる。

\begin{definition}
この三次元射影圏の反標準Hochschild--Serre環を
\[
 \operatorname{AC}(\mathcal C)
 :=\bigoplus_{n\ge0}
 \operatorname{Hom}_{D(\mathcal C^{\mathrm{op}}\otimes\mathcal C)}
 \left(I_{\mathcal C},\mathbb S_{\mathcal C}^{-n}[3n]\right)
\]
と定める。右辺では関手を表すdg双加群を用い、積は関手の合成から誘導する。
したがって、これはHochschild--Serre環の $(-n,3n)$ 対角である。
\end{definition}

この定義に、構造層、次数移動 $(1)$、例外列の標識は含めない。
dg Morita同値はSerre双加群を移し、合成を保つので、
$\operatorname{AC}(\mathcal C)$ は次数付き環として不変である。
三角関手間の任意の自然変換を用いる定義とは区別し、
以下では常に上のdg双加群のHomを計算する。

\subsection{局所Ore表示と一つの中心化環}

$q=xz-y^2$ とし、不変多項式 $J$ を
\[
 J=x^{r+1}w+q^j\sum_{k=0}^{r}
       \binom rk\frac{q^{r-k}y^{2k+1}}{2k+1}
\]
とする。$x,q,J$ は $N$ の核に属するので、それらの間の星積は通常の可換積である。

\begin{lemma}[局所Ore表示]
\label{lem:ore-centralizer-v35}
実際には任意の $\alpha\in\mathbb C$ に対して、$x$ は正規元であり、
\[
 A[x^{-1}]
 \simeq R[t;\delta],\qquad
 R=\mathbb C[x^{\pm1},q,J],\qquad t=x^{-1}\star y,
\]
\[
 \delta=s\left(\alpha x\partial_x
              +2(\alpha+1)q\partial_q
              +d(\alpha+1)J\partial_J\right)
\]
である。ここでは $A=\mathcal S_{j,r}^{(q)}(\alpha,s)$ と解釈する。
$u=q/x^2$ と置くと
\[
 [t,u]=2su,\qquad
 \operatorname{Cent}_{A[x^{-1}]}(u)=R.
\]
\end{lemma}

\begin{proof}
$T_0=1+sN$ と置くと $a\star x=T_0^\alpha(a)x$ なので、
$x$ は正規でありOre局所化が存在する。
Laurent多項式環 $\mathbb C[x^{\pm1},y,z,w]$ 上でも $N$ は局所冪零で、
CGG積は各入力について有限和である。
通常の可換座標として
\[
 \mathbb C[x^{\pm1},y,z,w]
 =\mathbb C[x^{\pm1},q,J,t],\qquad t=y/x,
\]
と書ける。実際、$y=xt$、$z=(q+y^2)/x$、最後に $J$ の式から $w$ を解けばよい。
この座標では $N(t)=1$、$N(R)=0$ であるから、$f\in R$ に対し
\[
 t\star f-f\star t=sX_\alpha(f)=\delta(f).
\]
$t$ の星冪は通常の $t$ の冪に関してmonicな三角変換なので、
これで上のOre表示を得る。$\delta(u)=2su$ も直接従う。

$a=\sum_{k=0}^{h}a_kt^k\in R[t;\delta]$、$a_h\ne0$ とする。
$h>0$ なら $[a,u]$ の $t^{h-1}$ の係数は
$h a_h\delta(u)=2sh a_hu$ であり、整域性と標数零から非零である。
したがって $u$ と可換な元は $h=0$、すなわち $R$ の元に限る。
\end{proof}

\subsection{双加群Homを多項式中心へ戻す}

整数 $D\ge0$ に対して、次数移動 $(D)$ を表す $B$ 双加群を
\[
 M_D=\operatorname{RHom}(T,T(D))
\]
と書く。左作用は $T(D)$ の自己準同型、右作用は $T$ の自己準同型から来る。
AS正則性から
\[
 \operatorname{Ext}^p\bigl(\mathcal O_A(i),\mathcal O_A(j+D)\bigr)
 \simeq
 \begin{cases}
 A_{j+D-i}&p=0,\\
 (A_{i-j-D-\ell})^*&p=3,\\
 0&\text{その他}
 \end{cases}
\]
である。$0\le i,j\le\ell-1$、$D\ge0$ では
$i-j-D-\ell<0$ なので、$M_D$ は次数零に集中する。

\begin{proposition}[正の移動に対する自然変換の完全計算]
\label{prop:positive-shift-center-v35}
$s\ne0$、$D\ge0$ とする。上の族の任意の $\alpha$ について、自然な写像は同型
\[
 \boxed{
 Z(A)_D\xrightarrow{\ \sim\ }
 \operatorname{Hom}_{D(B^{\mathrm{op}}\otimes B)}(B,M_D)
 }
\]
を与える。これらの同型は次数移動による積と両立する。
\end{proposition}

\begin{proof}
$B,M_D$ が次数零に集中するため、右辺は通常の双加群準同型
$\operatorname{Hom}_{B^{\mathrm{op}}\otimes B}(B,M_D)$ である。
この準同型は、各冪等元に対応する対角成分
\[
 c_i\in\operatorname{Hom}(\mathcal O_A(i),\mathcal O_A(i+D))=A_D,
 \qquad 0\le i\le\ell-1,
\]
によって定まり、条件はすべての $a\in A_{j-i}$ に対する
\[
 c_j a=a c_i\qquad(0\le i\le j\le\ell-1)
\]
である。この式以降、積は $A$ の積である。

隣接する対象間の射 $x$ から、局所化環の中で
\[
 c_i=x^ic_0x^{-i}
\]
を得る。対象 $0$ から $2$ への射 $q$ を使うと
\[
 (x^2c_0x^{-2})q=qc_0,
 \qquad [c_0,x^{-2}q]=0.
\]
補題\ref{lem:ore-centralizer-v35}より $c_0\in R$ となり、特に $c_0x=xc_0$ である。
よってすべての $c_i$ は一つの元 $c=c_0\in A_D$ に等しい。
対象 $0$ から $1$ への射 $y,z$、および対象 $0$ から $m$ への射 $w$ を使うと、
$c$ は $y,z,w$ とも可換である。$m\le\ell-1$ なので最後の射も列内に存在する。
従って $c\in Z(A)_D$ である。

逆に斉次中心元は各対象に同じ乗法を与え、全可換図式を満たす。
二つの中心元の自然変換の積はその中心元の積なので、積との両立性も従う。
\end{proof}

\begin{theorem}[反標準環の $A/D$ 型分類]
\label{thm:anticanonical-ADE-v35}
通常CY代表の射影圏について
\[
 \operatorname{AC}(\mathcal C)
 \simeq\bigoplus_{n\ge0}Z(A)_{n\ell}
\]
が次数付き環として成り立つ。右辺は $\ell$-Veronese部分環で、次数を $\ell$ で割る。
その具体形は次の通りである。
\[
\begin{array}{c|c|c}
 \text{条件}&\operatorname{AC}(\mathcal C)&\text{生成元の反標準次数}\\\hline
 r\ge2&\mathbb C&0\\[1mm]
 r=1&\mathbb C[S,P,R]/(R^2-SP^2-\frac{16}{81}S^{m+3})
       &(2,m+2,m+3)\\[1mm]
 r=0&\mathbb C[U,V,W]/(UV-W^{m+3})&(2,m+1,1)
\end{array}
\]
$r=1$ の場合は $D_{m+4}$ 型、$r=0$ の場合は $A_{m+2}$ 型のDu Val曲面である。
特に
\[
 \bigl(\dim\operatorname{AC}^1,\dim\operatorname{AC}^2\bigr)
 =\begin{cases}
 (0,0)&r\ge2,\\
 (0,1)&r=1,\\
 (1,2)&r=0.
 \end{cases}
\]
\end{theorem}

\begin{proof}
$\mathbb S_{\mathcal C}^{-n}[3n]=(n\ell)$ に
命題\ref{prop:positive-shift-center-v35}を適用すると最初の同型を得る。
$r\ge2$ では $Z(A)=\mathbb C$ である。
$r=1$ では既に求めた中心生成元の斉次次数は
$2\ell,(m+2)\ell,\ell^2$ なので、中心全体がVeroneseに入り、表示もそのまま残る。

$r=0$ では $m$ は偶数であり、$Z(A)=\mathbb C[q,J]$、
$(\deg q,\deg J)=(2,m+1)$ である。
$\ell=m+3$ は奇数なので
\[
 2a+(m+1)b\equiv0\pmod\ell
 \iff a\equiv b\pmod\ell.
\]
従ってVeroneseに属するすべての単項式は
\[
 U=q^\ell,\qquad V=J^\ell,\qquad W=qJ
\]
で生成される。指数半群の唯一の基本関係は $UV=W^\ell$ である。
例えば $UV$ を除く単項式の像の指数は相異なるので、追加の関係はない。
各生成次数を $\ell$ で割ると表示した次数を得る。
最後の低次数次元とDu Val型は各表示から従う。
\end{proof}

\begin{remark}[Hilbert級数]
上の反標準次数を使うと
\[
 H_{\operatorname{AC}}(t)=
 \begin{cases}
 1&r\ge2,\\[1mm]
 \displaystyle\frac{1-t^{2\ell}}
 {(1-t^2)(1-t^{\ell-1})(1-t^\ell)}&r=1,\\[3mm]
 \displaystyle\frac{1-t^\ell}
 {(1-t)(1-t^2)(1-t^{\ell-2})}&r=0.
 \end{cases}
\]
も得る。$r=0$ の座標環中心は滑らかな二変数多項式環だが、
射影圏から選び出す反標準次数の部分環は $A_{m+2}$ 特異点になる。
両者は異なる環である。
\end{remark}

\subsection{反標準環が定めるorbifold曲線}

正次数部分を持つ二つの場合について、環 $R_{\mathrm{ac}}=\operatorname{AC}(\mathcal C)$ から
\[
 \mathfrak C_{\mathrm{ac}}
 :=[(\operatorname{Spec}R_{\mathrm{ac}}\setminus\{0\})/\mathbb G_m]
\]
を作る。作用は反標準次数による。このstacky Projも射影導来圏に内在する。

\begin{proposition}
\label{prop:anticanonical-orbifold-v35}
$r=1$ では $\mathfrak C_{\mathrm{ac}}$ は粗空間 $\mathbb P^1$、
三つのstacky点の安定化群位数が $(2,2,m+2)$ である滑らかなorbifold曲線となる。
$r=0$ では
\[
 \mathfrak C_{\mathrm{ac}}\simeq\mathbb P_{\mathrm{stack}}(2,m+1).
\]
\end{proposition}

\begin{proof}
$r=1$ とし、$d=m+2$ は奇数、$\beta=4/9$ と置く。
環は $R^2-SP^2-\beta^2S^{d+1}=0$、重量は $(S,P,R)=(2,d,d+1)$ である。
このaffine曲面の特異点は原点だけであり、正次数の作用の安定化群は有限なので、
原点を除いた商は滑らかなDeligne--Mumford曲線である。
$S\ne0$ では $S=1$ とする切片により
\[
 [\operatorname{Spec}\mathbb C[P,R]/(R^2-P^2-\beta^2)\,/\,\mu_2]
\]
を得る。残余作用は $P\mapsto-P$、$R\mapsto R$ である。
従って粗空間は $R$ を座標とする $\mathbb A^1$ で、$R=\pm\beta$ に
安定化群 $\mu_2$ があり、他の点の安定化群は自明である。

$S=0$ では $P\ne0,R=0$ となり、商には一点だけを加える。
$P=1$ の切片で関係式の $S$ 偏微分はこの点で $-1$ なので、$R$ が局所座標となる。
残余群 $\mu_d$ は $R$ に重量 $d+1\equiv1\pmod d$ で作用する。
よってこの点の安定化群は $\mu_d$ である。
粗空間は上の $\mathbb A^1$ の滑らかな完備化なので $\mathbb P^1$ となる。

$r=0$ では $\ell=m+3$ が奇数であり、
$R_{\mathrm{ac}}=\mathbb C[q,J]^{(\ell)}$ である。
$\mu_\ell$ は $(q,J)$ に重量 $(2,\ell-2)$ で作用する。
$\gcd(\ell,2)=\gcd(\ell,\ell-2)=1$ なので、原点を除けばこの作用は自由である。
したがって $\operatorname{Spec}R_{\mathrm{ac}}\setminus\{0\}$ は
$(\mathbb A^2\setminus\{0\})/\mu_\ell$ と同型。
残余群 $\mathbb G_m/\mu_\ell$ による商は、元の重量 $(2,m+1)$ の
$\mathbb G_m$ 商と同じである。
\end{proof}

この命題は反標準環から内在的な曲線を取り出すものである。
元の非可換射影圏からこの曲線への全域的な射や、
中心上有限な代数によるfibrationを構成したという主張は含めない。

\subsection{同じ重みの異なる共鳴を導来圏で分離する}

\begin{corollary}
\label{cor:derived-separation-v35}
$s\ne0$ とする。
\begin{enumerate}
\item 固定した奇数 $m\ge3$ に対し、$r=1$ の射影圏は、
同じ重みの任意の $r\ge3$ の射影圏と $\mathbb C$ 線形導来同値ではない。
\item 固定した偶数 $m\ge4$ に対し、$r=0$ の射影圏は、
同じ重みの任意の $r\ge2$ の射影圏と $\mathbb C$ 線形導来同値ではない。
\end{enumerate}
これらの結論は、各族の任意の $\alpha$ に対して成り立つ。
\end{corollary}

\begin{proof}
CY代表については前定理の反標準環が異なる。
この不変量はdg Morita同値を排除する。
通常の三角同値の不存在も従うことを確認する。
仮に二つの有界導来圏が三角同値なら、標準tilting対象の像は
相手の有限次元Beilinson代数上のtilting complexである。
Rickardのderived Morita定理により、両Beilinson代数を結ぶ両側tilting complexが存在する
\cite{RickardMorita1989,RickardDerived1991}。
それはdg Morita同値を与えるので矛盾する。
ここでは、最初に与えられた三角関手そのものが標準関手であると仮定する必要はない。

全 $\alpha$ は同じ $j,r$ のCY代表とZhang twistで結ばれ、
対応する $\operatorname{qgr}$ 圏は同値である。
従って結論は全 $\alpha$ に移る。
\end{proof}

\begin{example}[重み三の二つの十八次元成分]
$m=3$ では $r=3$ の滑らかな接線成分の代表と、
$r=1$ の一般長さ二の接線成分の代表に対し、
反標準環はそれぞれ $\mathbb C$ と $D_7$ 型の環である。
従って、二つの量子重み付き射影三次元圏は導来同値でない。
両者の次数零量子関数斜体は同じ第一Weyl斜体であるから、
この区別はその双有理表示では検出できない。
\end{example}

\begin{remark}[この不変量の射程]
$r\ge2$ 同士では反標準環は常に $\mathbb C$ なので、この対角だけでは区別できない。
また、ここで用いたSerre関手は外側の射影圏全体のものである。
四項のVK切り出しのSerre関手は射影を含むため、この証明だけから
四項部分圏そのものの非同値を結論してはならない。
その内在的な区別には、第\ref{sec:vk-intrinsic-classification-v36}節でGerstenhaber作用を用いる。
この節は固定非零量子代数のminiversal変形基底の計算にも依存しない。
\end{remark}

\subsection{非零CGG族の射影圏の完全な離散分類}

\begin{theorem}[本族内部での射影圏の同値分類]
\label{thm:CGG-categorical-classification-v35}
$m=2j+r\ge3$、$j,r\ge0$、$s\ne0$ とし、
\[
 \mathcal Q_{m,j}(\alpha,s)
 :=\operatorname{qgr}\mathcal S_{j,m-2j}(\alpha,s)
\]
と書く。同様の条件を満たす二組のパラメータに対し、次は同値である。
\begin{enumerate}
\item $m=m'$ かつ $j=j'$。
\item $\mathcal Q_{m,j}(\alpha,s)$ と
 $\mathcal Q_{m',j'}(\alpha',s')$ は $\mathbb C$ 線形アーベル圏として同値である。
\item その有界導来圏は $\mathbb C$ 線形三角圏として同値である。
\end{enumerate}
従って固定した重み $m$ で得られる相異なる射影圏の個数は
$\lfloor m/2\rfloor+1$ である。
\end{theorem}

\begin{proof}
同じ $(m,j)$ では、全ての $\alpha$ がZhang twistで結ばれ、
全ての非零 $s$ が次数付き同型で結ばれるので、(1)から(2)が従う。
(2)から(3)は同値を複体に適用すればよい。

(3)を仮定する。標準full exceptional collectionの長さは $m+3$ であり、
その類はGrothendieck群の基底をなすので
\[
 K_0\bigl(D^b(\mathcal Q_{m,j}(\alpha,s))\bigr)
 \simeq\mathbb Z^{m+3}.
\]
従って $m=m'$ である。
通常の三角同値が存在すれば、そのtilting対象の像とRickardの定理により
二つの有限Beilinson代数の間にdg Morita同値も存在する。

両方の $r=m-2j,r'=m-2j'$ が正なら、
定理\ref{thm:fixed-quantum-hull-v35}によるminiversal環の長さは
それぞれ $j+1,j'+1$ であり、導来不変性から $j=j'$ を得る。
一方のみ $r=0$ なら、定理\ref{thm:endpoint-hull-v36}により
そのhullは次元 $m+2$ の正則環である一方、他方のhullはArtin環である。
これはhullの導来不変性に矛盾する。
両方が零なら $j=j'=m/2$ である。これで(3)から(1)が従う。
\end{proof}

この分類の対象は、明示した接線共鳴CGG族とその $r=0$ 端点である。
全ての非可換重み付き射影空間の分類を主張するものではない。
重み三の四対象VK部分圏は、第\ref{sec:vk-intrinsic-classification-v36}節で別に分類する。

\begin{corollary}[可換点を含めた族]固定した $m$ において $s=0$ も許すと、
本族の射影圏は正確に $\lfloor m/2\rfloor+2$ 個の導来同値類を持つ。
\end{corollary}
\begin{proof}$s=0$ では全て同じ可換射影スタックの圏になる。
その $HH^1$ の次元は $N_m+9$、非零量子化では一であるため、
可換圏は全ての非零族から区別される。\end{proof}

\section{固定量子VK成分の三十五次元変形と標準補完}
\label{sec:quantum-vk-corners-v35}

この節では元の重み三へ戻り、標準六項列と四項部分列を同時に扱う。
まず四項部分圏の変形環と標準補完の局所像を決定し、
次節でGerstenhaber作用により四項圏自体を区別する。

$m=3$、$j=0,1$、$r=3-2j$、$s\ne0$ とし、
$\mathcal S=\mathcal S_{j,r}(\alpha,s)$ と置く。標準列の有限代数を
\[
 \Lambda=\operatorname{End}\left(\bigoplus_{a=0}^{5}\mathcal O_{\mathcal S}(a)\right),
 \qquad e=e_0+e_1+e_3+e_4,\qquad B=e\Lambda e
\]
とする。ここで $B$ は階数五十六の量子VK代数である。
変形は原始冪等元を持ち上げて固定した結合代数変形とする。

\begin{theorem}[固定量子VK代数と標準補完の局所環]
\label{thm:quantum-vk-hull-v35}
上の二つの共鳴について、全ての $\alpha$ と $s\ne0$ で
\[
 (\dim HH^0(B),\dim HH^1(B),\dim HH^2(B),\dim HH^3(B))
       =(1,1,35,39)
\]
かつ $HH^p(B)=0$（$p>3$）であり、そのminiversal環は
\[
 \boxed{R_B\simeq\mathbf C[[z_1,\ldots,z_{35}]]}
\]
となる。標準補完の制限写像について
\[
 HH^1(\Lambda)\xrightarrow{\sim}HH^1(B),\qquad
 HH^2(\Lambda)\hookrightarrow HH^2(B)
\]
が成立する。適合的なminiversal基底を選べば、補完の写像は閉埋込み
\[
 \operatorname{Spf}R_\Lambda\hookrightarrow\operatorname{Spf}R_B,
 \qquad R_\Lambda=\mathbf C[[\gamma]]/(\gamma^{j+1})
\]
として表せる。従って、その局所像は $j=0$ では被約点、
$j=1$ では三十五次元の滑らかな基底の中の長さ二の部分スキームである。
\end{theorem}

\begin{proof}
まずCY値の $\alpha$ を取り、量子化係数を $s=h$ とする。
特殊ファイバーを $\Lambda_0,B_0$ と書く。
公表稿VIII\cite{VKAmbient32}の標準補完の計算により、
\[
 HH^i(\Lambda_0)\xrightarrow{\sim}HH^i(B_0)\quad(i=0,1,2),
 \qquad \dim HH^1(B_0)=19,
 \qquad R_{B_0}\simeq\mathbf C[[u_1,\ldots,u_{53}]]
\]
である。ここで $i=0$ の両辺はスカラー中心である。

相対正規化Hochschild複体の制限は、各可合成列が属する対象を用いて定義され、
次数ごとに全射である。従って、有限自由 $\mathbf C[[h]]$ 加群の複体の完全列
\[
 0\longrightarrow K_h^\bullet
 \longrightarrow C_E^\bullet(\Lambda_h)
 \longrightarrow C_{eE}^\bullet(B_h)
 \longrightarrow0
\]
を得る。特殊ファイバーの長完全列と上の三つの同型から
$H^1(K_0)=H^2(K_0)=0$ となる。
微分行列の階数はgeneric fiberで減少しないので、
$K=\mathbf C((h))$ 上でも $H^1(K_h\otimes K)=H^2(K_h\otimes K)=0$ である。
非零係数の同型で固定した $s\ne0$ へ移すと、長完全列から
主張した $HH^1$ の同型と $HH^2$ の単射が従う。
定理\ref{thm:fixed-quantum-hull-v35}により、特に $\dim HH^1(B)=1$ である。

次に $B$ の変形環を決める。この四対象の相対複体の次元は
\[
 (\dim C^0,\dim C^1,\dim C^2,\dim C^3)=(4,620,1584,972),
 \qquad C^p=0\quad(p>3)
\]
であり、underlying双加群が一定なので全族で変わらない。
中心はどのファイバーでもスカラーであるため、$\operatorname{rank}d^0=3$ である。
可換基点での基底変換軌道の次元は
\[
 b_0=620-19-3=598.
\]
軌道方向を分離する形式切片により、$B_0$ における結合積スキームの完備局所環は
五十三個と五百九十八個の滑らかな変数、計六百五十一変数の形式冪級数環となる。
従ってこの点は次元六百五十一の滑らかな点である。
滑らかな軌跡の開性、または補題\ref{lem:arc-translation-v35}を使うと、
$B_h\otimes K$ も次元六百五十一の滑らかな点である。
一方、その軌道次元は
\[
 b_\eta=620-1-3=616.
\]
よってminiversal環は $651-616=35$ 変数の形式冪級数環である。
非零係数の同型と忠実な体拡大から、固定した $B$ についても同じ結論を得る。

この結果から $\dim HH^2(B)=35$ である。
複体のEuler特性は $4-620+1584-972=-4$ なので
\[
 1-1+35-\dim HH^3(B)=-4
\]
となり、$\dim HH^3(B)=39$ を得る。

最後にminiversal基底の写像を選ぶ。
$R_\Lambda$ 上のminiversal族をcornerへ制限し、
$R_B$ 上の族の形式的versalityによって各Artin商上で適合する写像を帰納的に選ぶ。
これにより連続局所準同型 $R_B\to R_\Lambda$ を得る。
その余接写像は $HH^2$ 制限の双対であり全射である。
完備局所環に対するtopological Nakayamaの議論から、環写像そのものが全射となる。
従って形式スキームの写像は閉埋込みである。
$j=1$ では形式座標を取り直すことで
\[
 R_B\longrightarrow R_\Lambda,\qquad
 z_1\longmapsto\gamma,\qquad z_2,\ldots,z_{35}\longmapsto0,
 \qquad \gamma^2=0
\]
と書ける。$j=0$ は剰余体への写像である。

全 $\alpha$ への拡張には、式\eqref{eq:finite-twist-iso-v35}の有限代数同型を用いる。
$T=1+sN$、$\alpha' =\alpha+\beta$ とし、通常のEnd合成 $b\star a$ に対して
各 $i\to k$ の射を $G_{ik}(a)=T^{-\beta i}(a)$ へ送る。
この同型は六頂点の冪等元を全て固定するため、四項cornerとその制限写像も同時に保つ。
従って結論が全 $\alpha$ へ移る。
\end{proof}

\begin{remark}[基点を移した後の補完条件]
可換基点では四対象と六対象の接空間は五十三次元の同型である。
これに対し、固定した非零量子化では四対象の接空間は三十五次元となり、
標準六項へ延長できる一次方向の空間は、$j=0$ で零次元、$j=1$ で一次元である。
後者の非零方向も二次で障害を持つ。
この変化は、可換基点と固定量子化で同じ接写像を使えないことを具体的に示す。
\end{remark}

本定理の閉埋込みは、適合的に選んだminiversal基底に関する非標準的な表示である。
自己同型を含む変形スタック全体の閉埋込みや、別の例外列による全ての補完の分類を
主張するものではない。二つの四項部分圏の導来非同値には、次節の内在的な作用計算を用いる。

\section{量子VK成分そのものの導来分類}
\label{sec:vk-intrinsic-classification-v36}

標準六項補完の変形環は、四項成分に付加した外側のデータを測る。
この節では四項圏だけから定まる不変量を用いて、
二つの非零量子VK成分の導来非同値を証明する。
同じHochschild群の次元と同じ滑らかな変形環を持つ二つの圏を、
外的な無限小自己同型の変形方向への作用によって区別する。

\subsection{一次元の $HH^1$ が作用する接空間}

\begin{lemma}[Hochschild作用の階数とJordan分割]
\label{lem:one-dimensional-HH1-action-v36}
$A$ を標数零の体上の有限次元代数とし、$\dim HH^1(A)=1$ とする。
非零元 $\xi\in HH^1(A)$ と整数 $p\ge0$ に対し
\[
 \rho_p(A)=\operatorname{rank}\bigl([\xi,-]:HH^p(A)\longrightarrow HH^p(A)\bigr)
\]
は $\xi$ の選択によらない。
この作用が冪零なら、そのJordan分割 $\lambda_p(A)$ も選択によらない。
これらはGerstenhaber構造を保つHochschildコホモロジーの同型で不変であり、
特に導来Morita不変量である。
\end{lemma}

\begin{proof}
別の非零元は $c\xi$（$c\ne0$）と書ける。
対応する作用は $c$ 倍なので、その階数は変わらない。
冪零作用のJordan分割は全ての冪の階数によって決まり、
$(c[\xi,-])^a=c^a[\xi,-]^a$ であるから同じ分割を持つ。
Gerstenhaber括弧を保つ同型はこれらの作用を共役する。
最後の主張にはHochschild複体の $B_\infty$ 構造の導来Morita不変性
\cite[Theorem~3.2 and Corollary]{KellerDerived2003}を使う。
\end{proof}

$HH^1(A)$ は外的導分の空間、$HH^2(A)$ は代数変形の接空間である。
従って $\rho_2(A)$ は、唯一の外的な無限小自己同型が動かす一次変形方向の次元を測る。
変形環が滑らかであっても、この作用は零である必要はない。

\subsection{量子VK代数上の対数導分}

$m=3$ とし、$j=0,1$、$r=3-2j$ とする。
\[
 B_{j,r}(\alpha,s)
 =\operatorname{End}\bigl(
 \mathcal O_{\mathcal S}\oplus\mathcal O_{\mathcal S}(1)
 \oplus\mathcal O_{\mathcal S}(3)\oplus\mathcal O_{\mathcal S}(4)\bigr),
 \qquad \mathcal S=\mathcal S_{j,r}(\alpha,s)
\]
と置く。$s\ne0$ では、式\eqref{eq:finite-twist-iso-v35}と量子化係数の規格化により
$\alpha=0,s=1$ を選べる。この代表を単に $B_{j,r}$ と書く。

四対象上の計算では、経路順の積 $a\star b$ を用いる。
$N=x\partial_y+2y\partial_z-q^jz^r\partial_w$、$q=xz-y^2$ に対して
$T=1+N$ の全てのスカラー冪は星積の自己同型である。
従って
\begin{equation}
 \delta=\left.\frac{d}{dt}\right|_{t=0}T^t
       =\log(1+N)
       =\sum_{k\ge1}\frac{(-1)^{k+1}}{k}N^k
 \label{eq:log-derivation-v36}
\end{equation}
は $B_{j,r}$ の導分となる。この和は各射空間上で有限である。
ここで $N$ 自体を星積の導分として用いていない。

$E=\bigoplus_{a\in\{0,1,3,4\}}\mathbb C e_a$、$J=\operatorname{rad}B$ とし、
相対正規化Hochschild複体を
\[
 C_E^p(B)=\operatorname{Hom}_{E\text{-}E}(J^{\otimes_E p},B)
\]
とする。導分 $\delta$ は全冪等元を零へ送り、各射空間を保つので、
この複体上に
\begin{equation}
 (L_\delta f)(a_1,\ldots,a_p)
 =\delta(f(a_1,\ldots,a_p))
 -\sum_{i=1}^p f(a_1,\ldots,\delta(a_i),\ldots,a_p)
 \label{eq:cochain-lie-action-v36}
\end{equation}
を定める。これは $[\delta,f]$ のGerstenhaber括弧であり、
Hochschild微分と可換する。

\begin{lemma}[冪零性と合成規約]
\label{lem:log-action-nilpotent-v36}
$L_\delta$ は各 $C_E^p(B)$ 上で冪零である。
また、通常のEnd合成 $b\star a$ へ規約を戻しても、
$HH^p(B)$ 上の作用の階数とJordan分割は変わらない。
\end{lemma}

\begin{proof}
補助次数を $\operatorname{wt}(x,y,z,w)=(0,1,2,m+r+1)$ とすると、
$N$ は次数を一つ下げる。
単項式からなる各cochain基底に、出力の補助次数から入力の補助次数の和を引いた整数を付ける。
式\eqref{eq:log-derivation-v36}、\eqref{eq:cochain-lie-action-v36}から
$L_\delta$ はこの整数を厳密に下げる。
各cochain空間は有限次元なので冪零性を得る。

反対代数上で入力を反転する写像を
$R_p f(a_1,\ldots,a_p)=f(a_p,\ldots,a_1)$ と置くと
\[
 d_{\mathrm{op}}R_p=(-1)^{p+1}R_{p+1}d,
 \qquad L_{\delta,\mathrm{op}}R_p=R_pL_\delta
\]
である。従って $R_p$ はcocycleとboundaryをそれぞれ対応させ、
コホモロジー上で作用を共役する。
\end{proof}

\subsection{有理数上の作用計算}

\begin{proposition}[量子VK成分のHochschild作用]
\label{prop:vk-HH-action-v36}
$B_{0,3}$、$B_{1,1}$ のいずれでも $[\delta]$ は $HH^1(B)$ の非零生成元である。
その $HH^2$ と $HH^3$ への作用のJordan分割は次の通りである。
\[
\begin{array}{c|c|c|r}
 p&(j,r)&\lambda_p(B)&\rho_p(B)\\\hline
 2&(0,3)&(13,10,7,5)&31\\
 2&(1,1)&(11,9,6,5,3,1)&29\\
 3&(0,3)&(14,11,8,5,1)&34\\
 3&(1,1)&(13,10,7,5,3,1)&33
\end{array}
\]
従って $HH^2$ 上の固定部分空間の次元は、それぞれ四と六である。
\end{proposition}

\begin{proof}
以下は有理数上の有限行列計算による証明である。
対象を $a_0<\cdots<a_p$（$a_i\in\{0,1,3,4\}$）とし、
各入力には重み付き次数 $a_i-a_{i-1}$ の単項式、
出力には次数 $a_p-a_0$ の単項式を取る。
これにより全cochain基底を一意に列挙できる。
\[
 (\dim C^0,\dim C^1,\dim C^2,\dim C^3)=(4,620,1584,972),
 \qquad C^p=0\quad(p>3).
\]
星積\eqref{eq:CGG-product-v35}をこの基底で評価し、
\[
\begin{split}
 (df)(a_1,\ldots,a_{p+1})={}&a_1\star f(a_2,\ldots,a_{p+1})\\
 &+\sum_{i=1}^p(-1)^i
 f(a_1,\ldots,a_i\star a_{i+1},\ldots,a_{p+1})\\
 &+(-1)^{p+1}f(a_1,\ldots,a_p)\star a_{p+1}
\end{split}
\]
から微分行列を作る。式\eqref{eq:cochain-lie-action-v36}が作用行列を与える。
全ての係数は有理数であり、打切り誤差や数値的階数判定はない。
両方の代数で
\[
 \operatorname{rank}d^1=616,\qquad
 \operatorname{rank}d^2=933
\]
となる。微分の二乗が零であることと、$L_\delta$ が微分と可換することを確認し、
$\ker d^2/\operatorname{im}d^1$ および $C^3/\operatorname{im}d^2$ 上へ作用を降ろす。
得られた三十五次および三十九次の有理行列の各冪の階数は、次表の通りである。
\[
\begin{array}{r|rr|rr}
 &\multicolumn{2}{c|}{HH^2}&\multicolumn{2}{c}{HH^3}\\
 a &(0,3)&(1,1)&(0,3)&(1,1)\\\hline
0&35&35&39&39\\
1&31&29&34&33\\
2&27&24&30&28\\
3&23&19&26&23\\
4&19&15&22&19\\
5&15&11&18&15\\
6&12&8&15&12\\
7&9&6&12&9\\
8&7&4&9&7\\
9&5&2&7&5\\
10&3&1&5&3\\
11&2&0&3&2\\
12&1&0&2&1\\
13&0&0&1&0\\
14&0&0&0&0\\
\end{array}
\]
$r_a=\operatorname{rank}L_\delta^a$ とすると、長さ $a$ のJordan blockの個数は
$r_{a-1}-2r_a+r_{a+1}$ であり、表示した分割を得る。

さらに、第一階数は商座標を用いない方法でも確認できる。
$K$ を $Z^2=\ker d^2$ の全基底を列に持つ行列とすると
\[
 \rho_2(B)=\operatorname{rank}[\,d^1\mid L_\delta K\,]
                        -\operatorname{rank}d^1.
\]
両者で $\dim Z^2=651$、拡大行列の階数は順に $647,645$ となり、
$647-616=31$、$645-616=29$ を得る。
この独立な計算も厳密な有理数演算で行える。
計算コードと、全作用行列および階数列を含む検証データを付属資料に収める。

内的導分はHochschildコホモロジーに零作用するので、
得られた非零作用から $[\delta]\ne0$ が従う。
定理\ref{thm:quantum-vk-hull-v35}の $\dim HH^1(B)=1$ と合わせれば、
$[\delta]$ が生成元であることも分かる。
\end{proof}

\subsection{三つの圏の完全な区別}

\begin{theorem}[明示的な重み三CGG族内のVK分類]
\label{thm:vk-three-categories-v36}
$j,j'\in\{0,1\}$、$r=3-2j$、$r'=3-2j'$、
$\alpha,\alpha',s,s'\in\mathbb C$ とする。
$B=B_{j,r}(\alpha,s)$、$B'=B_{j',r'}(\alpha',s')$ について、次は同値である。
\begin{enumerate}
\item $s=s'=0$、または $ss'\ne0$ かつ $j=j'$ である。
\item $B$ と $B'$ は $\mathbb C$ 代数として同型である。
\item $\operatorname{mod}B$ と $\operatorname{mod}B'$ は
$\mathbb C$ 線形アーベル圏として同値である。
\item $D^b(\operatorname{mod}B)$ と $D^b(\operatorname{mod}B')$ は
$\mathbb C$ 線形三角圏として同値である。
\end{enumerate}
従って、この明示族のVK成分は、可換成分と二つの非零共鳴成分からなる
ちょうど三つの導来同値類を持つ。
\end{theorem}

\begin{proof}
$s=s'=0$ のとき星積は通常積であり、$j,\alpha$ に依存しない。
$ss'\ne0$ かつ $j=j'$ なら、全非零係数の規格化と
式\eqref{eq:finite-twist-iso-v35}の冪等元固定同型が $B\simeq B'$ を与える。
従って(1)から(2)、また明らかに(2)から(3)、(3)から(4)が従う。

(4)を仮定する。
これらの代数は四つの順序付けられた頂点を持つ有向代数であり、
相対bar分解は長さ三以下である。特に有限大域次元を持ち、
$D^b(\operatorname{mod}B)=\operatorname{Perf}(B)$ である。
通常の三角同値が存在すれば、その同値によるtilting対象 $B$ の像を取り、
Rickardの定理\cite{RickardMorita1989,RickardDerived1991}により
$B$ と $B'$ を結ぶ両側tilting複体、従ってdg Morita同値が存在する。
ここで、与えられた三角同値そのものが標準関手であることを仮定する必要はない。

可換の場合の $\dim HH^1$ は十九、非零量子化では一であるから、
$s=0$ と $s'\ne0$ は同値になれない。
両方が非零で $j\ne j'$ なら、命題\ref{prop:vk-HH-action-v36}より
$\rho_2$ は三十一と二十九になる。
補題\ref{lem:one-dimensional-HH1-action-v36}の導来不変性に矛盾する。
従って(1)が成立する。
\end{proof}

この比較では、三角同値に標準四項列や外側への埋込みを保存する条件を課していない。
非同値性は四項圏の内部だけで検出される。
分類対象は明示した $m=3$ のCGG族であり、全ての非可換VK圏や、
異なる例外列を使った全ての半直交埋込みの分類ではない。

\begin{corollary}[同じ滑らかな変形環を持つ非同値なVK圏]
\label{cor:vk-same-hull-distinct-action-v36}
二つの非零成分は、Hochschild群の次数別次元
$(1,1,35,39)$ とminiversal環 $\mathbb C[[z_1,\ldots,z_{35}]]$ を共有するが、
互いに導来同値ではない。
外的な無限小自己同型が固定する一次変形方向は、それぞれ四次元と六次元である。
\end{corollary}

\begin{corollary}[滑らかなhullと非可換なHochschild dg Lie構造]
\label{cor:vk-not-homotopy-abelian-v36}
二つの非零量子VK代数について、全相対Hochschild dg Lie代数
\[
 \mathfrak g_B=C_E^\bullet(B)[1]
\]
はhomotopy abelianではない。また、二つの $\mathfrak g_B$ は
 dg Lie代数としてquasi-isomorphicではない。
ここでは $C_E^0(B)$ を次数 $-1$ に含めた全複体を使う。
\end{corollary}

\begin{proof}
$H^0(\mathfrak g_B)=HH^1(B)$、$H^1(\mathfrak g_B)=HH^2(B)$ であり、
両者の括弧は命題\ref{prop:vk-HH-action-v36}により非零である。
アーベルdg Lie代数のコホモロジー括弧は零であり、
dg Lie quasi-isomorphismのzigzagはコホモロジーの次数付きLie構造を保つ。
従ってhomotopy abelianではない。
また二つの括弧の階数は異なるので、両dg Lie代数間のquasi-isomorphismも存在しない。
\end{proof}

この系では $C_E^{\ge1}(B)[1]$ への素朴な切捨てを行っていない。
その切捨てでは次数零コホモロジーが外的導分でなく、冪等元を固定する導分全体となり、
上で使った一次元の $HH^1$ と一致しなくなる。
無障害な変形のminiversal環と、自己同型の作用も保持する全Hochschild構造を
区別する必要がある。

この結果は、標準補完の局所像が被約点と長さ二の点に分かれるという
前節の結果を内在化する一つの方法を与える。
その証明では補完を使わず、変形方向に対する自己同型作用を保ったHochschild構造を用いた。
miniversal環だけを取り出すと失われる情報が、四項圏の同値分類に実際に寄与する。

\section{外的導分の消失と対称性保持の二次障害}
\label{sec:vk-symmetry-jump-v37}

前節の二つの量子VK代数は、いずれも三十五次元の滑らかなminiversal基底を持つ。
しかし、基点の一次元の外的導分が変形後にも残る条件には障害がある。
本節ではその条件をFittingイデアルで定義し、一般変形における外的導分の消失と、
対称性を保持する部分の二次近傍を決定する。
ここで対称性の保持とは、基点の非零外的導分を持ち上げる導分が存在することをいう。
指定した $\mathbf G_a$ 作用や、導分の局所冪零性そのものを保持する条件とは区別する。

\subsection{一列の微分が記述する外的導分の跳躍}

$B=B_{0,3}$ または $B_{1,1}$ とし、miniversal族 $B_R$ を選ぶ。
その基底は $R=\mathbf C[[t_1,\ldots,t_{35}]]$ である。
冪等元を固定した相対Hochschild複体は、有限自由 $R$ 加群の複体
\[
 C_R^\bullet=C_E^\bullet(B_R),\qquad
 (\operatorname{rank}C_R^0,\operatorname{rank}C_R^1,
 \operatorname{rank}C_R^2,\operatorname{rank}C_R^3)=(4,620,1584,972)
\]
であり、四次以上は零である。

\begin{theorem}[外的導分の跳躍を定める一列]
\label{thm:vk-jump-column-v37}
ある列 $a=(a_1,\ldots,a_{35})^t\in\mathfrak m_R R^{35}$ により
\begin{equation}
 C_R^\bullet\simeq
 R[0]\oplus
 \bigl(R\xrightarrow{\ a\ }R^{35}\bigr)_{1,2}
 \oplus R^{39}[-3]
 \label{eq:vk-minimal-relative-complex-v37}
\end{equation}
という自由複体のホモトピー同値がある。
添字 $1,2$ は二項を置くコホモロジー次数を表す。
次数一の基底を $[\delta]$ とし、$T_0\operatorname{Spf}R=HH^2(B)$ と同定すると
\[
 (da)_0(u)=-[\delta,u].
\]
イデアル
\begin{equation}
 I_B=(a_1,\ldots,a_{35})
     =\operatorname{Fitt}_{34,R}H^2(C_R^\bullet)
 \label{eq:vk-symmetry-Fitting-v37}
\end{equation}
は縮約と基底の選択によらない。
$Z_B=\operatorname{Spf}(R/I_B)$ と置くと、その接空間は
\[
 T_0Z_B=\ker\bigl([\delta,-]:HH^2(B)\to HH^2(B)\bigr).
\]
特に接空間の次元は $B_{0,3}$ で四、$B_{1,1}$ で六である。
\end{theorem}

\begin{proof}
四項代数の結合積スキームは、相対二cochainのアフィン空間に結合則を課して得られる。
その接写像は、符号を除いて $d^2:C^2\to C^3$ である。
定理\ref{thm:quantum-vk-hull-v35}の証明により、$B$ はこのスキームの
次元六百五十一の滑らかな点である。
従って、その滑らかな近傍では結合則のJacobianの像は階数
$1584-651=933$ の直和因子となる。
これを $R$ へ引き戻すと、$d^2$ の階数九百三十三の部分を可縮複体として除ける。
単に特殊ファイバーの階数だけから、残余微分が零だと推論しているのではない。

$d^0$ は頂点上のスカラーの差であり、全族で階数三の直和因子を持つ。
さらに基点での $d^1$ の階数六百十六の単元小行列を消す。
残る自由加群の階数は次数零から順に $1,1,35,39$ であり、
上の $d^2$ の一定階数から、残る微分は次数一から二への一列だけである。
これで式\eqref{eq:vk-minimal-relative-complex-v37}を得る。

積を $\mu+\varepsilon u$ と変形すると、固定した導分 $\delta$ に対して
\[
 d^1_{\mu+\varepsilon u}\delta
 =-\varepsilon[\delta,u]
\]
である。単元消去をコホモロジーへ降ろせば、表示した一次項を得る。
二次コホモロジー加群は $\operatorname{coker}(a:R\to R^{35})$ なので、
その第三十四Fittingイデアルは列の全成分が生成するイデアルである。
Fittingイデアルの表示によらない性質\cite[Tags 07Z8, 07ZA]{StacksFitting37}から
縮約の選択によらないことが従う。
最後の接空間の記述は $a$ の線形項から得られる。
\end{proof}

有限局所基底への写像 $R\to A$ について、$a$ の像が零であることは、
基点の $[\delta]$ を生成元として持つ外的導分を $A$ 上に持ち上げられることと同値である。
実際、縮約後の次数一加群は階数一であり、基点で非零の元は単元である。
その元が列 $a$ によって零へ写ることは、$a=0$ と同値になる。
したがって $Z_B$ は、外的導分の持上げの存在条件を全Artin基底上で記述する。
一方、持上げる導分そのものをデータとして選ぶ変形関手とは同一ではない。

\begin{corollary}[一般変形のHochschild群]
\label{cor:vk-generic-symmetry-breaking-v37}
$\operatorname{Spec}R$ の $V(I_B)$ の外では、各剰余体上のHochschild次元は
\[
 (1,0,34,39)
\]
であり、その点のminiversal基底は三十四次元の滑らかな形式空間である。
$V(I_B)$ 上では次元は $(1,1,35,39)$ となる。
また、初速度 $u\notin\ker[\delta,-]$ を持つ任意の形式弧では、
一般ファイバーの外的導分は零である。
\end{corollary}

\begin{proof}
一次項の階数は三十一または二十九なので $a\ne0$ である。
式\eqref{eq:vk-minimal-relative-complex-v37}を各剰余体へtensorすれば、
$a$ が零か非零かによって表示した二つの次元列を得る。
結合積スキームは滑らかな近傍の中にある。
その軌道の次元は $620-3-\dim HH^1$ であるから、
外的導分が消えた点の変形基底は $651-617=34$ 次元となる。
形式弧については $a(t)$ の一次係数が非零なので、$\mathbf C((t))$ 上で
列の階数が一になる。
\end{proof}

ここで $R$ 上の相対 $H^1(C_R^\bullet)$ は零である。
実際、$R$ は整域で $a\ne0$ だから $R\to R^{35}$ は単射である。
それでも特殊ファイバーの $HH^1$ は一次元であり、
この跳躍は式\eqref{eq:vk-minimal-relative-complex-v37}の導来基底変換に保持される。
相対コホモロジー加群を先に取ってから通常のtensor積を取る操作では検出できない。

\subsection{対称性保持の二次障害}

$L=[\delta,-]$ とし、$K_B=\ker(L:HH^2\to HH^2)$、
$Q_B=\operatorname{coker}(L:HH^2\to HH^2)$ と置く。
両空間の次元は同じであり、$k_B=4$ または $6$ である。
一次項の独立な $35-k_B$ 個を形式的に消去すると
\begin{equation}
 R/I_B\simeq
 \mathbf C[[v_1,\ldots,v_{k_B}]]/(f_1,\ldots,f_{k_B}),
 \qquad f_i\in(v_1,\ldots,v_{k_B})^2
 \label{eq:vk-jump-small-equations-v37}
\end{equation}
となる。これらの二次項を以下で計算する。

相対二cochainについて
\[
 (u\circ u)(a,b,c)=u(u(a,b),c)-u(a,u(b,c))
\]
とする。$u\in Z^2(B)$ が $K_B$ の元を表すとき、
\[
 d^1\eta=L_\delta u,\qquad d^2v=u\circ u
\]
を満たす $\eta\in C^1$、$v\in C^2$ を選べる。
第二の可解性は代数変形の無障害性による。
積と導分を
\[
 \mu_t=\mu+tu+t^2v,\qquad
 \delta_t=\delta+t\eta+t^2\theta
\]
と書けば、$t^2$ における導分条件は
\[
 d^1\theta=L_\delta v+L_\eta u
\]
である。従って二次障害は
\begin{equation}
 \mathfrak q_B(u)=
 [L_\delta v+L_\eta u]\bmod\operatorname{im}L
 \quad\in Q_B
 \label{eq:vk-secondary-symmetry-obstruction-v37}
\end{equation}
で与えられる。

\begin{lemma}[二次障害の独立性]
\label{lem:vk-quadratic-obstruction-v37}
式\eqref{eq:vk-secondary-symmetry-obstruction-v37}は、二次写像
$\mathfrak q_B:K_B\to Q_B$ を定める。
その成分は式\eqref{eq:vk-jump-small-equations-v37}の二次項である。
\end{lemma}

\begin{proof}
$v$ の選択を変えると差は二cocycleとなり、表示した類の差は
$\operatorname{im}L$ に入る。
$\eta$ の選択を変えると差は導分である。
$HH^1(B)=\mathbf C[\delta]$ であり、内的導分はコホモロジーに零作用するため、
差の $HH^2$ 作用は $[\delta,u]=0$ のスカラー倍となる。
従って選択によらない。
基底変換による $u$ の代表の変更は、積と導分を同時に共役することで
対応する持上げを与えるので、この存在障害を変えない。

二cocycleの代表と二つの線形方程式の解を線形に選べば、
$v$ は $u$ に二次、$\eta$ は一次に依存する。
これで $\mathfrak q_B$ が二次写像であることが分かる。
第二次の積に任意の二cocycleを加える自由度が、
一列 $a$ の線形像を消去する自由度と一致する。
従って $Q_B$ への射影が、独立な一次方程式の消去後に残る二次方程式を与える。
可換な複数の変数について偏極すれば、全二次項について同じ結論を得る。
\end{proof}

\begin{theorem}[外的導分を保つ部分の二次近傍]
\label{thm:vk-symmetry-second-neighborhood-v37}
$B_{0,3}$ では $\mathfrak q_B=0$ である。
$B_{1,1}$ では $\mathfrak q_B$ の像は一次元であり、
\[
 \mathfrak q_B(u)=0
 \quad\Longleftrightarrow\quad
 u\in\ker L\cap\operatorname{im}L.
\]
後者の交叉は五次元である。
従って、$Z_B$ の原点の極大イデアルを $\mathfrak n$ とすると、
\[
 \widehat{\mathcal O}_{Z_B,0}/\mathfrak n^3\simeq
 \begin{cases}
 \mathbf C[v_1,\ldots,v_4]/(v_1,\ldots,v_4)^3,&B=B_{0,3},\\[2pt]
 \mathbf C[v_1,\ldots,v_6]/\bigl(v_6^2,(v_1,\ldots,v_6)^3\bigr),&B=B_{1,1}.
 \end{cases}
\]
特に $Z_{B_{1,1}}$ は原点で正則でなく、次元は五以下である。
\end{theorem}

\begin{proof}
前節と同じ有理数基底で、$d^1,d^2,L_\delta$ を作る。
$HH^2$ の商座標上の作用行列からその核を取り、各核ベクトルについて
$d^1\eta=L_\delta u$ を解く。
四次元または六次元の核の全対称二次単項式について
$d^2v=u\circ u$ を解き、式\eqref{eq:vk-secondary-symmetry-obstruction-v37}を評価する。
必要な二次単項式の個数はそれぞれ十と二十一である。

境界のpivotによる商座標と通常の行簡約で基底を一意に選ぶと、
$B_{0,3}$ の四成分は全て零となる。
$B_{1,1}$ の六成分は全て $v_6^2$ の定数倍となり、例えば第二成分は
\[
 \frac{432000}{2468041}v_6^2
\]
である。自然写像 $K_B\to Q_B$ も第六座標だけに依存し、その第三成分は
$\frac{2700}{1571}v_6$ となる。
従って、この非零な二次障害の零集合は
自然写像 $\ker L\to\operatorname{coker}L$ の核、すなわち
$\ker L\cap\operatorname{im}L$ に一致する。
これは前節のJordan分割に長さ一のblockがちょうど一つあることとも整合する。

独立な検算では、二次の積補正を先に選ばず、未知数 $(v,\theta)$ に対する
\[
 \begin{pmatrix}d^2&0\\L_\delta&-d^1\end{pmatrix}
 \begin{pmatrix}v\\\theta\end{pmatrix}
 =\begin{pmatrix}u\circ u\\-L_\eta u\end{pmatrix}
\]
を全cochain空間で直接調べる。
この行列は $2556\times2204$、階数は順に $1580,1578$ である。
第一の場合は十個の右辺が全て像に入り、
第二の場合は $v_6^2$ の右辺だけが像に入らず、
全右辺を加えたときの階数の増分は一となる。
従って商座標による計算と、全空間における同時可解性の計算が一致する。
全ての持上げ方程式、cocycle条件、行列と有理係数は付属のコードと証明書に保持する。

補題\ref{lem:vk-quadratic-obstruction-v37}により、これらは跳躍部分の全二次項である。
第二の場合には二次項が非零の関係があるため、
埋込次元六に対してKrull次元は五以下となり、正則でない。
\end{proof}

この定理は、$B_{1,1}$ の不変な一次方向のうち、
$\ker L\cap\operatorname{im}L$ に入らない方向では、
代数自身を二次へ延長できても、基点の外的導分を二次へ持ち上げられないことを意味する。
$B_{0,3}$ では全ての不変な一次方向について二次までの持上げが存在する。
しかし、$B_{0,3}$ の $Z_B$ が滑らかであることや、
$B_{1,1}$ の全完備局所環が単に $(v_6^2)$ だけで定まることは、ここからは従わない。
また、二次項が平方であることだけでは、全 $Z_{B_{1,1}}$ の非被約性も結論できない。
三次以上の残余方程式と、標準補完の一次方向がこの内在的な五次元部分に属するかは、
次に調べるべき別の問題として残る。

\section{対称性保持部分の滑らかな成分と有限な横断構造}
\label{sec:vk-symmetry-families-v38}

前節の二次障害は、対称性保持部分の接空間のうち、実際の形式族で実現できる
方向を区別する。本節では、その二次障害が消える全方向を同時に積分する。
得られる族は形式的な逐次構成ではなく、積の係数がパラメータの高々二次式である
平坦な代数族である。しかも、外的導分の存在だけでなく、明示的な
$\mathbf G_a$ 作用を保持する。

\begin{theorem}[対称性を保持する全次数の族]
\label{thm:vk-Ga-families-v38}
標準量子VK代数 $B_{0,3}$ と $B_{1,1}$ について、次が成り立つ。
\begin{enumerate}
\item $B_{0,3}$ の対称性保持部分全体は
\[
 Z_{B_{0,3}}\simeq\operatorname{Spf}\mathbf C[[u_1,u_2,u_3,u_4]]
\]
である。四つのパラメータに線形な積と導分を持つ平坦な代数族が、この同型を実現する。
\item $Z_{B_{1,1}}$ には五次元の滑らかな閉形式部分 $Y$ が存在し、
\[
 T_0Y=\ker L_\delta\cap\operatorname{im}L_\delta
       \subset HH^2(B_{1,1})
\]
である。この部分は、五つのパラメータに高々二次な積と線形な導分を持つ
平坦な代数族で実現される。特に
\[
 \dim Z_{B_{1,1}}=5,\qquad
 \operatorname{edim}_0 Z_{B_{1,1}}=6.
\]
\item 上の両族で、保持する導分 $\delta_{\boldsymbol u}$ は
$\delta_{\boldsymbol u}^{12}=0$ を満たす。従って
$\exp(t\delta_{\boldsymbol u})$ は基底全体上の $\mathbf G_a$ 作用を定める。
次数一の射では $\delta_{\boldsymbol u}(y)=x$ であり、
各基底上で基点の非零外的導分を持ち上げる。
\end{enumerate}
\end{theorem}

\subsection{原始射の片側の合成だけを動かす}

四対象の次数を $(0,1,3,4)$ とする。
$w_0\in\operatorname{Hom}(0,3)$ は原始射、すなわち
$\operatorname{Hom}(0,1)\operatorname{Hom}(1,3)$ の像に含まれない射である。
次数一の空間を $V=\langle x,y,z\rangle$、
純粋な $x,y,z$ の四次単項式の空間を $P_4$ と書く。
相対二cochainの部分空間
\[
 W=\operatorname{Hom}_{\mathbf C}(\mathbf C w_0\otimes V,
             \operatorname{Hom}(0,4)),\qquad
 W_0=\operatorname{Hom}_{\mathbf C}(\mathbf C w_0\otimes V,P_4)
\]
を、それ以外の入力では零とすることで全cochain空間へ埋め込む。
その次元はそれぞれ $54,45$ である。

\begin{lemma}[原始射に支持を持つ積補正]
\label{lem:vk-primitive-supported-products-v38}
$F,G\in W$ に対して
\[
 d^2F=0,\qquad F\circ G=0
\]
である。従って任意の可換基底上で、任意の $W$ 値の補正 $F$ により
$\mu+F$ は結合積となる。
また、一次cochain $E$ が $w$ を含む射を純粋多項式へ送り、
純粋多項式上で零であるとする。このとき $G\in W_0$ に対して
$[E,G]=0$ である。
\end{lemma}

\begin{proof}
三本の非恒等射の可合成列は $0\to1\to3\to4$ だけである。
その最初の二本の積は純粋な $x,y,z$ の三次式であり、$w_0$ 成分を持たない。
従って $F$ は、この三重合成の途中に現れない。
同じ支持の理由から、$F\circ G$ も零である。
最後の主張では、$E$ は $G$ の純粋な出力を零へ送り、
$G$ は $E(w_0)$ を第一入力に持つと零となる。
次数一の入力に対する $E$ も零なので、全項が消える。
\end{proof}

\subsection{有限個の恒等式による全次数の証明}

前節で固定した有理数基底の不変cocycleを $U_i$、
その導分の一次補正を $E_i$ とする。
$B_{0,3}$ では四本全てを用い、$B_{1,1}$ では
$\ker L_\delta\cap\operatorname{im}L_\delta$ の五本を用いる。
これらは
\[
 d^1E_i=[\delta,U_i]
\]
を満たす。
全 $U_i$ は $W$ に属し、全 $E_i$ は $w$ を含む射を純粋多項式へ送る。
$B_{0,3}$ ではさらに全 $U_i$ が $W_0$ に属する。
従って補題\ref{lem:vk-primitive-supported-products-v38}から
\[
 \mu_{\boldsymbol u}=\mu+\sum_i u_iU_i,
 \qquad
 \delta_{\boldsymbol u}=\delta+\sum_i u_iE_i
\]
が、第一の場合の求める族である。

第二の場合には
\[
 R_{ij}=
 \begin{cases}
 [E_i,U_j]+[E_j,U_i],&i<j,\\
 [E_i,U_i],&i=j
 \end{cases}
\]
が $W_0$ に属する。
$L_\delta:W_0\to W_0$ の行列は $45\times45$、階数は $38$ であり、
全ての $R_{ij}$ はその像に入る。
そこで有理数上で
\[
 [\delta,V_{ij}]=-R_{ij},\qquad V_{ij}\in W_0
\]
を解き、
\begin{equation}
 \mu_{\boldsymbol u}
 =\mu+\sum_{i=1}^{5}u_iU_i
       +\sum_{1\le i\le j\le5}u_iu_jV_{ij},
 \qquad
 \delta_{\boldsymbol u}=\delta+\sum_{i=1}^{5}u_iE_i
 \label{eq:vk-exact-quadratic-family-v38}
\end{equation}
と置く。
積補正は全て $W$ に属するので、結合則は補題から全次数で成立する。
導分条件の一次項と二次項は上の方程式で零となり、
三次項は $[E_i,V_{ab}]=0$ により零となる。
それより高い次数の項は存在しない。
従って、これは二次までの打切りではなく、多項式基底全体上の厳密な族である。

この構成で必要な有限行列の検証は、次の通りである。
\[
\begin{array}{c|c|c|c|c}
 (j,r)&\#\{U_i\}&\operatorname{rank}[d^1\mid U_i]&
 \operatorname{rank}\langle R_{ij}\rangle&\#\{V_{ij}\text{ の非零係数}\}\\\hline
 (0,3)&4&620&0&0\\
 (1,1)&5&621&7&141
\end{array}
\]
ここで $\operatorname{rank}d^1=616$ である。
従って両族のKodaira--Spencer写像は単射であり、
その像はそれぞれ $\ker L_\delta$ と
$\ker L_\delta\cap\operatorname{im}L_\delta$ に一致する。
行列 $L_\delta|_{W_0}$ の階数だけで可解性を推論しているのではなく、
全ての右辺について表示した方程式を有理数上で検証する。
全 $U_i,E_i,V_{ij}$、支持、cocycle条件、全次数を保証する消滅恒等式は
付属の読み取り専用検証コードと証明書に保存する。

\begin{proof}[定理\ref{thm:vk-Ga-families-v38}の証明の完結]
各族の導分は、純粋多項式からなる部分と、$w$ を一つ含む商からなる
二段のフィルトレーションを保つ。
その随伴次数上の作用は基点の $\delta=\log(1+N)$ と同じであり、
追加した $E_i$ はこのフィルトレーションを真に下げる。
次数 $d\le4$ の純粋部分では冪零指数は $2d+1$ 以下、
$w$ を含む商では $2(d-3)+1$ 以下である。
従って全空間で十二乗が零となる。
次数一での $\delta(y)=x$ は補正されない。
頂点冪等元を固定する内的導分は各矢空間にスカラーとして作用するので、
この導分は各基底上で基点の非零外的導分を持ち上げる。

以上の族は $Z_B$ へ因子化する。
第一の場合、四次元の滑らかな形式空間から $Z_B$ への写像は接空間で同型である。
$Z_B$ 自体も埋込次元四であるから、完備局所環上の形式的逆関数定理により
この写像は同型となる。
より明示的には、四変数形式冪級数環からの表示と、この族への写像を合成すると、
一次項が可逆な四変数環の自己準同型となるため、表示の核は零である。

第二の場合は、五変数族の接写像の単射性により閉形式埋込みを得る。
従って $\dim Z_{B_{1,1}}\ge5$ である。
前節の非零二次関係による上界 $5$ と合わせて次元は五に等しい。
埋込次元六は前節の接空間の計算による。
\end{proof}

\subsection{六変数の未知関係を一変数の有限な構造へまとめる}

五次元の滑らかな部分が実際に存在することは、
残る高次構造を次の限られた形へ整理する。
ここでは被約性や平坦性を仮定しない。

\begin{proposition}[対称性保持部分の有限な正規形]
\label{prop:vk-jump-finite-normal-form-v38}
$A=\mathbf C[[u_1,\ldots,u_5]]$ とする。
ある $b\in\mathfrak m_A^2$ とイデアル $J\subset\mathfrak m_A^2$ が存在し、
\begin{equation}
 \widehat{\mathcal O}_{Z_{B_{1,1}},0}
 \simeq A[[t]]/\bigl(t(t-b),\ tJ\bigr)
 \simeq A\oplus(A/J)t,\qquad t^2=bt
 \label{eq:vk-jump-finite-normal-form-v38}
\end{equation}
となる。
この表示で $Y$ は $t=0$ に対応する。
従って $Z_{B_{1,1}}\to\operatorname{Spf}A$ は切断を持つ有限な射であり、
その加群は高々二個の元で生成される。
この主張は $A$ 上で階数二の平坦族であることを意味しない。
\end{proposition}

\begin{proof}
前節の六変数表示を $S=\mathbf C[[u_1,\ldots,u_5,t]]$ とする。
滑らかな閉形式部分 $Y$ を座標変更によって $t=0$ とすることができる。
その接空間は二次障害の零超平面であるから、定義イデアル $I$ は
$(t)$ に含まれ、全ての二次項は $t^2$ の定数倍である。
特に $f=tg\in I$ で $g=t+O((u,t)^2)$ となる元が存在する。
形式的陰関数定理により
\[
 g=\text{単元}\cdot(t-b(u)),\qquad b\in\mathfrak m_A^2.
\]
単元を除いて $t(t-b)\in I$ としてよい。

$I=tK$ と書くと $t-b\in K$ である。
$S/(t-b)\simeq A$ における $K$ の像を $J$ と置けば
$K=(t-b,J)$、従って $I=(t(t-b),tJ)$ となる。
$I$ の二次項に $tu_i$ がないことから $J\subset\mathfrak m_A^2$ である。
Weierstrass除法により、まず $S/(t(t-b))$ の全ての剰余類は
$a(u)+c(u)t$、$a,c\in A$ と一意に書ける。
さらに $tJ$ で割ると、第二係数だけを $J$ で割ることになる。
これで表示した $A$ 加群と積を得る。
\end{proof}

表示\eqref{eq:vk-jump-finite-normal-form-v38}において、
$J=0,b=0$ なら滑らかな五次元部分の二重構造であり、
$J=0,b\ne0$ なら接している二つの滑らかな五次元枝となる。
$J\ne0$ の場合には $t=b, J=0$ 上の低次元の追加支持や埋込構造があり得る。
五次元の既約成分は高々二つである。
この段階の正規形だけからは、これらの可能性を区別できない。
本稿の $B_{1,1}$ については、定理\ref{thm:vk-double-jump-v39}で
平方零方向を含む全次数の族を構成し、実際に $b=J=0$ と決定する。
本節では、その証明に必要な滑らかな被約部分と全次数の $\mathbf G_a$ 族、
および全環を制限する有限な正規形を準備した。

\section{全ての重みで内在的に残る純粋部分}
\label{sec:vk-all-weights-filtered-v40}

重み三で得られた還元は、有限行列の偶然ではない。
本節では $j,r\ge0$、$m=2j+r\ge3$ とし、$s=1$, $\alpha=0$ における
CGG代数 $\mathcal S_{j,r}$ の $(0,1,m,m+1)$ 切出しを $B_{j,r}$ と書く。
$x,y,z$ だけからなる三変数部分の同じ切出しを $N_m$ とする。
$N_m$ は $j,r$ によらない。以下の主張には $r=0$ の偶数端点も含む。
\[
 n_d=\binom{d+2}{2},\qquad
 k_m=3n_{m+1}-2n_m+7=\frac{m^2+9m+28}{2}
\]
と置く。変形と同型は、特殊ファイバーの同定および四つの頂点冪等元を固定する。

\begin{theorem}[全重みの標準部分と滑らかな相対変形]
\label{thm:vk-all-weights-filtered-v40}
次が成り立つ。
\begin{enumerate}
\item $B_{j,r}$ の任意のArtin局所基底上の変形において、$N_m$ は自由直和因子の
部分代数として一意に持ち上がる。その持上げは同型と基底変換に自然である。
具体的には、二つの隣接二段合成と三段合成の像が、それぞれ階数
$n_m,n_m,n_{m+1}$ の純粋部分を与える。
\item $w$ の個数によるフィルトレーションを保存する表示を取れば、
保存的な二つの変形を結ぶ標識付き同型、および頂点冪等元を零にする導分も保存的である。
\item 純粋部分の変形を一つ固定し、その部分との同定も固定した相対変形問題は、
形式的に滑らかな次元 $k_m$ の形式空間である。
相対自己同型は存在しない。
従って変形groupoid間の制限は、表現可能な形式的に滑らかな射であり、その相対次元は $k_m$ である。
\end{enumerate}
\end{theorem}

最後の主張では純粋部分の同定も固定する。
通常のminiversal環の分解には、純粋部分の無限小自己同型を持ち上げる条件も必要である。
本節の最後にその条件も全ての重みで証明する。

\subsection{可換極限での一様な反対称化}

まず $V$ を次元 $n\ge3$ のベクトル空間とし、
$\operatorname{Sym}(V)[w]$ の同じ四対象切出しを考える。
正の $w$ 重みの二cochainは
\[
 \begin{split}
 a:V\otimes\operatorname{Sym}^{m-1}V\to\mathbf C,&\quad
 b:\operatorname{Sym}^{m-1}V\otimes V\to\mathbf C,\\
 c:V\otimes\operatorname{Sym}^{m}V\to V,&\quad
 d:\operatorname{Sym}^{m}V\otimes V\to V
 \end{split}
\]
であり、cocycle条件は
\begin{equation}
 c(v;t z)-d(vt;z)=a(v;t)z-b(t;z)v
 \quad (v,z\in V,\ t\in\operatorname{Sym}^{m-1}V)
 \label{eq:all-weight-positive-cocycle-v40}
\end{equation}
である。

\begin{lemma}[任意の間隔での正の重みの消滅]
\label{lem:all-weight-commutative-vanishing-v40}
$m\ge3$, $n\ge3$ ならば、可換切出しの正の $w$ 重みの複体は
次数一と二で完全である。
\end{lemma}

\begin{proof}
$t=abC$ と書き、$C\in\operatorname{Sym}^{m-3}V$ を固定する。
式\eqref{eq:all-weight-positive-cocycle-v40}は、$v,a,b,z$ の四変数に関して
重み三の場合と全く同じ恒等式となる。
その左の二変数および右の二変数を反対称化する。
$K_{ij,k}=a(i;jkC)-a(j;ikC)$、
$L_{i,kl}=b(ikC;l)-b(ilC;k)$ と書けば
\[
 K_{ij,k}\delta_l^p-K_{ij,l}\delta_k^p
 -L_{j,kl}\delta_i^p+L_{i,kl}\delta_j^p=0.
\]
$a$ の後二変数と $b$ の前二変数の対称性は、さらに
\[
 K_{ij,k}+K_{jk,i}+K_{ki,j}=0,\qquad
 L_{i,jk}+L_{j,ki}+L_{k,ij}=0
\]
を与える。これらの巡回恒等式を用いて
$p=l$ および $p=i$ の縮約から、それぞれ
$(n-1)K_{ij,k}+L_{k,ij}=0$ と
$K_{ij,k}+(n-1)L_{k,ij}=0$ を得る。
従って $n(n-2)K=0$ であり、$K=L=0$ となる。
$C$ は任意なので、$a,b$ は全ての $m$ 個の入力に関して対称である。
二つの一次cochain $\operatorname{Sym}^mV\to\mathbf C$ の境界によりこれらを消す。
残った $c,d$ の共通テンソルは、最初の $m$ 個と最後の $m$ 個の変数に関して対称なので、
全 $m+1$ 変数について対称である。
従って $\operatorname{Sym}^{m+1}V\to V$ の境界で消える。

一次cocycleは隣接射空間上で零であり、それらの積は純粋部分の全非対角成分を生成する。
従って一次cocycle自身が零である。
\end{proof}

\subsection{CGG積へのフィルトレーションによる移行}

$C_F^\bullet\subset C^\bullet(B_{j,r})$ を $w$ フィルトレーションを保つ
正規化cochain複体とし、$Q^\bullet=C^\bullet/C_F^\bullet$ と置く。
補助重みを
\[
 \operatorname{wt}(x,y,z,w)=(0,1,2,m+r+1)
\]
で定める。$q=xz-y^2$ とすると
\[
 N=x\partial_y+2y\partial_z-q^jz^r\partial_w
\]
はこの重みを一つ下げ、$X_0$ は重みを保つ。
従ってCGG積の先頭項は可換積であり、残りの項は補助重みを厳密に下げる。
有限次元複体 $Q^\bullet$ に出力重みから入力重みを引いた有限フィルトレーションを入れると、
そのassociated gradedは補題\ref{lem:all-weight-commutative-vanishing-v40}の複体である。
このため有限フィルトレーションのspectral sequenceから
\begin{equation}
 H^1(Q^\bullet)=H^2(Q^\bullet)=0,
 \qquad \dim H^3(Q^\bullet)=4m^2-6m-7
 \label{eq:all-weight-positive-cohomology-v40}
\end{equation}
を得る。最後の次元は
\[
 \dim Q^1=2n_m+3n_{m+1},\quad
 \dim Q^2=6n_{m-1}+18n_m,\quad
 \dim Q^3=27n_{m-1}
\]
のEuler特性から従う。

各小拡大で、積の最初の非保存項は $Q^2$ のcocycleとなる。
式\eqref{eq:all-weight-positive-cohomology-v40}により、基底変換でそれを消せる。
保存的表示同士を結ぶ同型、保存的部分の埋込みの差、および正規化導分の最初の非保存項は
$Q^1$ のcocycleとなり、従って零である。
これを極大イデアルの次数について帰納すれば、定理の第一、第二の主張を得る。
純粋部分の三つの合成は特殊ファイバーで全射であり、保存的表示ではNakayamaの補題により全射を保つ。
これが合成像の内在的記述と、その平坦性を与える。

\subsection{原始射を動かす相対自己同型が存在しない理由}

制限 $C_F^\bullet\to C^\bullet(N_m)$ の核を $K^\bullet$ とする。
原始射 $w$ を含む三本の非恒等射の可合成列は存在しないので、$K^3=0$ である。
また
\begin{equation}
 \dim K^1=2n_m+3n_{m+1}+11,
 \qquad \dim K^2=6n_{m+1}+18.
 \label{eq:all-weight-relative-dimensions-v40}
\end{equation}
次の単射性が、単なる係数の数え上げを相対変形次元の公式へ変える。

\begin{lemma}[全共鳴での原始射の無限小安定化群]
\label{lem:all-weight-relative-stabilizer-v40}
$H^1(K^\bullet)=0$ である。
\end{lemma}

\begin{proof}
純粋部分上で零の正規化導分を考える。
二つの原始射に対する値を $a_iw_i+p_i$、$p_i\in\operatorname{Sym}^m\langle x,y,z\rangle$ と書く。
$x$ はCGG積で中心的であり、$x\star w_1=w_0\star x=xw$ である。
この等式から $a_0=a_1=a$, $p_0=p_1=p$ を得る。
残る値は左合成 $V\otimes\mathbf Cw_1\to Vw$ により一意に決まる。

三変数のCGG代数を $A$ と書く。
全座標代数は $A$ と $w$ で生成され、その三つの混合関係は
$w\star v-(1+N)^{-d}(v)\star w\in A$、$v=x,y,z$、$d=m+r+1$ である。
これらが全ての混合関係を生成する。
実際、それらで $w$ を右へ移動でき、得られる $A\otimes\mathbf C[w]$ の順序付き単項式は
CGG積の三角性により線形独立である。
従って上の導分は、純粋部分を固定し $w\mapsto aw+p$ とする座標代数の導分へ延びる。

$P=q^jz^r$ と置く。$Nq=0$, $X_0q=2q$, $Nw=-P$ から
\[
 w\star q-q\star w=-(2P+NP)q.
\]
これに導分を施すと
$(2+N)(Np-aP)q=0$ となる。
$N$ は次数 $m$ の純粋部分上で冪零であり、$2+N$ は可逆だから
\[
 Np=aP.
\]
さらに
$w\star y-(y-dx)\star w=-Py$ の保存から
\[
 yNp-xX_0p+dxp=aPy,
 \qquad\text{従って}\qquad X_0p=dp
\]
を得る。

この二条件が $a=p=0$ を強制することを、係数で直接示す。
$X_0p=dp$ および $\deg p=m=2j+r$ により
\[
 p=\sum_{b=0}^{j-1}c_b x^b y^{2j-1-2b}z^{b+r+1}.
\]
$j=0$ ではこの和は空である。
$q^jz^r$ の単項式基底
$x^b y^{2j-2b}z^{b+r}$、$0\le b\le j$ において、$Np=aP$ の係数式は
\begin{equation}
 2(b+r+1)c_b+(2j+1-2b)c_{b-1}
   =a(-1)^{j-b}\binom jb,
 \quad c_{-1}=c_j=0.
 \label{eq:all-weight-extremal-obstruction-v40}
\end{equation}
$\ell_0=1$、
\[
 \ell_{b+1}=-\frac{2(b+r+1)}{2j-1-2b}\ell_b
 \quad (0\le b<j)
\]
とすると、左辺を $\ell_b$ で縮約した和は零である。
右辺の $a$ の係数は
$(-1)^j\sum_{b=0}^j|\ell_b|\binom jb$ という非零有理数になる。
従って $a=0$ であり、式\eqref{eq:all-weight-extremal-obstruction-v40}の三角性から全 $c_b=0$ となる。
これは $r=0$ にもそのまま適用できる。
\end{proof}

\subsection{相対形式空間と自己同型の接続写像}

純粋部分の任意の変形に対し、二つの原始射を含む左右の合成を任意に指定してよい。
三重合成は全て純粋部分に属するので、追加の結合条件はない。
左合成は特殊ファイバーで同型であり、基底変換により標準形へ固定できる。
残る右合成には $3(n_{m+1}+3)$ 個の係数がある。
左合成を固定した相対基底変換には $2(n_m+1)$ 個の係数があり、
補題\ref{lem:all-weight-relative-stabilizer-v40}によりその無限小作用は単射である。
形式的逆関数定理でこの作用に横断するスライスを取れば、
相対変形空間は滑らかで次元
\[
 3(n_{m+1}+3)-2(n_m+1)=k_m
\]
となる。相対自己同型の最初の非零項も $H^1(K^\bullet)$ に属するため、
任意のArtin基底上で恒等に限られる。これで定理の第三の主張を得る。

普通のHochschild二次コホモロジーとの関係は、次の完全列で正確に表される。
\begin{equation}
 \begin{split}
 0\longrightarrow HH^1(B_{j,r})\longrightarrow HH^1(N_m)
 &\xrightarrow{\partial}\mathbf C^{k_m}\\
 &\longrightarrow HH^2(B_{j,r})\longrightarrow HH^2(N_m)\longrightarrow0.
 \end{split}
 \label{eq:all-weight-relative-long-exact-v40}
\end{equation}
ここで $\partial$ は、純粋部分の無限小自己同型を原始射の合成へ持ち上げる障害である。
従って
\[
 \dim HH^2(B_{j,r})-\dim HH^2(N_m)=k_m-\operatorname{rank}\partial.
\]
次に $\partial=0$ を全ての重みで示す。

\subsection{純粋部分の外的導分と通常の変形環}

\begin{lemma}[純粋部分の外的導分は全重みで一次元]
\label{lem:all-weight-pure-HH1-v40}
非零CGG点の純粋部分に対して
\[
 HH^1(N_m)=\mathbf C\,[\log(1+N)]
\]
である。
\end{lemma}

\begin{proof}
まず可換点 $N_{m,0}$ の正規化導分を考える。
左右の次数一射への作用を $F,H\in\operatorname{End}(V)$、中間の
$\operatorname{Sym}^{m-1}V$ への作用を $G$ と書く。
三段合成 $vtz=ztv$ を微分すると
\[
 t\bigl((F-H)(v)z-(F-H)(z)v\bigr)=0
\]
となるので、$F-H$ はスカラーである。
$F$ が誘導する通常の多項式導分と頂点スカラーによる内的導分を引けば、$F=H=0$ にできる。
残る条件は
\[
 xG(yt)=yG(xt)\quad(t\in\operatorname{Sym}^{m-2}V)
\]
であり、$x,y$ は任意の一次式でよい。
$x$ を反転して単項式の因子を一つずつ入れ替えると、$G$ は
$G(x^{m-1})/x^{m-1}$ という有理関数の乗法となる。
この乗法が全ての次数 $m-1$ の多項式を同じ空間へ送るため、その関数は定数である。
この定数も二つの頂点block間の内的導分で消せる。
従って正規化導分は $\mathfrak{gl}_3$ と三次元の内的導分で生成され、
両者の共通部分はEuler導分の一次元である。よって
\[
 HH^1(N_{m,0})=\mathfrak{pgl}_3,\qquad \dim HH^1(N_{m,0})=8.
\]

次に、三本の隣接射 $x,x^{m-1},x$ を可逆にする普遍局所化を考える。
可換点では
\begin{equation}
 N_{m,0}[x,x^{m-1},x]^{-1}\simeq M_4(\mathbf C[Y,Z]),
 \qquad Y=y/x,\quad Z=z/x
 \label{eq:all-weight-pure-localization-v40}
\end{equation}
となる。
実際、これら三本で四頂点を同定すると、中間の $x^{m-1}$ は $1$ となる。
左右の二段合成の関係から、左右の $y,z$ はそれぞれ同じ $Y,Z$ になり、
中間の単項式 $x^ay^bz^c$ は $Y^bZ^c$ になる。
$m\ge3$ なので $yzx^{m-3}$ を両順序で生成する関係が存在し、$YZ=ZY$ を得る。
これで両方向の準同型を構成できる。

CGG係数を $h$ とし、$A=\mathbf C[h]/(h^n)$ 上で同じ局所化を $U_h$ と書く。
中央元 $x$ を反転した三変数座標環の次数零部分を $C_h$ とする。
これは $A$ 上自由であり、$Y,Z$ で生成されるので
$U_h\twoheadrightarrow M_4(C_h)$ が存在する。
法 $h$ では式\eqref{eq:all-weight-pure-localization-v40}に一致する。
核を $K$ とすると、右辺が $A$ 上平坦であることから
\[
 0\longrightarrow K/hK\longrightarrow U_h/hU_h
  \longrightarrow M_4(C_h)/hM_4(C_h)\longrightarrow0
\]
は完全である。従って $K=hK$、さらに $h^n=0$ から $K=0$ となる。
ここでは $U_h$ の平坦性を仮定していない。

可換点の導分が $h^2=0$ まで持ち上がるなら、その局所化も導分になる。
法 $h$ の中心上で誘導する射影的ベクトル場は、
CGG積の半古典括弧
\[
 \pi_{\mathrm{aff}}=2(Z-Y^2)\partial_Y\wedge\partial_Z
\]
を保たなければならない。
これは導分を交換子 $[Y,Z]$ へ適用して $h$ の係数を比較した必要条件である。
射影的ベクトル場を
\[
 \begin{split}
 v(Y)&=b_0+(b_1-a_0)Y+b_2Z-a_1Y^2-a_2YZ,\\
 v(Z)&=c_0+c_1Y+(c_2-a_0)Z-a_1YZ-a_2Z^2
 \end{split}
\]
と書くと、$\mathcal L_v\pi_{\mathrm{aff}}=0$ の係数比較は
\[
 c_0=0,\ c_1=2b_0,\ b_1=c_2=a_0,\ a_1=a_2=b_2=0
\]
を与える。スカラー行列に対応する $a_0$ を除けば、
$\partial_Y+2Y\partial_Z$ の一次元だけが残る。

有限Hochschild複体を $\mathbf C[[h]]$ 上で考え、頂点の三次元の内的導分を先に除く。
この内的導分の像は自由直和因子である。実際、三本の隣接する $x,x^{m-1},x$ の
自己係数を取れば、頂点スカラーの隣接差を読む定数行列が得られ、その階数は三である。
従ってその像で割った一次cochainも自由であり、二次cochainへの微分が誘導される。
その核は自由加群内の飽和部分加群だから、法 $h$ への写像は単射である。
核の元の法 $h$ の像は全て一次で持上げ可能な導分類であるため、
上で計算した一次元の部分空間に属する。
従って一般ファイバーの $HH^1$ の次元は高々一である。
一方、$h^{-1}\log(1+hN)$ は全次数で定義された導分であり、$y\mapsto x$ なので外的に非零である。
従って一般ファイバーの次元は一である。
補助重みによる $y\mapsto h y$, $z\mapsto h^2z$ の変換は、
$\mathbf C((h))$ 上でこの一般ファイバーを $h=1$ のものへ同定する。
これで主張を得る。
\end{proof}

\begin{corollary}[同じ重みの全共鳴は同じ通常変形環を持つ]
\label{cor:all-weight-ordinary-hulls-v40}
全ての $j,r\ge0$, $m=2j+r\ge3$ について
\[
 HH^1(B_{j,r})\xrightarrow{\ \sim\ }HH^1(N_m)=\mathbf C,
 \qquad
 \dim HH^2(B_{j,r})=\dim HH^2(N_m)+k_m.
\]
適合するminiversal環を選べば
\begin{equation}
 R_{B_{j,r}}\simeq R_{N_m}[[u_1,\ldots,u_{k_m}]]
 \label{eq:all-weight-ordinary-hull-factorization-v40}
\end{equation}
である。特に、$m$ を固定すると、$r=0$ の端点を含む全共鳴の通常miniversal環は互いに同型となる。
\end{corollary}

\begin{proof}
純粋部分の生成導分 $\log(1+N)$ は元の四変数CGG代数上にも存在し、
従って各 $B_{j,r}$ へ持ち上がる。
補題\ref{lem:all-weight-pure-HH1-v40}と
式\eqref{eq:all-weight-relative-long-exact-v40}により $\partial=0$ となる。
純粋部分のminiversal族へ定理の滑らかな相対形式空間を付加すると、
全cornerの形式的にversalな族を得る。その接写像は上の次元公式により同型なのでminiversalであり、
式\eqref{eq:all-weight-ordinary-hull-factorization-v40}を得る。
\end{proof}

この同型は通常の変形環についてのものであり、外的導分の作用を含む対称性保持部分を同一視しない。
重み三の二つの共鳴で、通常hullは同じでも対称性保持部分の被約性が異なることは既に見た。
一般の重みの $R_{N_m}$ 自体の決定、および対称性保持曲線の切詰め次数の決定は、引き続き未決である。

付属検証器は $3\le m\le7$ の全共鳴に対して $Q^\bullet,K^\bullet$ を
CGG積から独立に構成し、有理数上で複体条件と上記の階数を確認する。
また式\eqref{eq:all-weight-extremal-obstruction-v40}の消去機構と、
射影的Poisson安定化群を定める係数行列の階数七を直接照合する。
全整数についての定理は、これらの有限標本ではなく、上の反対称化・有限フィルトレーション・係数漸化式による。

\section{純粋三変数部分と原始射への全形式的還元}
\label{sec:vk-filtered-reduction-v39}

対称性保持部分の六番目の接方向は、五次元族の原始射の合成だけでは
捉えられない。本節では、その理由を全変形問題の構造から説明する。
原始射の右合成には $54$ 個の係数があるが、Hochschild二次コホモロジーに
現れる方向は $32$ 個である。残りの三方向は純粋な $x,y,z$ の部分の変形であり、
この分離は接空間に限らず、全ての形式変形と導分の持上げに適用できる。

\subsection{フィルトレーションを破る変形は基底変換で除去できる}

$B=B_{j,r}$、$(j,r)=(0,3),(1,1)$ とする。
本節の基点は量子化パラメータ $s=1$ の二つの固定された非零量子点である。
以下の全形式近傍に関する主張は、この二点での $35$ 次元変形hullを対象とする。
四対象 $(0,1,3,4)$ の各射空間に、$w$ の個数による二段の
フィルトレーション $F_0\subset F_1$ を入れる。
$F_0$ は純粋な $x,y,z$ の部分であり、その四対象代数を $N$ と書く。
この $N$ は二つの場合で同じ代数である。
頂点冪等元を固定した正規化Hochschild複体を $C^\bullet(B)$、
フィルトレーションを保つcochainの部分複体を $C_F^\bullet(B)$ と書く。
単項式基底では、後者の条件は
\[
 \deg_w(\text{出力})\le
 \sum\deg_w(\text{入力})
\]
である。積とGerstenhaber括弧はこの条件を保つ。

\begin{theorem}[全形式変形のフィルトレーションによる還元]
\label{thm:vk-filtered-reduction-v39}
特殊ファイバーの同定と頂点冪等元を固定した変形を考える。
\begin{enumerate}
\item 任意の形式変形は、特殊ファイバー上で恒等である基底変換によって
$w$ のフィルトレーションを保つ変形にできる。
フィルトレーションを保つ二つの変形の間の、この種の代数同型も
フィルトレーションを保つ。
\item フィルトレーションを保つ変形上の、頂点冪等元を零にする正規化導分は、
フィルトレーションを保つ。
任意の外的導分類は内的導分を差し引いて正規化できるので、
基点の外的導分 $\delta$ の持上げを含めた変形問題にも同じ還元が適用できる。
\item $N$ の変形hullは滑らかな三次元形式空間である。
$B$ から $N$ への制限の相対変形空間は滑らかで、次元は $32$ である。
従ってhullを適切に選べば
\[
 R_N\simeq\mathbf C[[a_1,a_2,a_3]],\qquad
 R_B\simeq R_N[[z_1,\ldots,z_{32}]]
\]
となる。この積分解は座標の選択に依存する。
\end{enumerate}
\end{theorem}

\begin{proof}
商複体 $Q^\bullet=C^\bullet(B)/C_F^\bullet(B)$ の次数 $1,2,3$ における
次元と微分の階数は、それぞれ
\begin{equation}
 (\dim Q^1,\dim Q^2,\dim Q^3)=(65,216,162),\qquad
 (\operatorname{rank}d_Q^1,\operatorname{rank}d_Q^2)=(65,151).
 \label{eq:vk-filtered-quotient-ranks-v39}
\end{equation}
特に $H^1(Q^\bullet)=H^2(Q^\bullet)=0$ である。
この計算は二つの $B$ で同じであり、付属コードは有理数上の行基本変形により
部分複体条件、$d_Q^2d_Q^1=0$、および表示した二つの階数を直接検証する。

変形が基底の極大イデアルの $n$ 乗を法としてフィルトレーションを保つとする。
次の次数における非保存部分は、結合則により $Q^2$ のcocycleである。
$H^2(Q^\bullet)=0$ により、それは $d_Q^1$ の像であり、同じ次数の
基底変換で除去できる。これを繰り返せば第一の存在を得る。
二つの保存変形を結ぶ基底変換について、その最初の非保存項は
$\ker d_Q^1$ に属するが、式\eqref{eq:vk-filtered-quotient-ranks-v39}から零である。
同じ議論を次数ごとに適用すれば、基底変換自体が保存的である。
正規化導分の最初の非保存項についても $d_Q^1$ の核に属することから第二の主張を得る。
一般の導分 $D$ は、$a=\sum_iD(e_i)e_i$ に対する内的導分 $[a,-]$ を差し引けば
全ての $e_i$ を零にする。
いずれの議論もArtin局所基底の小拡大に適用でき、その逆極限で形式基底にも適用できる。

第三の主張のために、制限
$C_F^\bullet(B)\to C^\bullet(N)$ の核を $K^\bullet$ と書く。
必要な計算は
\[
\begin{array}{c|c|c}
 \text{複体}&(\dim C^1,\dim C^2,\dim C^3)
   &(\operatorname{rank}d^1,\operatorname{rank}d^2)\\\hline
 C_F^\bullet(B)&(555,1368,810)&(551,782)\\
 C^\bullet(N)&(479,1260,810)&(475,782)\\
 K^\bullet&(76,108,0)&(76,0)
\end{array}
\]
である。従って
\[
 HH^*(N)=(1,1,3,28),\qquad H^1(K^\bullet)=0,\qquad
 H^2(K^\bullet)\simeq\mathbf C^{32}.
\]

相対問題には高次の結合条件が存在しない。
実際、三本の非恒等射の可合成列は $0\to1\to3\to4$ だけであり、
その全ての射と中間積は $N$ に属する。
従って $N$ の任意の変形に対して、原始射 $w_0,w_1$ を含む左右二つの合成を
任意に指定すれば、$B$ の保存的変形が得られる。
相対基底変換の微分は単射であるため、相対商は形式的に滑らかで次元 $108-76=32$ である。

既に証明した $B$ の全変形の無障害性から、$N$ の無障害性も従う。
実際、$N$ のArtin基底上の変形を上記の方法で $B$ へ拡大し、
$B$ を基底の小拡大へ持ち上げ、第一の主張を用いて保存的に直してから制限すればよい。
$HH^2(N)$ は三次元なので $R_N$ は滑らかな三変数環である。
相対的に滑らかな $32$ 変数を付け加えることで、表示したhullを得る。
\end{proof}

\begin{corollary}[合成からの純粋部分の復元]
\label{cor:vk-intrinsic-pure-part-v39}
任意の標識付き形式変形で、二つの隣接二段合成と三段合成の像の階数は
それぞれ $10,10,15$ に保たれ、いずれも射空間の平坦な直和因子となる。
これらを $N_{03},N_{14},N_{04}$ とし、残りの射空間をそのまま用いると、
純粋部分 $N$ の変形を合成だけから復元できる。
従って $B$ から $N$ への制限は、表示の選択によらない変形関手の射であり、
定理の相対的に滑らかな $32$ 次元の射を与える。
\end{corollary}

\begin{proof}
基点ではこれらの合成は純粋部分へ全射であり、階数は $10,10,15$ である。
保存的な変形では同じ純粋部分への全射性がNakayamaの補題で保たれる。
従って合成の像は、それぞれ自由な直和因子 $N_{03},N_{14},N_{04}$ そのものである。
定理\ref{thm:vk-filtered-reduction-v39}により全ての変形がそのように表示できるので、
合成の像自体が表示によらない平坦な直和因子を与える。
\end{proof}

\subsection{原始射の合成が占める三十二方向}

左合成 $V\otimes\mathbf Cw_1\to\operatorname{Hom}(0,4)/N_{04}$ は
基点で同型なので、形式近傍全体で基底変換により固定できる。
残る右合成の空間は
\[
 W=\operatorname{Hom}_{\mathbf C}
   (\mathbf Cw_0\otimes V,\operatorname{Hom}(0,4)),\qquad \dim W=54
\]
である。左合成を固定した後に残る基底変換は
\[
 w_0\longmapsto c_0w_0+p_0,\qquad
 w_1\longmapsto c_1w_1+p_1,
 \qquad p_i\in N_3,\quad c_i\in\mathbf C^\times
\]
で、形式近傍で $2+10+10=22$ 次元である。
その無限小作用は単射なので、$W$ のHochschild像 $V_{j,r}$ は $32$ 次元となる。
純粋四次式に出力を限定した $45$ 次元部分 $W_0$ の像は $24$ 次元である。

\begin{proposition}[純粋部分への同変な完全列]
\label{prop:vk-pure-exact-sequence-v39}
自然な制限により、$L_\delta$ と可換な完全列
\begin{equation}
 0\longrightarrow V_{j,r}\longrightarrow HH^2(B_{j,r})
   \xrightarrow{\rho_N}HH^2(N)\longrightarrow0
 \label{eq:vk-pure-exact-sequence-v39}
\end{equation}
を得る。それぞれの $L_\delta$ のJordan分割は次の通りである。
\[
\begin{array}{c|c|c|c}
 (j,r)&V_{j,r}&HH^2(B_{j,r})&HH^2(N)\\\hline
 (0,3)&(12,9,7,4)&(13,10,7,5)&(3)\\
 (1,1)&(11,8,6,4,3)&(11,9,6,5,3,1)&(3)
\end{array}
\]
さらに
\[
 \operatorname{rank}\bigl(\ker L_\delta|_{HH^2(B_{j,r})}
          \longrightarrow\ker L_\delta|_{HH^2(N)}\bigr)=j.
\]
\end{proposition}

\begin{proof}
定理の商複体の消滅により $HH^2(B)$ を保存的複体で計算できる。
$K^\bullet$ の長完全列と $H^3(K^\bullet)=0$ から
式\eqref{eq:vk-pure-exact-sequence-v39}を得る。
この完全列は導分の作用と可換である。
付属コードは $W$ の $54$ 本の単位cochainの像、純粋複体の三本のcocycle、
および制限写像の行列を有理数上で構成し、完全性と同変性を直接確認する。
表はその三つの作用の冪の階数から得られ、最後の階数も制限行列で検証する。
\end{proof}

\subsection{純粋部分では対称性保持曲線が全次数に存在する}

純粋部分については、対称性保持方向は一変数の平坦族へ積分できる。
$D=1-\gamma$ と置き、$\mathbf C[[\gamma]]$ 上の代数 $C_\gamma$ を
\begin{align}
 yx&=\frac{1+\gamma}{D}xy-\frac1D x^2-\frac{2\gamma}{D}xZ,
 \notag\\
 Zx&=\frac{1+\gamma}{D}xZ-\frac2Dxy+\frac1D x^2,
 \label{eq:vk-pure-gamma-relations-v39}\\
 Zy&=yZ+\frac1Dxy-\frac1D xZ
 \notag
\end{align}
で定める。
特殊ファイバーは $\alpha=-1$ の純粋CGG代数において $Z=z+y$ としたものである。
既出の有限対象のZhang同型を特殊ファイバーの同定として用いれば、
その $(0,1,3,4)$ 切出しは $N$ の形式変形を与える。
$HH^1(N)$ は一次元であるため、前節までと同じ相対導分複体の単元消去により、
外的導分の持上げ条件は完備hull上の残余列 $a_N:R_N\to R_N^3$ の消滅で表される。
この列の線形項は $-L_\delta$ である。
従って $Z_N$ はNoether完備な閉形式部分であり、その接空間は $\ker L_\delta$ となる。

\begin{proposition}[純粋対称性保持部分の滑らかさ]
\label{prop:vk-pure-symmetry-curve-v39}
$N$ の外的導分の持上げの存在を表す形式部分は
\[
 Z_N\simeq\operatorname{Spf}\mathbf C[[\gamma]]
\]
である。上記の $C_\gamma$ の切出しがこの曲線を実現し、保持される導分は
\begin{equation}
 D_\gamma(x)=0,\qquad D_\gamma(y)=x+2\gamma Z,\qquad
 D_\gamma(Z)=2y
 \label{eq:vk-pure-gamma-derivation-v39}
\end{equation}
である。
\end{proposition}

\begin{proof}
規則 $yx,Zx,Zy$ を順序付き単項式へ書き換えると、唯一の重なりは $Zyx$ である。
その二つの簡約は $\mathbf Q(\gamma)$ 上で一致し、分母は $1-\gamma$ の冪だけである。
従って $x^ay^bZ^c$ は $\mathbf C[[\gamma]]$ 上のPBW基底となる。
式\eqref{eq:vk-pure-gamma-derivation-v39}を三つの二次関係へ適用すると、
同じ簡約によって全て零になる。これら四個の有理式恒等式は付属コードで直接検証する。
従って積と導分は全次数で定義される。

特殊ファイバーの導分は $\delta|_N$ である。
この族のKodaira--Spencer類は $HH^2(N)$ で非零であり、
$\ker L_\delta$ の一次元空間を張る。
一方、命題\ref{prop:vk-pure-exact-sequence-v39}から $Z_N$ の接空間も一次元である。
従って、この族への完備局所環の写像は接空間で同型である。
一変数形式冪級数環による表示と合成すれば一次係数が非零の自己準同型になるため、
表示の核は零である。これで $Z_N$ 全体の主張を得る。
\end{proof}

式\eqref{eq:vk-pure-exact-sequence-v39}に核・余核の列を適用すると、
二つのVK代数の相違が明確になる。
$B_{0,3}$ では純粋曲線の接方向は、原始射の合成を接合する時点で
一次の導分持上げを妨げる。
$B_{1,1}$ ではその接方向が持ち上がるが、既出の非零二次障害が残る。
この二次障害の $HH^2(N)/\operatorname{im}L_\delta$ への像は零であり、
$V_{1,1}/\operatorname{im}L_\delta$ に一意な非零の逆像を持つ。
従って二次障害は純粋三変数部分の量子化から生じるものではなく、
原始射の合成と導分を両立させる相対問題で生じる。
この同定も、既出の二次障害の代表と本節の制限行列から有理数上で直接検証する。

\section{外的導分の跳躍部分は滑らかな五次元部分の二重構造である}
\label{sec:vk-double-symmetry-v39}

定理\ref{thm:vk-Ga-families-v38}では $B_{1,1}$ の対称性保持部分に
滑らかな五次元の閉形式部分を構成し、
命題\ref{prop:vk-jump-finite-normal-form-v38}では全体の高次構造を
$A[[t]]/(t(t-b),tJ)$ の形に限定した。
本節では残った $b,J$ を決定する。
結論は $b=0$, $J=0$ であり、対称性保持部分は全体で非被約な二重構造となる。
証明は有限jetの安定からの推論ではなく、五つの自由パラメータと一つの
平方零パラメータを持つ全次数の代数族を構成することによる。

\begin{theorem}[対称性保持部分の全完備局所環]
\label{thm:vk-double-jump-v39}
$B=B_{1,1}$、$R_B\simeq\mathbf C[[t_1,\ldots,t_{35}]]$ とし、
$I_B=\operatorname{Fitt}_{34}HH^2_{R_B}(B_{R_B})$ と置く。
次が成り立つ。
\begin{enumerate}
\item 対称性保持部分の全環は
\begin{equation}
 R_B/I_B\simeq
 \mathbf C[[u_1,\ldots,u_5,\varepsilon]]/(\varepsilon^2).
 \label{eq:vk-double-jump-ring-v39}
\end{equation}
定理\ref{thm:vk-Ga-families-v38}の五次元の族は、その被約部分 $\varepsilon=0$ を与える。
従って上記の正規形における未知量は $b=J=0$ である。
\item $R_B$ に形式座標を取ることで、埋込みを含めて
\begin{equation}
 R_B=\mathbf C[[u_1,\ldots,u_5,\varepsilon,z_1,\ldots,z_{29}]],
 \qquad I_B=(z_1,\ldots,z_{29},\varepsilon^2)
 \label{eq:vk-double-embedded-ideal-v39}
\end{equation}
とできる。
特に $Z_B$ は次元五、埋込次元六、重複度二の局所完全交叉であり、
GorensteinかつCohen--Macaulayである。
その被約部分は滑らかで、埋込成分は存在しない。
\item 式\eqref{eq:vk-double-jump-ring-v39}上の族は、
非零外的導分を持ち上げる導分 $\mathcal D$ で
$\mathcal D^{24}=0$ を満たすものを持つ。
従って二重構造全体上で $\exp(s\mathcal D)$ という
$\mathbf G_a$ 作用が定義される。
\end{enumerate}
\end{theorem}

ここで $Z_B$ は導分の存在条件を表す部分であって、導分そのものを選んだ
変形関手の基底ではない。第三の主張は、この存在条件の全てを
実際の冪零導分を持つ一つの族で同時に実現できることを述べる。

\subsection{平方零方向を含む全次数の族}

$A_0=\mathbf Q[u_1,\ldots,u_5]$ とする。
定理\ref{thm:vk-Ga-families-v38}で構成した五変数族の積と導分を
$\mu_{\boldsymbol u},D_{\boldsymbol u}$ と書く。
積は高々二次、導分は一次であり、$D_{\boldsymbol u}^{12}=0$ である。
さらに、定理\ref{thm:vk-symmetry-second-neighborhood-v37}の計算で固定した
有理基底における第六の不変cocycleとその導分補正を
$U_6,E_6$ と書く。従って
\[
 d^2U_6=0,\qquad d^1E_6=L_\delta U_6,
 \qquad [U_6]\notin\operatorname{im}L_\delta.
\]

\begin{lemma}[有理係数による平方零横断族]
\label{lem:vk-rational-double-family-v39}
定数項一の多項式 $\Delta\in A_0$ と、多項式係数の相対cochain
\[
 P\in C_E^2(B)\otimes A_0,
 \qquad Q\in C_E^1(B)\otimes A_0
\]
を、$P(0)=Q(0)=0$ かつ次の条件を満たすように選べる。
\[
 \Phi=\Delta U_6+P,\qquad
 \Psi=\Delta E_6+Q
\]
と置くと、全多項式環 $A_0$ 上で
\begin{align}
 \mu_{\boldsymbol u}\circ\Phi+\Phi\circ\mu_{\boldsymbol u}&=0,
 \label{eq:double-associativity-identity-v39}\\
 [D_{\boldsymbol u},\Phi]+[\Psi,\mu_{\boldsymbol u}]&=0
 \label{eq:double-derivation-identity-v39}
\end{align}
が成り立つ。
$P$ は原始射 $w_0$ と次数一射の積だけを変え、二cochainの二十二座標に支持を持つ。
$Q$ の支持は一次cochainの百座標である。
$\Delta$ の全次数は十二、$P,Q$ の全次数は高々十四である。

従って
\begin{equation}
 \mathcal A=A_0[\Delta^{-1},\varepsilon]/(\varepsilon^2),
 \quad
 \widetilde\mu=\mu_{\boldsymbol u}+\varepsilon\Delta^{-1}\Phi,
 \quad
 \mathcal D=D_{\boldsymbol u}+\varepsilon\Delta^{-1}\Psi
 \label{eq:vk-rational-double-family-v39}
\end{equation}
は、階数五十六の平坦な代数族と、その導分を定める。
原点での六つのKodaira--Spencer類は独立であり、$\ker L_\delta$ の基底となる。
\end{lemma}

\begin{proof}
cochain基底は、対象 $(0,1,3,4)$、各次数の重み付き単項式基底、
その辞書式順序によって固定する。
五変数族と $U_6,E_6$ は、上記二定理の有理係数の証明書から取る。
未知補正を上記の二十二座標と百座標へ制限すると、
式\eqref{eq:double-associativity-identity-v39}と
\eqref{eq:double-derivation-identity-v39}は、百二十二個の未知関数についての
有限な線形方程式となる。
その原点の係数行列は列階数百二十二である。
従って独立な百二十二行を選び、原点で単元となる行列式を逆にして解くことができる。
得られた解の分母を払い、定数項が一となるよう規格化したものが
$\Delta,P,Q$ である。

選択した行だけでなく、全ての残余方程式を検証する必要がある。
ここでは、構成に用いた線形消去から独立に、元のCGG積と全cochainを再構成して
式\eqref{eq:double-associativity-identity-v39}の九百七十二成分と、
式\eqref{eq:double-derivation-identity-v39}の千五百八十四成分を評価する。
分母を除いた各成分は $\mathbf Q[u_1,\ldots,u_5]$ の多項式であり、
全二千五百五十六成分が恒等的に零となる。
同じ直接計算で $\mu_{\boldsymbol u}$ 自身の結合則と
$D_{\boldsymbol u}$ の導分条件も確認する。
全有理係数を固定した証明書と読み取り専用検証コードに保存する。
検証は多項式全体について行い、有限jetへの打切りや数値代入を用いない。

検証データの大きさを明記すると、$\Delta$ は四百五十五個の単項式を持ち、
百二十二個の補正分子の単項式数の合計は十三万五千百十五である。
この展開の大きさは選んだcochain基底と消去に由来し、
結論の局所環がその数のパラメータを持つという意味ではない。
検証の本体は、固定した多項式を二つの表示式へ代入する有限な恒等式計算である。

$\varepsilon^2=0$ であるため、
式\eqref{eq:double-associativity-identity-v39}は $\widetilde\mu$ の結合則の
全ての新しい項を、式\eqref{eq:double-derivation-identity-v39}は
$\mathcal D$ の導分条件の全ての新しい項を与える。
元の単位と頂点冪等元は固定され、基礎加群は自由な五十六次元加群なので平坦である。

原点での六つの初速度は $U_1,\ldots,U_6$ である。
境界写像の階数は $616$、これら六本を併合した行列の階数は $622$ である。
従ってKodaira--Spencer類は六本とも独立である。
導分の一次持上げから各類は $\ker L_\delta$ に入り、その次元は六なので基底となる。
\end{proof}

\subsection{全局所環を決める形式的な判定}

\begin{lemma}[平方零族と非零平方関係による判定]
\label{lem:double-family-determines-hull-v39}
完備局所 $\mathbf C$ 代数 $R$ の埋込次元を $k+1$ とする。
その最小表示の定義イデアルは極大イデアルの二乗に含まれ、
非零な二次関係を持つとする。
さらに
\[
 R\longrightarrow
 T=\mathbf C[[u_1,\ldots,u_k,\varepsilon]]/(\varepsilon^2)
\]
という、余接空間上で同型となる局所準同型が存在するならば、
$R\simeq T$ である。
\end{lemma}

\begin{proof}
$S=\mathbf C[[x_1,\ldots,x_{k+1}]]$ からの最小表示 $S\twoheadrightarrow R$ を取る。
合成 $S\to T$ の生成元の像を $\mathbf C[[u_1,\ldots,u_k,\varepsilon]]$ へ持ち上げると、
その一次項は可逆である。形式的逆関数定理により、座標変更の後に
$S\to T$ を標準の剰余写像とすることができる。
従って $I=\ker(S\to R)$ は $(\varepsilon^2)$ に含まれる。
非零な二次項を持つ $f\in I$ を取ると、
$f=\varepsilon^2h$ で $h(0)\ne0$ である。
$h$ は単元なので $(\varepsilon^2)\subset I$ となり、両イデアルは一致する。
\end{proof}

\begin{proof}[定理\ref{thm:vk-double-jump-v39}の証明]
補題\ref{lem:vk-rational-double-family-v39}の族を原点で完備化する。
$\Delta(0)=1$ なので、その基底は
$T=\mathbf C[[u_1,\ldots,u_5,\varepsilon]]/(\varepsilon^2)$ である。
保持する導分は基点で $\delta$ に一致する。
次数一の $y$ から $x$ への係数は $1$ を法 $\varepsilon$ で持つため、単元である。
一方、頂点冪等元を固定する内的導分は各矢空間上でスカラーとして作用する。
従って、この族は $Z_B$ へ因子化する。
六つのKodaira--Spencer類が独立であるため、
$R_B/I_B\to T$ は余接空間上で同型となる。

定理\ref{thm:vk-symmetry-second-neighborhood-v37}により、
$R_B/I_B$ の埋込次元は六であり、その二次関係に非零な平方がある。
具体的には、選んだ $U_6$ について二次障害の第二成分が
\[
 \frac{432000}{2468041}\ne0
\]
である。
補題\ref{lem:double-family-determines-hull-v39}を適用して
式\eqref{eq:vk-double-jump-ring-v39}を得る。
平方零方向を零にした族は定理\ref{thm:vk-Ga-families-v38}の五変数族そのものなので、その像は被約部分に一致する。

残余列 $a$ の線形項は階数二十九である。
その独立な二十九成分を形式座標 $z_1,\ldots,z_{29}$ に取り、
それらを消した六変数環へ上の同型を適用すれば、
式\eqref{eq:vk-double-embedded-ideal-v39}を得る。
$z_1,\ldots,z_{29},\varepsilon^2$ は正則局所環の正則列である。
二重構造のCohen--Macaulay性、Gorenstein性、重複度、および埋込成分の不存在は、
この表示から従う。

最後に、$D=D_{\boldsymbol u}$ と $E=\Delta^{-1}\Psi$ と書けば
$\mathcal D=D+\varepsilon E$、$D^{12}=0$ である。
$\mathcal D^{24}$ の展開では、$\varepsilon$ を二個以上含む項は零である。
一個含む項は $D^a\varepsilon E D^b$、$a+b=23$ という形であり、
$a,b$ のどちらかが十二以上だから零である。
$\varepsilon$ を含まない項も零なので $\mathcal D^{24}=0$ となる。
従って指数写像は有限和で定義され、全基底上の $\mathbf G_a$ 作用を与える。
\end{proof}

\subsection{相対Hochschild複体と任意の形式弧の完全な次数公式}

\begin{corollary}[残余列の埋込正規形]
\label{cor:vk-relative-column-normal-form-v39}
式\eqref{eq:vk-double-embedded-ideal-v39}の形式座標と、
相対最小複体の次数二の自由基底を選ぶことで、残余列を
\begin{equation}
 a=(z_1,\ldots,z_{29},\varepsilon^2,0,\ldots,0)^t\in R_B^{35}
 \label{eq:vk-double-column-v39}
\end{equation}
とできる。
従って相対Hochschild複体は
\[
 R_B[0]\oplus
 \left(R_B\xrightarrow{(z_1,\ldots,z_{29},\varepsilon^2)^t}
                    R_B^{30}\right)_{1,2}
 \oplus R_B^5[-2]\oplus R_B^{39}[-3]
\]
にホモトピー同値である。
また
\[
 \operatorname{Ext}_{R_B}^1
       \bigl(HH^2_{R_B}(B_{R_B}),R_B\bigr)
       \simeq R_B/I_B
\]
であり、二重構造は相対 $HH^2$ 加群のこの双対からも回復できる。
\end{corollary}

\begin{proof}
列の独立な一次項を持つ二十九成分を $z_i$ に選ぶ。
他の成分から $z_i$ の倍を引く行基本変形により、残りは
$\mathbf C[[u_1,\ldots,u_5,\varepsilon]]$ に値を取る。
それらが生成するイデアルは $(\varepsilon^2)$ なので、少なくとも一つは
$\varepsilon^2$ の単元倍であり、他の五つはその倍数である。
単元による行の規格化と消去によって式\eqref{eq:vk-double-column-v39}を得る。
次数零と三の直和因子は既知の最小複体のままである。
列は $R_B$ 上で単射であるから、$HH^2$ の長さ一の自由分解を双対にすると、
第一Extはこの列の成分が生成するイデアルによる剰余環となる。
\end{proof}

\begin{corollary}[全次数の接触次数]
\label{cor:vk-contact-order-exact-v39}
任意の形式弧 $c:\operatorname{Spf}\mathbf C[[s]]\to\operatorname{Spf}R_B$ に対し、
$\operatorname{ord}_s0=\infty$ と約束すると
\begin{equation}
 \operatorname{ord}_s(c^*I_B)
  =\min\left\{
      \operatorname{ord}_s z_1(c),\ldots,
      \operatorname{ord}_s z_{29}(c),
      2\operatorname{ord}_s\varepsilon(c)
          \right\}.
 \label{eq:vk-contact-order-formula-v39}
\end{equation}
この値が有限な $\nu$ ならば、$\mathcal O=\mathbf C[[s]]$ 上で
\[
 HH^1_{\mathcal O}(B_c)=0,
 \qquad HH^2_{\mathcal O}(B_c)
       \simeq\mathcal O^{34}\oplus\mathcal O/(s^\nu).
\]
値が無限大ならば、弧は被約部分 $Y$ を通り、
$HH^1_{\mathcal O}\simeq\mathcal O$、
$HH^2_{\mathcal O}\simeq\mathcal O^{35}$ となる。
\end{corollary}

\begin{proof}
式\eqref{eq:vk-double-embedded-ideal-v39}をDVRへ引き戻すと、
生成元の付値の最小値がイデアルの付値である。
二重方向の生成元は $\varepsilon^2$ なので、その付値は二倍になる。
相対Hochschild群は列のDVR上の標準形から従う。
値が無限大の場合は $z_i(c)=0$ および $\varepsilon(c)^2=0$ である。
$\mathbf C[[s]]$ は被約なので $\varepsilon(c)=0$ となり、弧は $Y$ を通る。
\end{proof}

この公式は、初速度別の次数表を全高次の座標により統一する。
不変接空間の外ではある $z_i$ に一次項があり、次数は一である。
不変接空間の中で五次元超平面の外では $\varepsilon$ に一次項があり、次数は二となる。
五次元超平面内では二次以降の $z_i$ と $\varepsilon$ の選択によって、
既に構成した任意の有限次数または無限大が生じる。
今回はそれらの次数だけでなく、それを支配する全非被約構造が決定された。

\subsection{純粋三変数部分の曲線に対する相対的な意味}

\begin{corollary}[純粋対称性保持曲線の一次近傍]
\label{cor:vk-double-over-pure-curve-v39}
命題\ref{prop:vk-pure-symmetry-curve-v39}の純粋部分の対称性保持曲線を
$Z_N=\operatorname{Spf}\mathbf C[[\gamma]]$ とする。
$B_{1,1}$ から純粋部分への制限に関して、
式\eqref{eq:vk-double-jump-ring-v39}の同型は
\[
 R_B/I_B\simeq
 \mathbf C[[u_1,\ldots,u_5,\gamma]]/(\gamma^2)
\]
という $\mathbf C[[\gamma]]$ 代数の同型として取ることができる。
一方、$B_{0,3}$ の滑らかな四次元の対称性保持部分は、
$Z_N$ の原点 $\gamma=0$ へ写る。
\end{corollary}

\begin{proof}
原始射に支持を持つ既存の五変数族は純粋部分の積を全く変えないため、
$Z_B$ の被約部分 $Y$ は $Z_N$ の原点へ写る。
従って、$\gamma$ の引戻しは、既に求めた環の中で
$\varepsilon h(u)$ という形を持つ。
命題\ref{prop:vk-pure-exact-sequence-v39}により、
$\ker L_B\to\ker L_N$ は $B_{1,1}$ で全射である。
従って $h(0)\ne0$ であり、$\varepsilon h(u)$ 自体を横断座標に取ればよい。
$B_{0,3}$ では、純粋部分を固定する四変数族が対称性保持部分全体を実現するので、
制限は全次数で定数である。
\end{proof}

したがって、純粋な三変数部分では全次数に存在する対称性保持方向が、
原始射の合成を付加すると、第一の共鳴では原点に、第二の共鳴では
原点の平方零近傍に制限される。
この記述は対称性保持部分の形式基底に関するものであり、
五変数族の全てが標準六項列へ補完できるという主張ではない。

\begin{corollary}[標準六項補完像は横断長さ二の切片である]
\label{cor:vk-completion-transverse-slice-v39}
標準六項補完に対応するhull間の射を一つ選び、
$C_{\mathrm{std}}\subset\operatorname{Spf}R_B$ をその補完像、
$Y=(Z_B)_{\mathrm{red}}$ とする。
純粋対称性保持曲線の座標 $\gamma$ を保ったまま形式座標を選ぶことで、
三つの閉形式部分のイデアルを同時に
\begin{equation}
 \begin{aligned}
 I_{Z_B}&=(z_1,\ldots,z_{29},\gamma^2),\\
 I_Y&=(z_1,\ldots,z_{29},\gamma),\\
 I_{C_{\mathrm{std}}}&=(z_1,\ldots,z_{29},u_1,\ldots,u_5,\gamma^2)
 \end{aligned}
 \label{eq:vk-three-embedded-loci-v39}
\end{equation}
と書ける。
従って標準補完像は、五つの滑らかな方向を零にした、
二重構造の横断長さ二の切片そのものである。
標準六項のhullを $R_{\Lambda}$ と書けば、純粋座標 $\gamma$ を同定した上で
$R_{Z_B}\simeq R_{\Lambda}[[u_1,\ldots,u_5]]$ となり、
標準補完像はこの相対的に滑らかな五変数形式空間の零切断である。
\end{corollary}

\begin{proof}
標準六項のhullは長さ二の $\mathbf C[[\eta]]/(\eta^2)$ であり、
そのcornerへの接写像は非零である。
従って標準補完の形式基底から $\operatorname{Spf}R_B$ への射は閉埋込みとなり、
補完像は長さ二である。
定理\ref{thm:vk-completion-secondary-v38}により、その唯一の非零接方向には
導分の一次持上げが存在するので、像は $Z_B$ に入る。

純粋部分への接写像も非零であることを確認する。
命題\ref{prop:vk-pure-exact-sequence-v39}の不変接空間上の制限は階数一であり、
その核は、純粋部分を固定する五次元の $T_0Y$ である。
一方、定理\ref{thm:vk-completion-secondary-v38}の補完接線は $T_0Y$ に含まれない。
従って補完接線は純粋曲線の接線へ同型に写る。
有理cochainの基底でも、既存の補完接線の $35$ 座標へ
純粋部分への $3\times35$ 制限行列を掛けると
\[
 \left(\frac{229}{120},\frac{117}{20},-\frac{1571}{60}\right)^t
\]
となり、純粋族の $\gamma$ 方向の座標に一致する。
% Exact Fraction verification: code/verify_vk_completion_image_v39.py.

従って、純粋曲線から補完像への準同型は
$\mathbf C[[\gamma]]\twoheadrightarrow\mathbf C[[\eta]]/(\eta^2)$ であり、
$\gamma$ の像は $\eta$ の非零定数倍である。
この像を補完像の座標に取り直すと、
$R_{Z_B}=\mathbf C[[u_1,\ldots,u_5,\gamma]]/(\gamma^2)$ からの写像は
$\gamma\mapsto\gamma$、$u_i\mapsto a_i\gamma$ という形になる。
$u_i'=u_i-a_i\gamma$ とすれば、被約部分を保ったままこの像を $u_i'=0$ にできる。
この座標変更を $R_B$ へ持ち上げれば、式\eqref{eq:vk-three-embedded-loci-v39}を得る。
\end{proof}

\section{標準補完方向における外的導分の二次消失}
\label{sec:vk-completion-symmetry-v38}

本節では、$m=3$, $j=r=1$ の量子VK代数について、標準六項列へ補完できる
唯一の一次方向を、四項圏に内在する対称性保持条件の中で位置付ける。
結果は、補完できる一次方向こそが、二次では外的導分を保持できない方向であるというものである。
さらに、その方向を積分する任意の形式弧に沿って、相対Hochschild群の捩れの長さまで決まる。

$B=B_{1,1}$、$\Lambda=A_{1,1}$ とし、
\[
 \ell=\operatorname{im}\bigl(HH^2(\Lambda)\longrightarrow HH^2(B)\bigr),
 \qquad L=[\delta,-],\qquad
 K=\ker L,\quad Q=\operatorname{coker}L
\]
と置く。ここで $\delta=\log(1+N)$ であり、計算では $\alpha=0$, $s=1$ を用いる。
$\ell$ は一次元、$K,Q$ は六次元である。

\begin{theorem}[補完接線と内在的二次障害]
\label{thm:vk-completion-secondary-v38}
次が成立する。
\begin{enumerate}
\item $\ell\subset K$ であり、自然写像 $K\to Q$ の $\ell$ への制限は単射である。
 従って
 \[
  \ell\cap\operatorname{im}L=0.
 \]
\item 任意の $0\ne u\in\ell$ について
 \[
  \mathfrak q_B(u)\ne0.
 \]
 従って、$u$ を初速度とする代数変形は全次数に積分できるが、
 基点の非零外的導分を $\mathbf C[t]/(t^3)$ まで持ち上げることはできない。
\item $B_{\mathcal O}$ を $\mathcal O=\mathbf C[[t]]$ 上の任意の平坦な形式変形とし、
 その一次類が $\ell\setminus\{0\}$ に属するとする。
 対称性保持のFittingイデアルの引戻しは正確に
 \begin{equation}
  I_B\mathcal O=(t^2)
  \label{eq:completion-arc-Fitting-v38}
 \end{equation}
 となる。相対Hochschild群は、加群として非標準的に
 \begin{equation}
 \begin{aligned}
 HH^0_{\mathcal O}(B_{\mathcal O})&\simeq\mathcal O,&
 HH^1_{\mathcal O}(B_{\mathcal O})&=0,\\
 HH^2_{\mathcal O}(B_{\mathcal O})&\simeq
       \mathcal O^{34}\oplus\mathcal O/(t^2),&
 HH^3_{\mathcal O}(B_{\mathcal O})&\simeq\mathcal O^{39}
 \end{aligned}
 \label{eq:completion-arc-relative-HH-v38}
 \end{equation}
 と同型であり、四次以上は零である。
 特に一般ファイバーのHochschild次元は $(1,0,34,39)$ となる。
\end{enumerate}
\end{theorem}

定理\ref{thm:fixed-quantum-hull-v35}により、同じ非零一次方向は外側では
$\mathbf C[t]/(t^3)$ へ延長できない。
上の定理は、その方向を内側で延長したときに失われる構造を特定する。
二つの二次障害の値域はそれぞれ $HH^3(\Lambda)$ と $Q$ であり、
両者の自然同一視を必要とせずに、この対応が成立する。

\subsection{双対数上の標準補完の明示的構成}

$A=\mathbf Q[\varepsilon]/(\varepsilon^2)$ とし、以下の三変数代数を取る。
\begin{equation}
\begin{aligned}
 yx&=(1+2\varepsilon)xy-(1+\varepsilon)x^2-2\varepsilon xZ,\\
 Zx&=(1+2\varepsilon)xZ-2(1+\varepsilon)xy+(1+\varepsilon)x^2,\\
 Zy&=yZ+(1+\varepsilon)xy-(1+\varepsilon)xZ.
\end{aligned}
\label{eq:completion-base-relations-v38}
\end{equation}
この代数を $C_\varepsilon$ と書く。
次数付き辞書式順序 $Z>y>x$ に関する唯一の重なりは $Zyx$ であり、
二通りの簡約が一致する。従って $x^a y^bZ^c$ は $A$ 上の基底となる。
正則な正規元 $x$ の正規化自己同型は
\begin{equation}
\begin{aligned}
 \sigma(x)&=x,\\
 \sigma(y)&=(1+2\varepsilon)y-(1+\varepsilon)x-2\varepsilon Z,\\
 \sigma(Z)&=(1+2\varepsilon)Z-2(1+\varepsilon)y+(1+\varepsilon)x
\end{aligned}
\label{eq:completion-sigma-v38}
\end{equation}
であり、$ux=x\sigma(u)$ が成立する。
表示した写像が関係式を保存することは $x$ の正則性からも分かる。
次数一での行列は可逆であり、各次数の自由性によって自己同型となる。

PBW順序で書いた次数五の元を
\begin{equation}
\begin{aligned}
 H_0={}&x^2yZ^2+x^3yZ-2x^3Z^2+x^2y^2Z+x^3y^2
              -x^4y+x^4Z-2x^2y^3\\
      &-\frac53xy^3Z+\frac53xy^4+\frac23y^5,\\
 H_1={}&-\frac83y^5-\frac53y^3Z^2
       +\frac52xyZ^3+\frac52xy^2Z^2-5xy^3Z,
 \qquad H=H_0+\varepsilon H_1
\end{aligned}
\label{eq:completion-H-v38}
\end{equation}
とする。三つの関係式から、
\begin{equation}
 [x,H],\ [y,H],\ [Z,H]\ \in x^2 C_\varepsilon
 \label{eq:completion-divisibility-v38}
\end{equation}
を直接確認できる。
この確認には $H$ を $x^2C_\varepsilon$ で割った剰余
\[
 \frac23y^5+\frac53xy^4-\frac53xy^3Z
 +\varepsilon\left(-\frac83y^5-\frac53y^3Z^2
       +\frac52xyZ^3+\frac52xy^2Z^2-5xy^3Z\right)
\]
だけを用いればよい。

一般に、正則な正規元 $x$ が $ux=x\sigma(u)$ を満たし、
$[C,H]\subset x^2C$ ならば
\[
 D_H(u)=x^{-2}(uH-Hu)
\]
は $C$ に値を取る $\sigma^2$ 導分である。
右辺は一意な左割算であり、実際に局所化を取る必要はない。
積の展開と $ux^2=x^2\sigma^2(u)$ から
\[
 D_H(uv)=\sigma^2(u)D_H(v)+D_H(u)v
\]
を得る。従って
\begin{equation}
 \mathcal S_\varepsilon
   =C_\varepsilon[w;\sigma^2,D_H],\qquad \deg w=3
 \label{eq:completion-Ore-v38}
\end{equation}
は結合的であり、$x^a y^bZ^cw^d$ を $A$ 上の基底に持つ。
また $x^2w+H$ は中心元である。

特殊ファイバーと既存のCGG表示との対応を明記する。
$\alpha=-1$, $s=1$ とし、$Z\mapsto z+y$, $w\mapsto w$ と置くと、
式\eqref{eq:completion-base-relations-v38}の $\varepsilon=0$ への特殊化は
CGG代数の三変数部分と一致する。
この対応で
\[
 H_0\longmapsto q\left(xyz-\frac23y^3\right),\qquad
 x^2w+H_0\longmapsto
 J=x^2w+q\left(xyz-\frac23y^3\right).
\]
$N(J)=X_{-1}(J)=0$ であるから $J$ はCGG積で中心元となる。
三変数部分で $\sigma=T^{-1}$、$T=1+N$ であり、
$w$ との積の最高 $w$ 次数部分は $\sigma^2$ によって決まる。
中心性から残る項は一意に $D_{H_0}$ となるため、全特殊ファイバーは
$\mathcal S_{1,1}(-1,1)$ と一致する。

$\mathcal S_\varepsilon$ の次数 $0,\ldots,5$ の標準有限代数を取れば、
階数百四十一の平坦な六項代数変形 $\Lambda_\varepsilon$ を得る。
その $(0,1,3,4)$ cornerは階数五十六の変形 $B_\varepsilon$ となる。
以下の計算では、経路順の積の規約で
\[
 F_{ik}=T^{-k}:\operatorname{Hom}_{\alpha=0}(i,k)
                 \longrightarrow\operatorname{Hom}_{\alpha=-1}(i,k)
\]
を用い、さらに固定したPBW基底との変換を施して、
$B_\varepsilon$ の一次cochainを前節の $\alpha=0$ の基底へ移す。

\subsection{接線の位置と二次障害の厳密計算}

$U\in Z^2(B)$ を上の $B_\varepsilon$ の一次cochain、$u=[U]$ とする。
前節と同じ境界pivotと自由変数順序による $K$ の基底で、その座標は
\begin{equation}
 [u]_K=
 \left(\frac{41667}{4},\frac{324203}{180},\frac{1387}{2},
 -\frac{732409}{160},\frac{2794373}{480},-\frac{1571}{60}\right)^t.
 \label{eq:completion-K-coordinates-v38}
\end{equation}
特に $u\ne0$ なので、この変形は外側の唯一の非零一次方向を生成する。
$L(u)=0$ を確認し、自然写像 $K\to Q$ によって
\begin{equation}
 [u]_Q=
 \left(0,0,-45,-\frac{2519}{160},-\frac{3861}{80},
                  \frac{17281}{80}\right)^t\ne0
 \label{eq:completion-Q-coordinates-v38}
\end{equation}
を得る。
さらに
\begin{equation}
 \mathfrak q_B(u)=
 \left(0,120,1640,-\frac{5917}{12},-\frac{2269}{2},
                     \frac{42347}{6}\right)^t\ne0.
 \label{eq:completion-secondary-value-v38}
\end{equation}
これらの数値自体は選んだ有理基底に依存するが、三つの非零性と
$u\in K$ という結論は基底に依存しない。

計算は二通りの手順で照合した。
第一に、式\eqref{eq:completion-K-coordinates-v38}を前節の二次写像へ代入する。
第六座標の平方の第二成分は
\[
 \frac{432000}{2468041}\left(-\frac{1571}{60}\right)^2=120
\]
となる。
第二に、実際のOre族から得た $1584$ 成分の $U$ に対して
\[
 d^1\eta=L_\delta U,\qquad d^2V=U\circ U
\]
を有理数上で解き、$L_\delta V+L_\eta U$ のコホモロジー類を
$Q$ へ射影する。この手順でも式\eqref{eq:completion-secondary-value-v38}を得る。
前節の作用行列も同じ積から再計算して一致を確認する。
六項全体の特殊ファイバーは全 $588$ 可合成二入力について、
一次結合律は全 $1026$ 可合成三入力について確認した。
完全なcochain、二つの補正、商座標、有理係数、および読取専用の再計算コードを
付属の証明書に記録する。

\begin{proof}[定理\ref{thm:vk-completion-secondary-v38}の証明]
構成した変形が六項全体に定義され、$u\ne0$ であるため、$\ell=\mathbf C u$ である。
式\eqref{eq:completion-Q-coordinates-v38}と
\eqref{eq:completion-secondary-value-v38}が第一、第二の主張を与える。
二次障害の選択独立性は補題\ref{lem:vk-quadratic-obstruction-v37}による。
$R_B$ は滑らかなので、任意の一次方向に全形式弧が存在する。

第三の主張を示す。定理\ref{thm:vk-jump-column-v37}の最小複体を
$\mathcal O$ へ引き戻し、その残余列を $a(t)$ と書く。
一次方向が $K$ に入るため $a(t)\in t^2\mathcal O^{35}$ である。
$t^2$ の係数の $Q$ における像は、符号と基底変換を除いて
$\mathfrak q_B(u)$ である。
二次の積補正を変更しても、この像は変わらない。
従って、少なくとも一つの成分は $t^2$ で割った後に単元となり、
式\eqref{eq:completion-arc-Fitting-v38}を得る。
これは、どの形式弧を選んでも成立する。

DVR上の基底変換によって列 $a(t)$ は $(t^2,0,\ldots,0)^t$ となる。
従って引き戻した最小複体は
\[
 \mathcal O[0]\oplus
 \bigl(\mathcal O\xrightarrow{\,t^2\,}\mathcal O\bigr)_{1,2}
 \oplus\mathcal O^{34}[-2]\oplus\mathcal O^{39}[-3]
\]
に同型である。
コホモロジーを取れば式\eqref{eq:completion-arc-relative-HH-v38}を得る。
$t$ を可逆にすると二項複体が可縮となるため、一般ファイバーの次元列も従う。
\end{proof}

\begin{corollary}[初速度別の対称性消失次数]
\label{cor:vk-arcs-order-strata-v38}
$B=B_{1,1}$ の形式弧に沿って $I_B\mathcal O=(t^\nu)$ と書き、
$I_B\mathcal O=0$ のときは $\nu=\infty$ と定める。
固定した一次類 $u$ を持つ全形式弧について、可能な $\nu$ は正確に次の通りである。
\[
\begin{array}{c|c}
 \text{一次類の位置}&\text{可能な }\nu\\\hline
 u\notin K&1\\
 u\in K\setminus\operatorname{im}L&2\\
 u\in K\cap\operatorname{im}L&2,3,4,\ldots,\infty
\end{array}
\]
特に、定理\ref{thm:vk-completion-secondary-v38}の弧に関する結論は、
補完接線を含む全ての $K\setminus\operatorname{im}L$ の方向に成立する。
有限な $\nu$ に対しては
\[
 HH^1_{\mathcal O}(B_{\mathcal O})=0,\qquad
 HH^2_{\mathcal O}(B_{\mathcal O})
       \simeq\mathcal O^{34}\oplus\mathcal O/(t^\nu)
\]
となる。
\end{corollary}

\begin{proof}
$u\notin K$ では列 $a(t)$ の一次項 $-L(u)$ が非零なので $\nu=1$ である。
$u\in K\setminus\operatorname{im}L$ では
定理\ref{thm:vk-symmetry-second-neighborhood-v37}により
$\mathfrak q_B(u)\ne0$ であり、直前の証明から $\nu=2$ となる。

$u\in K\cap\operatorname{im}L$ の場合、一次項は零なので $\nu\ge2$ または
$\nu=\infty$ である。
定理\ref{thm:vk-Ga-families-v38}の滑らかな閉形式部分 $Y$ 上に、
一次類 $u$ の弧 $c(t)$ を取れば $\nu=\infty$ を得る。
$Y$ を座標部分空間とする $R_B$ の形式座標を選び、$v\notin K$ を固定する。
任意の整数 $n\ge2$ に対し、弧を $c(t)+t^n v$ へ変えると、一次類は $u$ のままで
\[
 a(c(t)+t^nv)=t^n\bigl(-L(v)+O(t)\bigr)
\]
となる。従って $\nu=n$ を得る。
相対Hochschild群は、一列の微分のDVR上の標準形から従う。
\end{proof}

$B_{0,3}$ でも、初速度が $K$ の外なら消失次数は一である。
$K$ の中の固定した初速度については、滑らかな四次元の $Z_{B_{0,3}}$ を用いる
同じ構成により、可能な次数は任意の整数 $n\ge2$ と $\infty$ となる。
五次元または四次元の不変な接空間に入ることだけでは、
任意に選んだ弧で対称性が三次まで保持されるとは限らない。
二次以降の積補正を変更する自由度が、この差を生む。

特に、補完接線を持つ任意の形式曲線と対称性保持部分との形式的な交叉は、
長さ二の $\operatorname{Spec}\mathbf C[t]/(t^2)$ である。
代数自身には滑らかな三十五次元の変形空間がありながら、
この接線に沿う外的導分の保持には正確に二次の限界がある。
本節で扱う補完は、固定した六項標準列とその四項cornerの同時変形である。

\section{外側の結合性障害から内側の対称性障害への推移}
\label{sec:obstruction-transgression-v40}

標準補完の非零接方向では、外側の結合性と内側の導分の持上げが
ともに二次で妨げられる。この対応を障害類の間の写像として説明するには、
結合性障害と、導分の一次補正を一緒に記録する必要がある。
本節ではその自然な比較を構成する。
通常の $HH^3$ 全体を内側の余核と同一視することは用いない。

代数と一つの導分を同時に変形する一般理論はAssDer対の理論であり、
そのcochain複体と障害類はDas--Mandal
\cite[\S2.1, Proposition~4.9, Theorem~4.10]{DasMandal2026}により与えられている。
以下ではこの既存の理論を全充満部分圏への制限と組み合わせ、
VKの標準補完で現れる二種類の障害を比較する。

\subsection{導分を記録した障害複体}

有限個の対象を持つ代数 $A$ と、頂点冪等元を零にする導分 $\delta_A$ を取る。
頂点の分離的代数に相対的な正規化Hochschild複体を $C^\bullet(A)$ と書く。
本節では符号を統一するため
\[
 d=[\mu,-],\qquad L_A=[\delta_A,-]
\]
とする。従って通常のHochschild微分は $C^n$ 上で $(-1)^{n-1}d$ であり、
コホモロジーは同じである。
$L_A$ は $d$ と可換なので
\begin{equation}
 \mathcal P_A^n=C^n(A)\oplus C^{n-1}(A),\qquad
 \partial(f,g)=(df,L_Af-dg)
 \label{eq:pair-cone-v40}
\end{equation}
は複体となる。これは $L_A$ のmapping coneを一つずらしたものである。
次数零まで含める規約は内的導分も記録するが、本節で用いる
$\mathcal P_A^2\to\mathcal P_A^3\to\mathcal P_A^4$ と三次コホモロジーには
AssDer複体の低次数の切詰めによる差はない。

coneの長完全列は、自然な短完全列
\begin{equation}
 0\longrightarrow
 \operatorname{coker}(L_A|_{HH^2(A)})
 \xrightarrow{\ i_A\ }H^3(\mathcal P_A)
 \xrightarrow{\ p_A\ }
 \ker(L_A|_{HH^3(A)})
 \longrightarrow0
 \label{eq:pair-obstruction-extension-v40}
\end{equation}
を与える。$i_A$ は $[q]\mapsto[(0,q)]$、$p_A$ は第一成分への射影である。
この完全列の分裂は選ばない。

一次の積補正 $a\in C^2(A)$ と導分補正 $\eta\in C^1(A)$ が
\[
 da=0,\qquad d\eta=L_Aa
\]
を満たすとする。二次の誤差を
\begin{equation}
 O_A=\frac12[a,a],\qquad Q_A=[\eta,a]
 \label{eq:enhanced-primary-obstruction-v40}
\end{equation}
と置くと、Jacobi恒等式から
$dO_A=0$、$dQ_A=L_AO_A$ である。
従って $(O_A,Q_A)$ は $\mathcal P_A^3$ のcocycleを定める。
積と導分をともに二次まで延長できることは、その類が零であることと同値である。
実際、二次補正 $(v,\zeta)$ が満たす式は
$\partial(v,\zeta)=-(O_A,Q_A)$ である。

\begin{proposition}[一次類だけから定まる増強された障害]
\label{prop:enhanced-canonical-obstruction-v40}
$HH^1(A)=\mathbf C[\delta_A]$ とする。
任意の $\bar a\in\ker(L_A|_{HH^2(A)})$ について、
式\eqref{eq:enhanced-primary-obstruction-v40}の類
\[
 \widetilde o_A(\bar a)=[(O_A,Q_A)]\in H^3(\mathcal P_A)
\]
は、$a$ と $\eta$ の選び方によらない二次写像である。
その $p_A$ による像は通常の結合性障害
$o_A(\bar a)=[\tfrac12[a,a]]$ である。
\end{proposition}

\begin{proof}
まず $a$ を固定する。二つの $\eta$ の差は導分であるから、仮定により
$c\delta_A+db$ と書ける。第二成分の差は
\[
 [c\delta_A+db,a]=d\bigl(c\eta+[b,a]\bigr)
\]
であり、coneの境界となる。

次に $a'=a+df$ とし、$\eta'=\eta+L_Af$ を選ぶ。
\[
 X=-[a,f]+\tfrac12[f,df],\qquad
 Y=-[\eta,f]+\tfrac12[f,L_Af]
\]
と置くと、graded Jacobi恒等式により
\[
 (O_A',Q_A')-(O_A,Q_A)=\partial(X,Y)
\]
となる。他の $\eta'$ は先ほどの議論で扱える。
従って代表と導分補正の両方に依存しない。
スカラー倍と和に対して式を展開すれば二次性も従う。
\end{proof}

\subsection{全充満部分圏への自然な比較}

$B=eAe$ を頂点の一部を残した全充満部分圏の代数とし、
$\delta_B=\delta_A|_B$ とする。
相対cochainの制限
\[
 \rho:C^\bullet(A)\longrightarrow C^\bullet(B)
\]
は微分とGerstenhaber括弧を保ち、$\rho L_A=L_B\rho$ を満たす。
これは全充満部分圏へのHochschild制限である。
一般に $a\mapsto eae$ 自体を代数準同型として扱う必要はない。
従って $\rho$ は式\eqref{eq:pair-cone-v40}と
\eqref{eq:pair-obstruction-extension-v40}の間の自然な写像を与える。

\begin{theorem}[結合性障害が消えた後の対称性障害]
\label{thm:obstruction-transgression-v40}
上記の一次対 $(a,\eta)$ に対し、$\rho[O_A]=0$ in $HH^3(B)$ とする。
\[
 dv=-\rho O_A
\]
を満たす $v\in C^2(B)$ を選ぶと
\begin{equation}
 q_B=L_Bv+[\rho\eta,\rho a]
 \label{eq:obstruction-transgression-v40}
\end{equation}
はHochschild二cocycleである。
その余核における類は $v$ の選び方によらず、自然な等式
\begin{equation}
 i_B([q_B])=\rho_*[(O_A,Q_A)]
 \quad\text{in }H^3(\mathcal P_B)
 \label{eq:obstruction-transgression-naturality-v40}
\end{equation}
で特徴付けられる。
特に右辺からの逆像は一意である。
$HH^1(A)=\mathbf C[\delta_A]$ ならば、これは一次類 $\bar a$ だけで定まる。
\end{theorem}

\begin{proof}
$dQ_A=L_AO_A$ と $dL_B=L_Bd$ から $dq_B=0$ となる。
別の選択との差 $v'-v$ はcocycleなので、$q_B'-q_B=L_B(v'-v)$ は
余核で零である。
さらにcochainの段階で
\[
 (\rho O_A,\rho Q_A)+\partial(v,0)=(0,q_B)
\]
となる。これが式\eqref{eq:obstruction-transgression-naturality-v40}を与え、
$i_B$ の単射性が一意性を与える。
最後の主張は命題\ref{prop:enhanced-canonical-obstruction-v40}による。
\end{proof}

固定した二次の積補正 $v$ に導分を持ち上げる障害は $[q_B]\in HH^2(B)$ である。
余核を取ることは、二次の積補正もcocycleだけ変更する自由を許すことに対応する。
従って、余核での非零性はどの二次積補正を選んでも導分を持ち上げられないことを意味する。
この区別により、特定の形式弧と全ての形式弧に関する主張を分けられる。

同じ比較は任意の小拡大に成立する。
核 $I$ が極大イデアルにより零化されるArtin基底の小拡大について、
積と導分の仮の持上げを $\widetilde\mu,\widetilde\delta$ とすれば、
誤差
\[
 O=\tfrac12[\widetilde\mu,\widetilde\mu]\in C^3(A)\otimes I,
 \qquad Q=[\widetilde\delta,\widetilde\mu]\in C^2(A)\otimes I
\]
は同じconeのcocycleである。持上げを変更するとconeの境界だけ変わる。
制限と、内側の結合性誤差を消す操作を順に適用すれば、
式\eqref{eq:obstruction-transgression-naturality-v40}と同じ一意な余核類を得る。
これは既に与えられた導分付き変形の延長に関する主張であり、
任意の外側の代数変形に導分が自動的に付随するとは仮定していない。

\subsection{標準六項補完の障害線への適用}

\begin{corollary}[二つの非零障害を結ぶ標準的な写像]
\label{cor:vk-primary-obstruction-line-v40}
$A=\Lambda=A_{1,1}$、$B=B_{1,1}$、$\delta=\log(1+N)$ とする。
標準六項代数の一次変形空間 $W=HH^2(\Lambda)$ は一次元であり、
\[
 O_{\mathrm{std}}=
 \operatorname{im}\bigl(\operatorname{Sym}^2W
       \xrightarrow{\ o_\Lambda\ }HH^3(\Lambda)\bigr)
\]
も一次元である。
固定した $\delta$ と標準cornerに関して、自然な非零線形写像
\begin{equation}
 \mathcal T_\delta:O_{\mathrm{std}}\longrightarrow
 \operatorname{coker}(L_B:HH^2(B)\to HH^2(B))
 \label{eq:primary-obstruction-line-map-v40}
\end{equation}
が存在し、全ての $a\in W$ について
\[
 \mathcal T_\delta(o_\Lambda(a))=\mathfrak q_B(\rho a)
\]
となる。右辺は既出の内在的二次障害である。
\end{corollary}

\begin{proof}
$\delta$ は有限個の射空間上で冪零なので、全cochainに対する $L_\Lambda$ も冪零である。
一次元の $HH^2(\Lambda)$ 上の作用は従って零であり、
任意の一次類に導分の一次補正を選べる。
$HH^1(\Lambda)=\mathbf C[\delta]$ なので、増強された障害は一次類だけで定まる。

六項のhullは $\mathbf C[[\gamma]]/(\gamma^2)$ であるから、
$o_\Lambda:\operatorname{Sym}^2W\to HH^3(\Lambda)$ は非零である。
従ってその像への写像は同型となる。
他方、四項の全変形は無障害なので、その結合性障害の制限は零である。
定理\ref{thm:obstruction-transgression-v40}を適用し、
上の同型の逆と合成すれば式\eqref{eq:primary-obstruction-line-map-v40}を得る。

外側で選んだ導分補正の制限は、既存の四項計算で選んだ補正と
Hochschild一cocycleだけ異なる。
$HH^1(B)=\mathbf C[\delta]$ であるため、その差は余核類を変えない。
従って既出の有理基底と接線の規格化では、写像の値は
\[
 \left(0,120,1640,-\frac{5917}{12},-\frac{2269}{2},
                    \frac{42347}{6}\right)^t
\]
であり、非零である。
元の四項bar積から式\eqref{eq:enhanced-primary-obstruction-v40}と
\eqref{eq:obstruction-transgression-v40}を再構成し、coneのcocycle条件、
代表変更の境界式、および余核への射影を有理数上で直接検証した。
\end{proof}

この写像は、補完方向をスカラー倍しても変わらない。
導分の規格化を変えれば写像は同じスカラー倍だけ変わるが、その像の直線は変わらない。
従って、外側で消せない結合性障害は、導分を記録した障害として制限すると、
内側で消せない対称性障害へ移る。
本節で得た標準的な線形写像の定義域は $O_{\mathrm{std}}$ であり、
七次元の $HH^3(\Lambda)$ 全体への拡張は主張していない。
一般の重みで高次の障害を比較する場合にも、小拡大ごとに導分付き障害を追う必要がある。

\section{検証と残る変形問題}

全重みのhullについては、Poissonスキームの全局所環、形式弧の再中心化、
Hochschild一次差分と軌道切片を用いた。有限個の階数からの外挿ではない。
端点の全量子化はOre構造によって各ファイバーでAS正則かつCYとなる。
関数斜体の分類は先頭単項式の不変量と整数指数の可逆変換による。
PI次数の有限個の円分計算は、全位数の証明の確認として位置付ける。

VKの全次数の対称性保持族は、原始射に支持を持つ積補正の合成が零となることと、
平方零横断方向に関する有限な線形方程式によって構成する。
分母が原点で単元であることを確認し、得られた有理式について
全二千五百五十六個の結合則・導分方程式を多項式恒等式として照合した。
これにより五次元部分の二重構造を全次数で決定し、
三十五変数基底の中のイデアルも $(z_1,\ldots,z_{29},\gamma^2)$ と同定する。
純粋部分への還元では、商cochain複体の一次・二次コホモロジーの消滅を使って、
全次数でフィルトレーションを復元する。
補完接線は双対数上の具体的なOre代数から求め、全六項の合成を検査した。
同じ二次障害を、既存二次写像への代入と、実際のcochainからの独立した持上げ計算で求める。
通常の検証実行は入力と証明書を変更せず、完全な結果を凍結した証明書と比較する。

この結果は、標準六項補完像と対称性保持部分の位置関係を含む。
補完像は二重構造の横断長さ二の切片であり、滑らかな五変数部分の全てが
標準六項列へ補完できるという意味ではない。
別の例外列を許す半直交補完可能性の分類は、なお未決である。
導分付き障害複体を用いることで、外側の結合性障害線から内側の対称性障害への自然な比較写像を構成した。
一般の $HH^3$ 全体への写像は主張しない。
全重みで純粋部分は内在的に持ち上がり、同定を固定した相対変形は滑らかとなる。
純粋部分の $HH^1$ は一次元で、その生成元が全共鳴へ持ち上がるため、
接続写像は零であり、通常hullは純粋部分のhullと滑らかな因子に分解する。
一般の純粋hull $R_{N_m}$ の決定は未決である。
一般の重み・共鳴で、純粋部分の対称性保持方向を原始射の接合が何次まで切り詰めるかは、
この全次数の決定を拡張する次の問題となる。

端点では関数斜体を全パラメータで分類できたが、同じ斜体を持つ混合項 $H$ の方向を
射影圏や導来圏の不変量で区別する問題が残る。
公表稿XI\cite{VKPrimitive36}では、$H$ に依存しない中心開部分に対し、
境界を含む別の中心開部分を一般化Weyl代数として表示する。
その全中心は、$H\bmod x$ の巡回ノルム $P_H$ を用いた
$\Xi\Omega=P_H$ というアフィン改変であり、
標識付き境界方程式はこの剰余類を巡回作用まで復元する。
全中心ファイバーの単純加群と根基も切断間隔から決まり、
同じ関数斜体と同じ中心開部分を持ちながら異なる境界を持つ具体例を得る。

同稿では、偶数位数の特異部分上で根基を含む代数そのものを降下し、
二つの単純加群を分ける連結な二重被覆を決定する。
滑らかな中心部分のBrauer類の周期と一般点の指数は、共に冪根の位数 $\ell$ である。
従って完成局所的な行列環や非可換クレパント解消の表示は、
元の中心上での自己準同型環表示を意味しない。

原稿XIでは、$H=f(Q,u)+xG$ の全中心も決定し、
重み四・位数五で標識を忘れた全中心が非同型になる二つの混合項を示す。
この例では全アフィン代数のMorita同値とdg導来Morita同値も排除できる。
混合項を消せる部分族では全中心ファイバーを分類し、$f$ が $Q$ の倍数なら
最深点のファイバーと完備局所指数がそのまま保たれることも示す。
任意の混合項の深い境界と、混合項を動かした射影圏の同値分類は未決である。
CGG族内の導来分類を全miniversal族の分類へ拡大して述べることはしない。

\begin{thebibliography}{99}
\interlinepenalty=10000
\bibitem{VKPoisson34} 本研究、公表原稿IX、
\emph{重み付きPoisson変形の非被約成分とCalabi--Yau量子化}, v34.
\nolinkurl{weighted_poisson_components_ja_v34.pdf}.
\bibitem{VKAmbient32}本研究、公表原稿VIII、\emph{半直交成分の変形と重み付き射影空間への補完}, v32.
\nolinkurl{vk_ambient_completions_ja_v32.pdf}.

\bibitem{LecoutreSierra2019}
C.~Lecoutre and S.~J.~Sierra,
\emph{A new family of Poisson algebras and their deformations},
Nagoya Math. J. \textbf{233} (2019), 32--86,
\href{https://doi.org/10.1017/nmj.2017.29}{doi:10.1017/nmj.2017.29};
\href{https://arxiv.org/abs/1611.04306}{arXiv:1611.04306}.

\bibitem{Zhang1996}
J.~J.~Zhang,
\emph{Twisted graded algebras and equivalences of graded categories},
Proc. London Math. Soc. (3) \textbf{72} (1996), 281--311,
\href{https://doi.org/10.1112/plms/s3-72.2.281}{doi:10.1112/plms/s3-72.2.281}.

\bibitem{Schlessinger1968}
M.~Schlessinger, \emph{Functors of Artin rings},
Trans. Amer. Math. Soc.\ \textbf{130} (1968), 208--222.
\url{https://doi.org/10.1090/S0002-9947-1968-0217093-3}.
\url{https://math.uchicago.edu/~amathew/schlessingerdef.pdf}.

\bibitem{KellerDerived2003}
B.~Keller, \emph{Derived invariance of higher structures on the Hochschild complex},
preprint, August 2003, revised May 16, 2018, Theorem~3.2 and its corollary.
\url{https://webusers.imj-prg.fr/~bernhard.keller/publ/dih.pdf}.

\bibitem{Manetti2005}
M.~Manetti, \emph{Deformation theory via differential graded Lie algebras},
Seminari di Geometria Algebrica 1998--1999, Scuola Normale Superiore (1999);
corrected version arXiv:math/0507284 (2005), Corollary~3.2.
\url{https://arxiv.org/abs/math/0507284}.

\bibitem{LinZhangSerre2025}
X.~Lin and S.~Zhang,
\emph{Serre algebra, matrix factorization and categorical Torelli theorem for hypersurfaces},
Math. Ann. \textbf{391} (2025), 163--177.
\url{https://doi.org/10.1007/s00208-024-02915-8}.
\url{https://arxiv.org/abs/2310.09927}.

\bibitem{MinamotoMori2011}
H.~Minamoto and I.~Mori,
\emph{The structure of AS-Gorenstein algebras},
Adv. Math. \textbf{226} (2011), no.~5, 4061--4095.
\url{https://doi.org/10.1016/j.aim.2010.11.004}.
Preprint title: \emph{Structures of AS-regular algebras}.
\url{https://www.math.kyoto-u.ac.jp/preprint/2010/06Minamoto.pdf}.

\bibitem{RickardMorita1989}
J.~Rickard, \emph{Morita theory for derived categories},
J. London Math. Soc. (2) \textbf{39} (1989), no.~3, 436--456.
\url{https://doi.org/10.1112/jlms/s2-39.3.436}.

\bibitem{RickardDerived1991}
J.~Rickard, \emph{Derived equivalences as derived functors},
J. London Math. Soc. (2) \textbf{43} (1991), no.~1, 37--48.
\url{https://doi.org/10.1112/jlms/s2-43.1.37}.

\bibitem{VKPrimitive36} 本研究、公表原稿XI、\emph{重み付きCalabi--Yau量子化の中心と表現}, v40.
\nolinkurl{quantum_primitive_spectra_ja_v40.pdf}.
\bibitem{LWWEndpoint37}
L.~Liu, S.~Wang, and Q.~Wu,
\emph{Twisted Calabi--Yau property of Ore extensions},
Journal of Noncommutative Geometry \textbf{8} (2014), 587--609,
doi:\href{https://doi.org/10.4171/JNCG/165}{10.4171/JNCG/165},
arXiv:\href{https://arxiv.org/abs/1205.0893}{1205.0893}.

\bibitem{ReyesRogalski2022}
M.~L.~Reyes and D.~Rogalski,
\emph{Graded twisted Calabi--Yau algebras are generalized Artin--Schelter regular},
Nagoya Mathematical Journal \textbf{245} (2022), 100--153,
doi:\href{https://doi.org/10.1017/nmj.2020.32}{10.1017/nmj.2020.32}.

\bibitem{StacksFitting37}
The Stacks Project Authors,
\emph{The Stacks Project}, Tags 07Z8 and 07ZA,
Fitting ideals and base change.
\url{https://stacks.math.columbia.edu/tag/07Z8}.
\url{https://stacks.math.columbia.edu/tag/07ZA}.

\bibitem{AlevDumas1994Bir38}
J.~Alev and F.~Dumas,
\emph{Sur le corps des fractions de certaines alg\`ebres quantiques},
Journal of Algebra \textbf{170} (1994), 229--265,
doi:\href{https://doi.org/10.1006/jabr.1994.1336}{10.1006/jabr.1994.1336}.

\bibitem{RichardSolotar2006Bir38}
L.~Richard and A.~Solotar,
\emph{Isomorphisms between quantum generalized Weyl algebras},
Journal of Algebra and Its Applications \textbf{5} (2006), 271--285,
arXiv:\href{https://arxiv.org/abs/math/0402413}{math/0402413}.

\bibitem{Bresar2010Bir38}
M.~Bre\v{s}ar,
\emph{An alternative approach to the structure theory of PI-rings},
arXiv:\href{https://arxiv.org/abs/1003.5275}{1003.5275}, 2010.

\bibitem{DasMandal2026}
A.~Das and A.~Mandal,
\emph{Extension, deformation and categorification of AssDer Pairs},
Theory Appl. Categ. \textbf{45} (2026), no.~1, 1--32.
\url{https://tac.mta.ca/tac/volumes/45/1/45-01abs.html}.
\url{https://arxiv.org/abs/2002.11415}.

\end{thebibliography}

\end{document}
